\documentclass[paper=186mm:234mm,pagesize=auto,fleqn,11pt,dvipdfmx]{scrbook}
\usepackage[utf8]{inputenc}
\usepackage[T1]{fontenc}
\usepackage[ngerman]{babel}
\usepackage{microtype}
\usepackage{amsmath}
\usepackage{amssymb}
\usepackage{amsthm}
\usepackage{proof}
\usepackage{booktabs}
\usepackage{graphicx}
\usepackage[all]{xy}
\usepackage{enumitem}
\usepackage{fitch}

\usepackage{listings}
\lstset{basicstyle=\ttfamily}

\usepackage{mdframed}
\usepackage{lipsum}

\usepackage{libertine}
\usepackage[libertine]{newtxmath}
% \usepackage[scaled=0.80]{DejaVuSans}
\usepackage[scaled=0.78]{DejaVuSansMono}
% \renewcommand\ttdefault{lmvtt}
\addtokomafont{disposition}{\rmfamily}

\usepackage[automark]{scrlayer-scrpage}
\ohead{\pagemark}
\ihead{\headmark}
\chead{}
\cfoot{}
\ofoot{}

\usepackage{geometry}
\geometry{left=34mm,right=24mm,top=26mm,bottom=23mm,%
footskip=10mm, headsep=6mm}

\usepackage{color}
\definecolor{gray1}{RGB}{80,80,80}

% Für den Druck auf {0,0,0} stellen.
\definecolor{c1}{RGB}{0,40,80}

\usepackage[colorlinks=true,linkcolor=c1,urlcolor=c1,citecolor=c1]{hyperref}
% \usepackage[colorlinks=true,linkcolor=black]{hyperref}
\usepackage{pdfpages}

\newcommand{\strong}[1]{\textbf{#1}}

\newtheoremstyle{rmbox}%
  {0pt}% space above
  {0pt}% space below
  {}% bodyfont
  {}% indent
  {\normalfont\bfseries}% head font
  {~}% punctuation between head and body
  {0pt}% space after theorem head
  {\thmname{#1} \thmnumber{#2}\thmnote{ (#3)}.}

\theoremstyle{rmbox}
\newtheorem{Definition}{Definition}
\newtheorem{Satz}{Satz}
\newtheorem{Lemma}[Satz]{Lemma}
\newtheorem{Korollar}[Satz]{Korollar}

\numberwithin{Definition}{chapter}
\numberwithin{Satz}{chapter}

\newenvironment{Beweis}[1][Beweis]%
  {\par\noindent\strong{#1.}\;\ignorespaces}%
  {\par\addvspace{\topsep}}

\definecolor{greenblue}{rgb}{0.0,0.42,0.3}
\definecolor{grayblue}{rgb}{0.1,0.2,0.4}
\definecolor{bgreen}{rgb}{0.94,0.94,0.84}
\definecolor{bblue}{rgb}{0.9,0.92,0.94}
\definecolor{hltext}{rgb}{0.7,0,0.6}

% Muss je nach Einstellung des Drucksystems neu justiert werden.
\definecolor{bordercolor}{rgb}{0.5,0.5,0.5}

\surroundwithmdframed[topline=false,rightline=false,bottomline=false,%
  linecolor=bordercolor, linewidth=3.5pt, innerleftmargin=6pt,%
  innertopmargin=0pt, innerbottommargin=1pt,%
  innerrightmargin=6pt%
]{Definition}

\newcommand{\framedtheorem}[1]{%
\surroundwithmdframed[topline=false,rightline=false,bottomline=false,%
  linecolor=bordercolor, linewidth=3.5pt, innerleftmargin=6pt,%
  innertopmargin=0pt, innerbottommargin=1pt,%
  innerrightmargin=6pt%
]{#1}}

\framedtheorem{Satz}
\framedtheorem{Lemma}
\framedtheorem{Korollar}

\newcommand{\N}{\mathbb N}
\newcommand{\Z}{\mathbb Z}
\newcommand{\Q}{\mathbb Q}
\newcommand{\R}{\mathbb R}
\newcommand{\C}{\mathbb C}
\newcommand{\ui}{\mathrm i}
\newcommand{\ee}{\mathrm e}
\newcommand{\compc}{\mathsf c}
\newcommand{\defiff}{\;:\Longleftrightarrow\;}
\newcommand{\emdef}[1]{\emph{#1}}
\renewcommand{\qedsymbol}{\ensuremath{\Box}}
\newcommand{\doubleslash}{/\!/}
\newcommand{\category}[1]{\textsf{\textbf{#1}}}
\newcommand{\setsystem}[1]{\mathcal{#1}}
\newcommand{\cond}{\Rightarrow}
\newcommand{\bicond}{\Leftrightarrow}
\newcommand{\lnec}{\Box}
\newcommand{\lpos}{\Diamond}
\newcommand{\Cn}{\operatorname{Cn}}

\DeclareMathOperator{\id}{id}
\DeclareMathOperator{\sur}{sur}
\DeclareMathOperator{\real}{Re}
\DeclareMathOperator{\imag}{Im}
\DeclareMathOperator{\Bild}{Bild}
\DeclareMathOperator{\Abb}{Abb}
\DeclareMathOperator{\dom}{dom}
\DeclareMathOperator{\cod}{cod}
\DeclareMathOperator{\Hom}{Hom}
\DeclareMathOperator{\Ob}{Ob}

\newcommand{\infernote}[1]{\!\text{\scriptsize #1}}
\newcommand{\newlinefirst}{\mbox{}\\}
\renewcommand{\models}{\vDash}
\newcommand{\todo}[1]{{\color{hltext}#1}}
\newcommand{\kw}[1]{\ifmmode%
  \text{\texttt{\textbf{#1}}}%
  \else\texttt{\textbf{#1}}\fi}

\usepackage{makeidx}
\makeindex

\title{Vom Gefüge des Denkens}
\subtitle{Logische Systeme zur Grundlegung der Mathematik,\\
der Informatik und der Philosophie}
\date{Juni 2024}

\begin{document}
\renewcommand{\thepage}{C\arabic{page}}
\includepdf{Cover/Cover.pdf}

\newpage
\thispagestyle{empty}
\mbox{}

\newpage
\addtocounter{page}{-2}
\renewcommand{\thepage}{\arabic{page}}
\thispagestyle{empty}

\maketitle


\noindent
Dieses Buch steht unter der Lizenz Creative Commons CC0 1.0.

\tableofcontents


\chapter{Logisches Schließen}

\section{Grundbegriffe}

\subsection{Schlussregeln}

% Alles Schlussregeln und Axiome müssen unmittelbar einsichtig sein,
% müssen unbedenklich sein. Normalerweise will man in der Wissenschaft
% nicht so informell sein, aber hier sind wir am basalen Grund.
% Selbst in der Semantik wird der hier aufgestellte Formalismus
% benötigt. Der Korrektheitsbeweis würde somit zirkulär.

% Am Anfang ist noch nicht allzu viel zu rechnen. Wir müssen uns
% erst um die Schlussregeln und Axiome bemühen, bevor das erste Theorem
% bewiesen werden kann.

Logisches Schließen findet in einzelnen Schritten statt. Ein Schritt
stellt hierbei immer die Ableitung einer Schlussfolgerung aus einer
oder mehreren Voraussetzungen dar. Die Voraussetzungen heißen
\emph{Prämissen}\index{Praemisse@Prämisse}, die Schlussfolgerung
\emph{Konklusion}\index{Konklusion}. Darstellen wollen wir den Schritt
durch eine waagerechte Linie, wobei die Prämissen oberhalb befindlich
sein sollen, und die Konklusion unterhalb. Der Schritt
\[\dfrac{\text{Wenn es regnet, wird die Straße nass}\qquad\text{Es regnet}}
{\text{Die Straße wird nass}}\]
beschreibt beispielsweise, dass aus den Prämissen »Wenn es regnet, wird
die Straße nass« und »Es regnet« die Konklusion »Die Straße wird nass«
gefolgert wird.

Schlüsse wie der Obige treten in der Mathematik ständig auf. Ihnen allen
liegt ein bestimmtes Muster zugrunde, welches sich durch eine als
\emph{Modus ponens}\index{Modus ponens} oder
\emph{Abtrennungsregel}\index{Abtrennungsregel}
bezeichnete schematische \emph{Schlussregel}\index{Schlussregel}
beschreiben lässt. Es bezeichne hierzu $A\cond B$ die Implikation
»wenn $A$, dann $B$«. Es dürfen nun in
\[\dfrac{A\cond B\qquad A}{B}\]
für $A,B$ beliebige Aussagen eingesetzt werden. So darf »Es regnet»
für $A$ und »Die Straße wird nass« für $B$ eingesetzt werden.

\subsection{Sequenzen}

Das Schließen von Aussagen allein genügt nicht. Um freier argumentieren
zu können, würden wir gerne den Umstand beschreiben können, dass eine
Aussage unter bestimmten Annahmen abgeleitet werden konnte. Diese
Annahmen $A_k$ sind selbst Aussagen. Wir fassen sie zu einer endlichen
Ansammlung
\[\Gamma := [A_1,A_2,\ldots,A_n]\]
zusammen, worunter wir eine endliche Liste, oder auch eine endliche
Menge verstehen wollen, denn man soll mit dieser Liste umgehen können
wie mit einer Menge. Das heißt, es ist nicht von Bedeutung, wie
oft eine Aussage vorkommt oder in welcher Reihenfolge die Aussagen
stehen. Man bezeichnet $\Gamma$ als die \emph{Antezedenz}%
\index{Antezedenz} oder die Liste der \emph{Antezedenzen}. Es wird
$\Gamma$ auch der \emph{Kontext}\index{Kontext}
oder die \emph{Umgebung}\index{Umgebung} genannt, das sind auf die
Typentheorie zurückzuführende Sprechweisen, die einen ganz ähnlichen
Formalismus besitzt. Wir bezeichnen die Symbolik \[\Gamma\vdash A\]
als \emph{Sequenz}\index{Sequenz}. Sie drückt das \emph{Urteil}%
\index{Urteil} aus, dass die Aussage $A$ aus den Annahmen
vermittels Schlussregeln ableitbar ist. Der Modus ponens%
\index{Modus ponens} wird nun in der allgemeinen Form
\[\dfrac{\Gamma\vdash A\cond B\qquad\Gamma\vdash A}{\Gamma\vdash B}\]
beschrieben. Wir argumentieren beim Schließen also ab jetzt nicht mehr
mit den Aussagen selbst, sondern mit den Sequenzen. Dies hat einen wichtigen
Grund, nämlich dass die Berücksichtigung der Abhängigkeit von Annahmen
expliziter Teil des Schließens wird.

Ein Kontext kann auch eine leere Liste sein. Besitzt eine vermittels
Schlussregeln ableitbare Sequenz einen leeren Kontext, so bezeichnet
man die Antezedenz als ein \emph{Theorem}\index{Theorem} im engeren
Sinne. Ein Theorem ist also eine Aussage, die für sich allein gilt,
ohne dass dafür irgendwelche Annahmen getroffen werden müssen.

Für Sequenzen gilt die \emph{Abschwächungsregel}%
\index{Abschwaechungsregel@Abschwächungsregel}. Sie besagt, dass
falls die Aussage $A$ bereits aus $\Gamma$ ableitbar ist, diese
Aussage erst recht ableitbar ist, wenn zu $\Gamma$ weitere Annahmen
$\Gamma'$ hinzugefügt werden. Kurzum gilt die Regel
\[\dfrac{\Gamma\vdash A}{\Gamma,\Gamma'\vdash A}.\]
Hierbei bedeutet $\Gamma,\Gamma'$ die Konkatenation der Listen
$\Gamma$ und $\Gamma'$, also im Wesentlichen dasselbe wie die
Vereinigung $\Gamma\cup\Gamma'$, insofern man die Kontexte als
Mengen betrachtet.

\subsection{Zulässige Schlussregeln}

Wiewohl die Regeln des Schließens den Mechanismus zum Beweis
von Aussagen bilden, ist ihre Rolle sogar noch ein wenig tiefgreifender.
Wir können sie nämlich ebenfalls zur Ableitung \emph{weiterer Regeln}
nutzen. Das heißt, wir können sie dazu nutzen, den logischen Kalkül
selbst zu erweitern. Erweiterungen dieser Art nennen wir
\emph{zulässige Schlussregeln}%
\index{zulaessige Schlussregel@zulässige Schlussregel}.

Mit den bisherigen Regeln ist bereits die zulässige Regel
\[\dfrac{\Gamma\vdash A\cond B\qquad\Gamma'\vdash A}
{\Gamma,\Gamma'\vdash B}\]
ableitbar, die eine allgemeinere Form des Modus ponens darstellt. Man
erhält sie kurzerhand, indem den Prämissen des Modus
ponens jeweils die Abschwächungsregel vorgesetzt wird:
\[\infer{\Gamma,\Gamma'\vdash B}{
  \infer{\Gamma,\Gamma'\vdash A\cond B}{\Gamma\vdash A\cond B}
& \infer{\Gamma,\Gamma'\vdash A}{\Gamma'\vdash A}}
\]
Die einfache Form des Modus ponens erhält man mit $\Gamma':=\Gamma$ als
Spezialfall unter Anwendung der Kontraktionsregel.

\subsection{Implikationseinführung}

Ich möchte mich nun der Frage zuwenden, wie eine Implikation $A\cond B$
bewiesen wird. Intuitiv ist hierzu aus der Annahme $A$ die
Aussage $B$ zu folgern. Das heißt, es genügt die Ableitung
der Sequenz $A\vdash B$. Ein weiteres Mal gilt es noch zu
berücksichtigen, dass ein Beweis auch auf einen vorausgesetzten
Kontext $\Gamma$ beschränkt sein dürfen soll. Reflektiert man darüber
eine Weile, dürfte es der Überlegung nach wohl genügen, dass $A$
einfach dem Kontext $\Gamma$ hinzugefügt wird, woraus $B$ zu folgern
ist. Man gelangt zur Regel
\[\dfrac{\Gamma,A\vdash B}{\Gamma\vdash A\cond B}.\]
Wer diese Regel nicht so leicht fassbar findet, insbesondere nicht
direkt plausibel, ob sie bedenkenlos eingesetzt werden darf, der
ist nicht allein. Es gibt auch logische Kalküle, die diese Regel nicht
explizit enthalten. Sie tritt dennoch als \emph{Deduktionstheorem} in
Erscheinung, ein metalogisches Theorem, dessen Beweis erst erbracht
werden muss. Ich möchte diesen Weg allerdings aus einem bestimmten Grund
nicht gehen. Nämlich ist beim Beweis eigentlich natürliches Schließen
auf der metalogischen Ebene zu verwenden, wenn dies auch in informaler
Weise stattfinden mag. Aber nicht jeder Leser weiß zu diesem Zeitpunkt,
wie akkurates logisches Schließen geht. Der Leser benötigt am Anfang
etwas, um sich an den eigenen Haaren aus dem Sumpf zu ziehen.

\subsection{Axiome}

Zur Komplettierung des Kalküls gesellen sich schließlich auch noch
\emph{Axiome}\index{Axiom} hinzu, das sind gemachte Grundannahmen, die
nicht weiter bewiesen werden müssen. Sie sollten daher möglichst
plausibel, oder besser noch zweifelsfrei einsichtig sein. Für die Logik
selbst genügt das Axiom
\[A\vdash A.\]
Der Kalkül funktioniert dergestalt, dass für $A$ eine beliebige Aussage
eingesetzt werden darf, worunter auch zusammengesetzte Aussagen
wie fallen. Eine gern gewählter Weg der Definition des logischen
Kalküls sieht $A$ als eine metasprachliche Variable, für die eine
beliebige Formel eingesetzt werden darf. Unter dieser Sichtweise
spricht man von einem \emph{Axiomenschema}\index{Axiomenschema}. Wie
eine Schablone produziert es für jede Einsetzung einer konkreten
logischen Formel ein eigenes Axiom.

Anstelle $A,B,C$ werden für metasprachliche Variablen zuweilen auch
die griechischen Buchstaben $\varphi,\psi,\chi$ benutzt. Man muss sie
von atomaren logischen Variablen unterscheiden, für die wir in diesem
Buch, um Missverständnissen aus dem Weg zu gehen, kleine Buchstaben
$a,b,c$ oder $p,q,r$ verwenden werden. Sprachlich suggestiv steht
$\varphi$ für \emph{Formel}\index{Formel} oder \emph{formula}, $a$ für
\emph{Aussage} und $p$ für \emph{proposition}.

In diesem Sinne sind auch die Schlussregeln Schemata. Sofern man
Schlussregeln mit null Prämissen gestattet, lässt sich das
Axiomenschema auch als Regel
\[\dfrac{}{A\vdash A}\]
auffassen. In dieser Weise wollen wir die Anwendung von Axiomen in den
Beweisbäumen darstellen.

Axiome in der Form von Sequenzen heißen auch
\emph{Grundsequenzen}\index{Grundsequenz}.

Wir haben nun die Mittel in der Hand, um erste Theoreme beweisen
zu können. Es ist $A\cond A$ ein Theorem. Der Beweisbaum ist:
\[
\infer[\infernote{Implikationseinführung}]{\vdash A\cond A}{
  \infer[\infernote{Axiom}]{A\vdash A}{}}
\]
Unter der Lesung, dass $A$ eine Metavariable ist, handelt es
eigentlich nicht nur um ein Theorem, sondern um ein Schema von
Theoremen. Setzt man für $A$ bspw. die konkrete Formel $p\cond q$
ein, bekommt man das konkrete Theorem
\[\vdash (p\cond q)\cond (p\cond q).\]
% Es ist zu bemerken, dass die Unterscheidung zwischen Metavariablen
% und atomaren logischen Variablen später durch die Einsetzungsregel
% mehr oder weniger hinfällig wird. Hierbei handelt es sich aber um eine
% höhere Überlegung, deren Beweis der Logiker erbringen will. In der
% Wissenschaft, insbesondere in der Logik, will man den Dingen auf den
% Grund gehen, will alles genau auseinandernehmen. Da möchte man bestimmte
% Regeln nicht einfach so als gegeben voraussetzen.

\subsection{Junktoren}

Bisher traten zusammengesetzte Aussagen alleinig in Form einer
Implikation auf. Will man logische Zusammenhänge beschreiben können,
muss die logische Sprache um weitere Junktoren bereichert werden.
Unter einem \emph{Junktor}\index{Junktor} versteht man eine logische
Verknüpfung von Aussagen zu einer zusammengesetzten Aussage.

Wir werden einen Junktor durch \emph{Einführungsregeln}%
\index{Einfuehrungsregel@Einführungsregel} und
\emph{Beseitigungsregeln}\index{Beseitigungsregel}
charakterisieren. Die Regeln der Implikation
wurden bereits beschrieben; die Einführung geschieht per
Implikationseinführung, die Beseitigung per Modus ponens.
Für die restlichen Junktoren der Aussagenlogik lassen sich die Regeln
wahlweise in Form von Axiomenschemata oder in Form von Schlussregeln
darstellen. Ich möchte das per Schemata machen, weil diese ein wenig
kompakter sind, was sie hoffentlich ein wenig leichter einsichtig macht.
Die entsprechenden Schlussregeln leiten wir anschließend als zulässige
Regeln ab.

Die Konjunktion\index{Konjunktion} $A\land B$, auch Und"=Verknüpfung
genannt, sprich »$A$ und $B$«, ist charakterisiert durch die Sequenzen
\[A,B\vdash A\land B;\qquad A\land B\vdash A;\qquad A\land B\vdash B.\]
Aus dem Fall von sowohl Regen als auch Schnee ist der Fall von
Schneeregen ableitbar. Aus dem Fall von Schneeregen ist der Fall
von Regen ableitbar. Aus dem Fall von Schneeregen ist der Fall
von Schnee ableitbar. So sind diese Sequenzen zu verstehen.

Die Einführung der Konjunktion geschieht mit der Regel
\[\dfrac{\Gamma\vdash A\qquad\Gamma'\vdash B}{\Gamma,\Gamma'\vdash A\land B}.\]
Denn es findet sich der Beweisbaum:
\[
\infer[\infernote{MP}]{\Gamma,\Gamma'\vdash A\land B}{
  \infer[\infernote{MP}]{\Gamma\vdash B\cond A\land B}{
    \infer[\infernote{Impl-Einf.}]{\vdash A\cond (B\cond A\land B)}{
      \infer[\infernote{Impl-Einf.}]{A\vdash B\cond A\land B}{
        \infer[\infernote{Axiom}]{A,B\vdash A\land B}{}}}
  & \Gamma\vdash A}
& \Gamma'\vdash B}
\]
Es steht MP als Abkürzung für Modus ponens, und Impl-Einf. für
Implikationseinführung. Man schreibt alternativ auch das Kürzel
$\cond$E anstelle Impl-Einf. und das Kürzel $\cond$B anstelle
von MP. Hierbei steht E offenkundig für \emph{Einführung} und
B für \emph{Beseitigung}. Aber Vorsicht, in der englischsprachigen
Literatur sind das I für \emph{introduction} und E für
\emph{elimination}.

Die beiden Regeln zur Beseitigung der Konjunktion sind
\[\dfrac{\Gamma\vdash A\land B}{\Gamma\vdash A},
\qquad\dfrac{\Gamma\vdash A\land B}{\Gamma\vdash B}.\]
Denn es findet sich:
\[
\infer[\infernote{MP}]{\Gamma\vdash A}{
  \infer[\infernote{Impl-Einf.}]{\vdash A\land B\cond A}{
    \infer[\infernote{Axiom}]{A\land B\vdash A}{}}
& \Gamma\vdash A\land B}
\]
Die Disjunktion\index{Disjunktion} $A\lor B$, auch Oder"=Verknüpfung
genannt, sprich »$A$ oder $B$«, ist charakterisiert durch die Sequenzen
\[A\vdash A\lor B;\qquad B\vdash A\lor B;\qquad
A\lor B, (A\cond C), (B\cond C)\vdash C.\]
So ist »Die Erde des Beetes ist nass« ableitbar aus »Es hat geregnet
oder das Beet wurde gegossen«. Denn sowohl »Es hat geregnet« als auch
»Das Beet wurde gegossen« impliziert »Die Erde des Beetes ist nass«.

Die beiden Regeln zur Einführung der Disjunktion sind
\[\dfrac{\Gamma\vdash A}{\Gamma\vdash A\lor B},\qquad
\dfrac{\Gamma\vdash B}{\Gamma\vdash A\lor B}.\]
Die Regel zur Beseitigung der Disjunktion ist
\[\dfrac{\Gamma\vdash A\lor B\qquad\Gamma',A\vdash C\qquad\Gamma'',B\vdash C}
{\Gamma,\Gamma',\Gamma''\vdash C}.\]
Die Beweise dieser Regeln seien dem Leser überlassen.

Eine Aussage wie »Bertram wird seine Hausaufgaben nicht machen«
formuliert man gern in der Form »Wenn Bertram seine Hausaufgaben macht,
färbt sich der Mond grün«. In gleichartiger Weise lässt sich die
Verneinung auch in der formalen Logik definieren. Hierzu legt man als
Hilfsbegriff zunächst $\bot$ als die \emph{Kontradiktion}%
\index{Kontradiktion} fest, sie steht für eine widersprüchliche Formel.

Die Negation\index{Negation} $\lnot A$, auch Verneinung genannt, sprich
»nicht $A$«, definiert man als identisch mit $A\cond\bot$. Hierdurch
sind die Regeln zu ihrer Einführung und Beseitigung auf die der
Implikation zurückführbar. Es ergibt sich
\[\dfrac{\Gamma,A\vdash\bot}{\Gamma\vdash\lnot A},
\qquad\dfrac{\Gamma\vdash\lnot A\qquad\Gamma'\vdash A}
{\Gamma,\Gamma'\vdash\bot}.\]
Alternativ ließe sich die Negation durch die Sequenzen
\[(A\cond\bot)\vdash\lnot A;\qquad A,\lnot A\vdash\bot\]
charakterisieren. Man überzeuge sich, dass dies aufs selbe hinausläuft.

Die Äquivalenz\index{Aequivalenz@Äquivalenz} $A\bicond B$, sprich
»$A$ genau dann, wenn $B$«, definiert man als identisch mit
$(A\cond B)\land (B\cond A)$. Insofern sind die Regeln zu ihrer
Einführung und Beseitigung auf die der Konjunktion zurückführbar.
Es ergibt sich
\[\dfrac{\Gamma\vdash A\cond B\qquad\Gamma'\vdash B\cond A}
{\Gamma,\Gamma'\vdash A\bicond B},\qquad
\dfrac{\Gamma\vdash A\bicond B}{\Gamma\vdash A\cond B},\qquad
\dfrac{\Gamma\vdash A\bicond B}{\Gamma\vdash B\cond A}.\]
Die entsprechenden charakterisierenden Sequenzen sind
\[(A\cond B),(B\cond A)\vdash A\bicond B;\qquad (A\bicond B),A\vdash B;
\qquad (A\bicond B),B\vdash A.\]

\subsection{Zur Syntax}

So wie »Punktrechnung vor Strichrechnung« gilt, legt man für jeden Junktor
zur Einsparung von Klammern eine Stufe der Priorität fest. In
absteigender Rangfolge sind dies ${\lnot}, {\land}, {\lor}, {\cond},
{\bicond}$. So wird die Formel
\[\lnot A\land B\lor C\cond D\quad\text{gelesen als}\quad
(((\lnot A)\land B)\lor C)\cond D.\]
Weiterhin legt man die Implikation als rechtsassoziativ fest. So
wird
\[A\cond B\cond C\quad\text{gelesen als}\quad A\cond (B\cond C).\]
Manche Schüler haben Schwierigkeiten, die Struktur von Termen zu
durchschauen. Infolge kann es bei ihnen zu Flüchtigkeitsfehlern
bei der Ersetzung von Variablen durch Terme kommen. Sie vergessen,
dass ein Term vor der Einfügung zunächst geklammert werden muss.
Erst die Operatorrangfolge gewährt es, die Klammern unter Umständen
nachträglich entfallen zu lassen. Für diese Schüler mag es förderlich
sein, einen Term als \emph{abstrakten Syntaxbaum} darzustellen.
Gleichermaßen verhält es sich mit der Programmiersprache Lisp, die
Terme als Schachtelung von Listen darstellt, deren Klammern
obligatorisch sind.  Die Aussage $A\land B\cond C$ ist beispielsweise
beschreibbar als
\[\text{\texttt{(implies (and A B) C)}}.\]
Im Wesentlichen veranschaulicht diese Schachtelung
nichts anderes als den abstrakten Syntaxbaum. Man kann gewissermaßen
sagen, dass Lisp eine Programmiersprache ohne Syntax ist. Fast ohne,
im höheren Sinne ohne.

% In der Logik versucht man, die Objektsprache genau zu fassen.
% Daher nutzt man das Konzept der formalen Sprache. Terme werden
% in ihr durch Produktionsregeln beschrieben. Produktionsregeln
% lassen unter anderem in Form der EBNF notieren.

Schreibt man viele logische Formeln auf, drängt es, zumindest bei privaten
Notizen und Rechnungen, nach Kurzschreibweisen. In der Logik ist für das
Konditional $A\cond B$ auch die Schreibweise $A\rightarrow B$ gebräuchlich,
für das Bikonditional $A\bicond B$ entsprechend $A\leftrightarrow B$.
Insbesondere in der Schaltalgebra schreibt man auch $\overline A$
anstelle von $\lnot A$, $AB$ anstelle von $A\land B$ und $A+B$ anstelle
von $A\lor B$. Hierbei darf die Disjunktion $A+B$ allerdings nicht mit
der Kontravalenz $A\oplus B$ verwechselt werden.


\section{Natürliches Schließen}

\subsection{Darstellungsformen}

Abgeleitet werden soll das Theorem
\[\vdash (A\cond B)\cond (\lnot B\cond\lnot A).\]
Meine favorisierte und in diesem Buch genutzte Form der Darstellung
des natürlichen Schließens fügt die aus den Schlussregeln erhaltenen
Schlüsse wie Legosteine zu einem Baum zusammen, dem
\emph{Beweisbaum}\index{Beweisbaum} oder \emph{Herleitungsbaum}.
Im Eigentlichen stehen in den Blättern die Grundsequenzen, und in der
Wurzel das Theorem. Wie wir es bereits getan haben, arbeitet man
allerdings auch mit Exemplaren, die irgendwelche Sequenzen in
irgendwelche Sequenzen überführen, womit man zulässige Schlussregeln
erhält. Allgemeiner ginge ferner die Formulierung als gerichteter
azyklischer Graph, die bei einigen Beweisen ein wenig den
Schreibaufwand reduzieren würde.

Der Beweisbaum genannten Theorems ist:
\[
\infer[\infernote{$\cond$E}]{\vdash (A\cond B)\cond (\lnot B\cond\lnot A)}{
  \infer[\infernote{$\cond$E}]{A\cond B\vdash \lnot B\cond\lnot A}{
    \infer[\infernote{$\lnot$E}]{A\cond B, \lnot B\vdash\lnot A}{
      \infer[\infernote{$\lnot$B}]{A\cond B, \lnot B, A\vdash\bot}{
        \infer[\infernote{Axiom}]{\lnot B\vdash\lnot B}{}
      & \infer[\infernote{$\cond$B}]{A\cond B, A\vdash B}{
          \infer[\infernote{Axiom}]{A\cond B\vdash A\cond B}{}
        & \infer[\infernote{Axiom}]{A\vdash A}{}}}}}}
\]
Die Ausformulierung der Sequenzen verlangt langwieriges erneutes
Aufschreiben der Antezedenzen. Sobald man das Prozedere einmal
verstanden hat, erscheint es überausführlich. Man kann sich daher
verkürzte Darstellungen der Beweisbäume überlegen:
\[
\begin{tabular}{@{}l@{\qquad\quad}l}
\infer{\vdash (A\cond B)\cond (\lnot B\cond\lnot A)}{
  \infer{2\vdash \lnot B\cond\lnot A}{
    \infer{1, 2\vdash\lnot A}{
      \infer{1, 2, 3\vdash\bot}{
        \infer{1\equiv\lnot B}{}
      & \infer{2, 3\vdash B}{
          \infer{2\equiv A\cond B}{}
        & \infer{3\equiv A}{}}}}}}
&
\infer[\infernote{$\sim$2}]{(A\cond B)\cond (\lnot B\cond\lnot A)}{
  \infer[\infernote{$\sim$1}]{\lnot B\cond\lnot A}{
    \infer[\infernote{$\sim$3}]{\lnot A}{
      \infer{\bot}{
        \infer[\infernote{1}]{\lnot B}{}
      & \infer{B}{
          \infer[\infernote{2}]{A\cond B}{}
        & \infer[\infernote{3}]{A}{}}}}}}
\end{tabular}
\]
Die linke Form kürzt die Antezedenzen durch Nummern ab. In der rechten
Form entfallen die Antezedenzen vollständig. Stattdessen tauchen sie
als in den Blättern gemachte nummerierte \emph{Annahmen} auf, die im
Fortgang zur Wurzel irgendwann zu tilgen sind. Ihre Tilgung erscheint
nun als Randnotiz.

Eine weitere, sehr systematische Darstellung setzt den Beweis aus
einer Liste von Tabellenzeilen zusammen. Allerdings ist sie ein wenig
mühevoll zu lesen. Jede Zeile enthält eine Aussage und dahinter
zusätzlich die Information, wie und woraus die Aussage abgeleitet wurde.
Jede der Aussagen bekommt eine Nummer, siehe Tabelle \ref{tab:Kontraposition}.
Die Nummerierung der Abhängigkeiten ist in derselben Reihenfolge wie
zuvor bei den Bäumen angegeben. Wer die Liste genauer betrachtet, erkennt,
dass die jeweilige Zeile nichts anderes als die Sequenz
$\text{Abh.}\vdash\text{Nr.}$ darstellt.

\begin{table}
\begin{center}
\caption{Beweis in Form einer Liste von Tabellenzeilen}
\label{tab:Kontraposition}
\begin{tabular}{cclll}
\toprule
\strong{Abh.} & \strong{Nr.} & \strong{Aussage} & \strong{Regel} & \strong{auf}\\
\midrule
1 & 1 & $\lnot B$ & Axiom &\\
2 & 2 & $A\cond B$ & Axiom &\\
3 & 3 & $A$ & Axiom &\\
2, 3 & 4 & $B$ & $\cond$B & 2, 3\\
1, 2, 3 & 5 & $\bot$ & $\lnot$B & 1, 4\\
1, 2 & 6 & $\lnot A$ & $\lnot$E & 5\\
2 & 7 & $\lnot B\cond\lnot A$ & $\cond$E & 6\\
$\emptyset$ & 8 & $(A\cond B)\cond(\lnot B\cond\lnot A)$ & $\cond$E & 7\\
\bottomrule
\end{tabular}
\end{center}
\end{table}

Unabhängig von Gentzen entwickelte Jaśkowski das natürliche Schließen
einige Jahre zuvor. Während Gentzen Beweise als Bäume darstellte, nutzte
Jaśkowski zunächst eine grafische Darstellung, die später von Fitch
adaptiert wurde und in dieser Form nun als \emph{Fitch"=Style}%
\index{Fitch-Style} bekannt ist.

Zu guter Letzt muss die klassische Darstellung der Beweisführung
aufgeführt werden. Die in Worten. Sie zeichnet sich durch die Auslassung
mühseliger technischer Details und blumige Formulierungen aus,
soll aber genug Information enthalten, dass der Leser im Zweifel
eine Formalisierung des Beweises erstellen und verifizieren kann.

\begin{Satz}
Es gilt $(A\cond B)\cond (\lnot B\cond\lnot A)$.
\end{Satz}
\strong{Beweis.} Aus der Annahme von sowohl $A\cond B$ als
auch $\lnot B$ als auch $A$ ist ein Widerspruch abzuleiten.
Man erhält $B$ zunächst per Modus ponens aus $A\cond B$ und $A$.
Nun steht $\lnot B$ bereits im Widerspruch zu $B$.\,\qedsymbol

Als komfortablen Bonus erhält man mit dem Theorem nun im Anschluss
kurzerhand eine weitere zulässige Regel, die
\emph{Kontraposition}\index{Kontraposition}
\[\dfrac{\Gamma\vdash A\cond B}{\Gamma\vdash\lnot B\cond\lnot A},
\quad\text{denn}\quad
\dfrac{\vdash (A\cond B)\cond (\lnot B\cond\lnot A)\quad\Gamma\vdash A\cond B}
{\Gamma\vdash \lnot B\cond\lnot A}.\]
Fügt man ihr den Modus ponens an, findet sich der
\emph{Modus tollens}\index{Modus tollens}
\[\dfrac{\Gamma\vdash A\cond B\qquad\Gamma'\vdash\lnot B}
{\Gamma,\Gamma'\vdash\lnot A}.\]

\subsection{Widerspruchsbeweise}

Beim \emph{Beweis durch Widerspruch} widerlegt man eine Aussage, indem
gezeigt wird, dass die Annahme der Aussage zu einem logischen
Widerspruch führt. In manchen Situationen bietet diese Art der
Argumentation eine große Hilfe. So schreibt der britische Mathematiker
Godfrey Harold Hardy in seinem Essay \emph{A Mathematician's Apology}
die Worte
\begin{quote}
»The proof is by \emph{reductio ad absurdum}, and \emph{reductio ad
absurdum}, which Euclid loved so much, is one of a mathematician's
finest weapons. It is a far finer gambit than any chess gambit: a chess
player may offer the sacrifice of a pawn or even a piece, but a
mathematician offers \emph{the game}.«
\end{quote}
Zur Schaffung von Klarheit muss man zunächst zwei inhaltlich
verschiedene Arten des Widerspruchsbeweises unterscheiden.
Präzisieren lässt sich diese Unterscheidung anhand
der Regeln
\[\dfrac{\Gamma,A\vdash\bot}{\Gamma\vdash\lnot A},\qquad
\dfrac{\Gamma,\lnot A\vdash\bot}{\Gamma\vdash A}.\]
Die linke Regel stellt die stets verfügbare Negationseinführung dar,
die man auch als \emph{Widerlegung durch Widerspruch} bezeichnen kann.
In der rechten Regel, der klassischen \emph{Reductio ad absurdum},
gelangt man zunächst per Negationseinführung von
$\Gamma,\lnot A\vdash\bot$ zu $\Gamma\vdash\lnot\lnot A$, und daraufhin
zu $\Gamma\vdash A$. Die Beseitigung der Doppelnegation ist allerdings
lediglich in der klassischen Logik verfügbar, in der intuitionistischen
gilt sie dagegen als unzulässig.




\chapter{Semantik}

\section{Die klassische Semantik der Aussagenlogik}

\subsection{Die Erfüllungsrelation}

Bislang trat die Logik in der Form eines formalen Systems in
Erscheinung. Gegenstand eines solchen Systems sind im Allgemeinen
\emph{Wörter} einer formalen Sprache; im natürlichen Schließen sind das
die Sequenzen. Einige Wörter, die \emph{Axiome}, werden als gegeben
vorausgesetzt. Unter Anwendung von \emph{Ableitungsregeln}, auch
\emph{Inferenzregeln} genannt, das sind die Schlussregeln, leitet man
aus bereits abgeleiteten Wörtern weitere Wörter der Sprache ab.
In diesem Sinne handelt es sich um ein rein syntaktisches System.

Zum tieferen Verständnis muss man sich im Fortgang damit beschäftigen,
welche inhaltliche Bedeutung den logischen Aussagen beigemessen
wird. Der hierfür wesentliche Schritt besteht in der Definition
einer passenden \emph{Semantik}\index{Semantik}.

Gegenstand der Semantik der Logik ist der Wahrheitsgehalt von Aussagen.
Man hat gefunden, dass es sich mit der Frage nach dem Wesen der Wahrheit
schwierig verhält. Wir wollen daher an dieser Stelle gar nicht erst
versuchen, sie zu ergründen. Stattdessen tritt Wahrheit für uns zunächst
lediglich im leicht fassbaren Rahmen der zweiwertigen booleschen
Algebra auf.

In der klassischen Semantik der Aussagenlogik herrscht das
\emph{Bivalenzprinzip}\index{Bivalenzprinzip}, das besagt, dass jede
Aussage entweder \emph{wahr} oder \emph{falsch} sein muss, also einen
von zwei Wahrheitswerten\index{Wahrheitswert} haben muss. Eine Aussage
kann nicht \emph{ein wenig wahr} oder \emph{halbwegs wahr} sein, noch
kann sie eine von mehreren unterschiedlichen gleichwertigen Wahrheiten
haben. Wir schreiben kurz $0$ für falsch und $1$ für wahr. Enthält eine
Formel logische Variablen, kommt ihr ein Wahrheitswert zu, sobald alle
Variablen durch eine Interpretation\index{Interpretation} mit
einem Wahrheitswert belegt wurden.

Die Art und Weise, wie einer Formel ein Wahrheitswert zukommt,
präzisiert die Erfüllungsrelation. Sie wird als Rekursion über den
Formelaufbau definiert. Der Wahrheitswert einer Formel ist hierbei
einzig und allein durch die Wahrheitswerte ihrer Teilformeln bestimmt.

\begin{Definition}[Erfüllungsrelation]%
\label{def:sat}\newlinefirst
Eine \emph{Interpretation} $I$ ist eine Funktion, die jede atomare
logische Variable $P$ mit einem Wahrheitswert $I(P)\in\{0,1\}$ belegt.
Man definiert $I\models A$, sprich »$I$ erfüllt $A$«, rekursiv als
\begin{align*}
(I\models\bot) &:\Leftrightarrow 0,
& (I\models A\land B) &:\Leftrightarrow ((I\models A)\land (I\models B)),\\
(I\models\top) &:\Leftrightarrow 1,
& (I\models A\lor B) &:\Leftrightarrow ((I\models A)\lor (I\models B)),\\
(I\models P) &:\Leftrightarrow I(P),
& (I\models A\cond B) &:\Leftrightarrow ((I\models A)\cond (I\models B)),\\
(I\models\lnot A) &:\Leftrightarrow \lnot (I\models A),
& (I\models A\bicond B) &:\Leftrightarrow ((I\models A)\bicond (I\models B)).
\end{align*}
\end{Definition}
Die rechte Seite der jeweiligen Festsetzung ist metalogisch zu verstehen
und per Wahrheitstafel definiert, siehe Tabelle \ref{tab:Junktoren}.
Die Schreibweise $I\nvDash A$ ist gleichbedeutend mit
$\lnot (I\models A)$. Eine Interpretation wird auch als \emph{Modell}
bezeichnet. Man nennt sie \emph{Modell} einer Formel, falls sie die
Formel erfüllt. Andernfalls spricht man von einem \emph{Kontramodell}
oder \emph{Gegenmodell} der Formel.

\begin{table}
\caption{Wahrheitstafel der Junktoren}
\label{tab:Junktoren}
\centering
\begin{tabular}{cc@{\qquad}c@{\qquad}c@{\qquad}c@{\qquad}c@{\qquad}c}
\toprule
$A$ & $B$ & $\lnot A$ & $A\land B$ & $A\lor B$ & $A\cond B$ & $A\bicond B$\\
\midrule[\heavyrulewidth]
$0$ & $0$ & $1$ & $0$ & $0$ & $1$ & $1$\\
$1$ & $0$ & $0$ & $0$ & $1$ & $0$ & $0$\\
$0$ & $1$ & $1$ & $0$ & $1$ & $1$ & $0$\\
$1$ & $1$ & $0$ & $1$ & $1$ & $1$ & $1$\\
\bottomrule
\end{tabular}
\end{table}

\begin{Definition}\label{def:sat-context}
Für einen Kontext $\Gamma = \{A_1,\ldots,A_n\}$ setzt man
\[(I\models\Gamma)\,:\bicond\, (I\models A_1)\land\ldots\land (I\models A_n).\]
\end{Definition}

\subsection{Gültigkeit einer Formel}

Eine wichtige Rolle spielen \emph{allgemeingültige} Formeln, die man
in der Aussagenlogik auch als \emph{Tautologien}\index{Tautologie}
bezeichnet. Sie sind immer wahr, unabhängig davon, mit welchem
Wahrheitswert ihre logischen Variablen belegt werden.

Als allgemeinere Begrifflichkeit wollen wir auf einen Kontext $\Gamma$
bezogen gültige Formeln $A$ betrachten. Die Idee hierbei ist,
dass wenn die Formeln des Kontextes als wahr angenommen werden, die
Formel $A$ ebenfalls wahr sein muss. Trifft dies auf $A$ zu,
schreibt man $\Gamma\models A$, gelesen »im Kontext $\Gamma$ ist
$A$ gültig«, oder auch »$\Gamma$ zieht $A$ nach sich«. Die Bezeichnung
\emph{logische Folgerung} oder \emph{logische Konsequenz} ist
ebenfalls verbreitet.

\begin{Definition}[Gültige Formel]\label{def:valid}\newlinefirst
Eine Formel $A$ heißt \emph{gültig} im Kontext $\Gamma$,
wenn jede Interpretation, die sämtliche Formeln von $\Gamma$ erfüllt,
auch $A$ erfüllt. In metalogischer Symbolik,
\[(\Gamma\models A)\,:\bicond\,\forall I\colon (I\models\Gamma)\cond (I\models A).\]
\end{Definition}
Eine im leeren Kontext gültige aussagenlogische Formel $A$ nennt man
wie gesagt Tautologie. Statt $\emptyset\models A$ schreibt man auch
kurz $\models A$. Wie bei Sequenzen schreibt man auch $\Gamma,A,B\models C$
statt $\Gamma\cup\{A,B\}\models C$.

\subsection{Wahrheitstafeln}

Obgleich der Variablenvorrat unendlich groß sein darf, enthält eine
Formel von den Variablen nur endlich viele. Insofern sind für eine
Formel in einem Kontext auch nur endlich viele Interpretationen relevant.
Sind insgesamt $n$ Variablen vorhanden, sind es $2^n$ Interpretationen.

Eine Interpretation $I$ mit der Auswertung $I\models A$ ist nichts
anderes als eine Zeile der Wahrheitstafel\index{Wahrheitstafel} der
Formel $A$. Eine Formel ist genau dann tautologisch, wenn in der
Ergebnisspalte in jeder Zeile eine~1 steht.

Ein Beispiel. Die Wahrheitstafel \ref{tab:Tautologie-zur-Kontraposition}
bestätigt
\[\models (a\cond b)\bicond (\lnot b\cond\lnot a).\]

\begin{table}
\caption{Wahrheitstafel der Tautologie zur Kontraposition}
\label{tab:Tautologie-zur-Kontraposition}
\centering
\begin{tabular}{cc@{\quad\;\;}c@{\quad\;\;}c@{\quad\;\;}c@{\quad\;\;}c@{\quad\;\;}c}
\toprule
$a$ & $b$ & $\lnot b$ & $\lnot a$ & $a\cond b$ & $\lnot b\cond\lnot a$
& $(a\cond b)\bicond (\lnot b\cond\lnot a)$\\
\midrule[\heavyrulewidth]
$0$ & $0$ & $1$ & $1$ & $1$ & $1$ & $1$ \\
$1$ & $0$ & $1$ & $0$ & $0$ & $0$ & $1$ \\
$0$ & $1$ & $0$ & $1$ & $1$ & $1$ & $1$ \\
$1$ & $1$ & $0$ & $0$ & $1$ & $1$ & $1$ \\
\bottomrule
\end{tabular}
\end{table}

\noindent
Die Tafel führt zusätzlich die Teilformeln auf, was bei längeren
Formeln recht mühselig erscheinen mag. Eine geschickte Methode zur
Reduzierung des Schreibaufwands erspart die Teilformeln, und setzt ihre
Wahrheitswerte dafür schlicht unter die Junktoren, denn Ziffern
benötigen nicht viel Platz.

Die Prüfung einer logischen Folgerung per Wahrheitstafel
wird ermöglicht durch die metalogische Beziehung
\[(A_1,\ldots,A_n\models A)\,\bicond\, (\models A_1\land\ldots\land A_n\cond A).\]
Nämlich findet sich die äquivalente Umformung
\begin{align*}
(A_1,\ldots,A_n\models A)
&\;\Leftrightarrow_{\text{(1)}}\;
(\forall I\colon (I\models A_1)\land\ldots\land (I\models A_n)\cond (I\models A))\\
&\;\Leftrightarrow_{\text{(2)}}\;
(\forall I\colon (I\models A_1\land\ldots\land A_n)\cond (I\models A))\\
&\;\Leftrightarrow_{\text{(3)}}\;
(\forall I\colon I\models A_1\land\ldots\land A_n\cond A)\\
&\;\Leftrightarrow_{\text{(4)}}\;
(\models A_1\land\ldots\land A_n\cond A).
\end{align*}
Hierbei gilt (1), (4) gemäß Def. \ref{def:valid}, \ref{def:sat-context}
und (2), (3) gemäß Def. \ref{def:sat}.

\subsection{Korrektheit des natürlichen Schließens}

Jede Schlussregel und jedes Axiom besitzt eine semantische Entsprechung.
Die Beweise dafür werden nun auf metalogischer Ebene mittels natürlichem
Schließen selbst erbracht, was wie eine Art Zirkelschluss erscheinen
mag.

Zu den Grundsequenzen und zur Abschwächungsregel findet sich:
\[
\begin{array}{@{}l@{\qquad}l}
\begin{prooftree}[center = false, separation = 1em]
        \infer0[1]{I\models\Gamma\cup\{A\}}
      \infer1[Def. \ref{def:sat-context}]{(I\models\Gamma)\land (I\models A)}
    \infer1{I\models A}
  \infer1[$\sim$1]{(I\models \Gamma\cup\{A\})\cond (I\models A)}
\infer1[Def. \ref{def:valid}]{\Gamma\cup\{A\}\models A}
\end{prooftree}
&
\begin{prooftree}[center = false, separation = 1em]
        \hypo{\Gamma\models A}
      \infer1[Def. \ref{def:valid}]{(I\models\Gamma)\cond (I\models A)}
        \infer0[1]{I\models\Gamma\cup\Gamma'}
      \infer1[Def. \ref{def:sat-context}]{I\models\Gamma}
    \infer2{I\models A}
  \infer1[$\sim$1]{(I\models\Gamma\cup\Gamma')\cond (I\models A)}
\infer1[Def. \ref{def:valid}]{\Gamma\cup\Gamma'\models A}
\end{prooftree}
\end{array}
\]
Die Prüfung der restlichen Entsprechungen sei dem Leser überlassen.

\begin{Satz}[Korrektheit des natürlichen Schließens]\newlinefirst
Ist die Sequenz $\Gamma\vdash A$ ableitbar, so muss auch
$\Gamma\models A$ gelten.
\end{Satz}
\begin{Beweis}
Strukturelle Induktion über die Konstruktion von Beweisbäumen.
Induktionsanfänge sind die semantischen Entsprechungen
der Grundsequenzen. Induktionsschritte sind die semantischen
Entsprechungen der Schlussregeln. Die Beweise der Entsprechungen
wurden bereits diskutiert.\,\qedsymbol
\end{Beweis}

% #todo
% Korrektheit verschafft noch mehr. In Kontraposition,
% wenn ein Kontramodell $\Gamma\models A$ widerlegt,
% ist die Sequenz $\Gamma\vdash A$ nicht ableitbar.
% Solche Negativresultate sind ohne Semantik schwieriger
% zu erbringen.

\newpage
\subsection{Logische Äquivalenz}
\begin{Definition}[Äquivalente Formeln]\newlinefirst
Die Äquivalenz zweier Formeln $A,B$ ist definiert als
\[(A\equiv B)\,:\bicond\, (\models A\bicond B).\]
\end{Definition}

\noindent
Eine Äquivalenz besteht genau dann, wenn jede der beiden Formeln
eine logische Folgerung der anderen ist. Das heißt, es besteht die
metalogische Beziehung
\[(\models A\bicond B)\,\bicond\, (A\models B)\land (B\models A).\]
Mit den semantischen Entsprechungen der Schlussregeln
findet sich nämlich:
\[
\begin{array}{@{}l@{\qquad\quad}l}
\begin{prooftree}[center = false]
    \hypo{\models A\bicond B}
  \infer1{\models A\cond B}
  \infer0{A\models A}
\infer2{A\models B}
\end{prooftree}
&
\begin{prooftree}[center = false]
    \hypo{A\models B}
  \infer1{\models A\cond B}
    \hypo{B\models A}
  \infer1{\models B\cond A}
\infer2{\models A\bicond B}
\end{prooftree}
\end{array}
\]
\begin{Satz}\label{log-Aeq-ist-Aeqrel}
Es ist $A\equiv B$ eine Äquivalenzrelation. Das heißt, es gilt
\begin{align*}
& A\equiv A, &&\text{(Reflexivität)}\\
& (A\equiv B)\cond (B\equiv A), &&\text{(Symmetrie)}\\
& (A\equiv B)\land (B\equiv C)\cond (A\equiv C). &&\text{(Transitivität)}
\end{align*}
\end{Satz}
Der Beweis sei dem Leser als kleine Übung überlassen.

Der Satz \ref{log-Aeq-ist-Aeqrel} vermittelt, dass Formeln mit
Äquivalenzen so umgeformt werden dürfen, wie Terme mit Termumformungen.
Es finden sich im Fortgang eine Reihe von grundlegenden Äquivalenzen,
die Regeln der \emph{booleschen Algebra}\index{boolesche Algebra}. Sie
wurde erstmals in der Mitte des 19. Jahrhunders vom britischen
Mathematiker George Boole in seiner Abfassung \emph{The Mathematical
Analysis of Logic} und seinem späteren Buch \emph{An Investigation of
The Laws of Thought} beschrieben. Boole beschreibt allerdings, anders
als heute üblich, eine Einbettung der logischen Operationen in die
gewöhnliche Algebra, siehe Tabelle \ref{tab:Boole}.

\begin{table}
\centering
\caption{Einbettung in die gewöhnliche Algebra}
\label{tab:Boole}
\begin{tabular}{ccccc}
\toprule
\strong{Modern} & $\lnot a$ & $a\land b$ & $a\lor b$ & $a\cond b$\\
\strong{Boole} & $1-a$ & $ab$ & $a+b(1-a)$ & $1-a+ab$\\
\bottomrule
\end{tabular}
\end{table}

Logische Äquivalenz im absoluten Sinne ist nicht unter allen Umständen
der Weisheit letzter Schluss. In der Schaltalgebra tun sich
Problemstellungen auf, wo der Wahrheitswert einer Formel $A$ in
einzelnen Zeilen der Wahrheitstafel keine Rolle spielt, man spricht von
\emph{Don't-Care"=Zellen}. Es sei $X$ eine Formel, die genau in diesen
Zeilen wahr ist, in allen anderen falsch. Man möchte $A$ nun
beispielsweise zu $B$ vereinfachen. Unter normalen Umständen sollte
diese Vereinfachung die Äquivalenz $A\equiv B$ einhalten. Bei
Vorhandensein der irrelevanten Zellen genügt jedoch die weniger strenge
Forderung
\[\lnot X\models A\bicond B.\]
Wir haben es hier mit einer relativen Äquivalenz zu tun. Auch bei ihr
handelt es sich um eine Äquivalenzrelation, wobei $X$ beliebig ist,
aber fest sein muss.

\begin{table}
\begin{center}
\caption{Die Regeln der booleschen Algebra.}
\label{tab:boolesche-Algebra}
\begin{tabular}{@{}c@{\qquad}c@{\qquad}l@{}}
\toprule
\strong{Konjunktion}&
\strong{Disjunktion}&
\strong{Bezeichnung}\\
\midrule[\heavyrulewidth]
$A\land 0\equiv 0$ &
$A\lor 1\equiv 1$ &
Extremalgesetze\\

$A\land\lnot A\equiv 0$ &
$A\lor\lnot A\equiv 1$ &
Komplementärgesetze\\

$A\land A\equiv A$ &
$A\lor A\equiv A$ &
Idempotenzgesetze\\

$A\land 1\equiv A$ &
$A\lor 0\equiv A$ &
Neutralitätsgesetze\\
\midrule
$A\land B\equiv B\land A$ &
$A\lor B\equiv B\lor A$ &
Kommutativgesetze\\

$A{\land}(B{\land}C)\equiv (A{\land}B){\land}C$ &
$A{\lor}(B{\lor}C)\equiv (A{\lor}B){\lor}C$ &
Assoziativgesetze\\

$\lnot (A\land B)\equiv\lnot A\lor\lnot B$ &
$\lnot (A\lor B)\equiv\lnot A\land\lnot B$ &
De morgansche Gesetze\\

$A\land (A\lor B)\equiv A$ &
$A\lor (A\land B)\equiv A$ &
Absorptionsgesetze\\
\bottomrule
\end{tabular}
\end{center}
\end{table}

\subsection{Die Einsetzungsregel}

\begin{Satz}[Einsetzungsregel]\newlinefirst
Ist die Formel $A$ allgemeingültig, so führt die simultane Ersetzung
einiger atomarer Variablen durch Formeln bei ihr zu einer weiteren
allgemeingültigen Formel. Metalogisch
\[(\models A)\cond (\models A[P_1:=B_1,\ldots,P_n:=B_n]).\]
\end{Satz}
\begin{Beweis}
Die Bestimmung von $I\models A[\ldots]$ läuft in
derselben Weise ab wie die von $I\models A$, außer dass $I(P_k)$ durch
$I\models B_k$ zu ersetzen ist. Sofern $A$ allgemeingültig ist, gilt
$I\models A$ unabhängig davon, ob $I(P_k)$ wahr oder falsch ist.
Ergo muss $I\models A[\ldots]$ unabhängig davon gelten, ob $I\models B_k$
wahr oder falsch ist.\,\qedsymbol
\end{Beweis}

\noindent
Ich mag daran erinnern, dass bei der Substitution \emph{jedes}
Vorkommen der Variable durch dieselbe Formel zu ersetzen ist.

Zum Beispiel erhält man zu der simultanen Ersetzung
$a:=A$ und $b:=B$ aus
\[\models (a\cond b)\bicond (\lnot b\cond\lnot a)\;\;\text{das Schema}\;\;
\models (A\cond B)\bicond (\lnot B\cond\lnot A).\]
Mit dem Vollständigkeitssatz infolge das Theoremschema
\[\vdash (A\cond B)\bicond (\lnot B\cond\lnot A).\]
Hiermit gewinnt man kurzum die Regeln
\[\dfrac{\Gamma\vdash A\cond B}{\Gamma\vdash\lnot B\cond\lnot A},\qquad
\dfrac{\Gamma\vdash\lnot B\cond\lnot A}{\Gamma\vdash A\cond B}.\]
Die Vorgehensweise ermöglicht summa summarum die Auffindung von
zulässigen Schlussregeln durch geistloses Ausfüllen von Wahrheitstafeln,
was allerdings auf die klassische Aussagenlogik beschränkt bleibt.

\subsection{Wahrheitsfunktionen}

Mit der Bindung ihrer atomaren Variablen gehen aus den Formeln der
Aussagenlogik Wahrheitsfunktionen\index{Wahrheitsfunktion} hervor.
Die Disjunktion zweier Aussagen wird zum Beispiel vermittelt durch die Funktion%
\[f\colon\{0,1\}\times\{0,1\}\to\{0,1\},\quad f(a,b):=(a\lor b).\]
Jede Wahrheitsfunktion in $n$ Variablen ist durch ihre Wahrheitstafel
charakterisiert, die aus $2^n$ Zeilen besteht, da der Definitionsbereich
so viele unterschiedliche Tupel enthält. Ein Tupel von Wahrheitswerten
wird auch als Bitfolge betrachtet.

Zwei Formeln $A,B$ sind genau dann äquivalent, wenn sie durch dieselben
Interpretationen erfüllt werden, sich also in der Wahrheitstafel gleich
verhalten. Man überzeugt sich davon unschwer mit der metalogischen
Umformung%
\begin{align*}
(A\equiv B) &\iff (\models A\bicond B)
\iff (\forall I\colon I\models A\bicond B)\\
&\iff (\forall I\colon (I\models A)\bicond (I\models B)).
\end{align*}
Demnach charakterisieren äquivalente Formeln dieselbe
Wahrheitsfunktion. Man kann dies auch so betrachten, dass die
Wahrheitsfunktion die Äquivalenzklasse all ihrer Formeln repräsentiert.
Unter den Formeln gibt es nun einen besonderen Vertreter, die
\emph{disjunktive Normalform}\index{disjunktive Normalform}, kurz DNF,
die eigentlich nichts anderes als eine direkte Kodierung der
Wahrheitstafel ist. Insofern ließt sich an der Wahrheitstafel
unmittelbar die DNF der Formel ab.

Die Normalform einer Formel lässt sich unter Umständen
vereinfachen. Bei der DNF der Disjunktion ist bereits klar, welche
Formel Resultat der Vereinfachung sein müsste. Es findet sich
\begin{gather*}
(a\land\lnot b)\lor (\lnot a\land b)\lor (a\land b)
\equiv (a\land\lnot b)\lor ((\lnot a\lor a)\land b)\\
\equiv (a\land\lnot b)\lor (1\land b)
\equiv (a\land\lnot b)\lor b
\equiv (a\lor b)\land (\lnot b\lor b)\\
\equiv (a\lor b)\land 1 \equiv a\lor b.
\end{gather*}
In der Schaltalgebra ist die Vereinfachung der DNF von großer
Wichtigkeit, da sie den maßgeblichen Schritt zur Ermittelung von
Schaltungen mit einer minimalen Zahl von Gattern darstellt.
Aus dieser Anforderung heraus wurden systematische Verfahren
zur Vereinfachung entwickelt. Für wenige Variablen stellt das
sogenannte Karnaugh"=Veitch"=Diagramm ein geschicktes
Hilfsmittel dar.

Mit dem Verfahren nach Quine und McCluskey lassen sich Formeln mit
beliebig vielen Variablen mit einem Computer automatisch vereinfachen.
Der Algorithmus ist allerdings von exponentieller Laufzeit, dessen
Ausführung also von einer kombinatorischen Explosion
überschattet. Eine wesentliche Verbesserung darf man auch nicht erwarten,
da die Problemstellung der Vereinfachung als NP"=vollständig befunden
wurde. Für die Formeln mit sehr vielen Variablen liegt das Wissen über
die beste Vereinfachung somit im Schleier der Dunkelheit.

\subsection{Logisches Schließen vermittels boolescher Algebra}

Boole gelang es, vermittels seiner Algebra Schlussfolgerungen zu führen.
Es stellt sich hierbei die Frage, wie dies vonstattengeht, genauer,
in welcher Verbindung die boolesche Algebra zum natürlichen Schließen
steht. Die Herleitung ihrer Regeln stellt mehr eine kleine Übung dar,
als dass sie Schwierigkeiten bereitet. Uns interessiert daher die
Umkehrung, also ob und wie natürliches Schließen aus der booleschen
Algebra heraus gewonnen werden kann. Den Übergang schafft die folgende
Überlegung.

Insofern wir uns mit der booleschen Algebra auf der semantischen Ebene
bewegen, werden wir uns auf die semantische Folgerung statt der
syntaktischen Herleitbarkeit beziehen. Wir stellen $\models A$ hierbei
in der Form $A\equiv\top$ dar. Hierzu will ich unbedingt noch ein paar
Bemerkungen machen. Die semantische Folgerung $A\models B$ mag als
eine Art Ordnung gedeutet werden, wonach $A$ als kleiner oder gleich groß
$B$ angesehen wird. Allerdings handelt es sich lediglich um eine
Quasiordnung, nicht aber um eine Halbordnung, weil zwei Formeln nicht
gleich sein müssen, nur weil sie semantisch gegenseitig auseinander
folgen. Da $\equiv$ aber eine Äquivalenzrelation ist, kann die
Quotientenmenge bezüglich dieser gebildet werden.
Zwei äquivalente Formeln werden dahingehend als gleich
angesehen, wobei dies in gewisser Weise seltsam erscheinen mag, da
die Formeln durch diesen Schritt immerhin ihre eigentliche Bedeutung
verlieren. Die als gleichbedeutend mit $A\models B$ definierte
Ordnung $[A]\le [B]$ stellt nun wirklich eine Halbordnung dar,
wonach $[A]=[B]$, also $A\equiv B$, bereits aus $[A]\le [B]$ zusammen
mit $[B]\le [A]$ folgt. Mit $1:=[\top]$ ist $[A]\le 1$ allgemeingültig,
also $1$ das größte Element, wonach $[A]=1$ und $1\le [A]$ äquivalent sind.
Außerdem nimmt die Quotientenmenge die Gestalt einer algebraischen
Struktur an, indem%
\[\lnot [A] := [\lnot A],\quad [A]\land [B] := [A\land B]\]
usw. erklärt wird. Laut der Ersetzungsregel sind diese wohldefiniert,
womit $\equiv$ die Rolle einer Kongruenzrelation einnimmt.

Besteht über die Semantik noch Unklarheit, kann man die gerade aufgezeigte
Konstruktion auch über die Ableitbarkeitsrelation statt der semantischen
Folgerung vornehmen. Hierbei spricht man, sofern sie aufgeht, von der
\emph{Lindenbaum"=Tarski"=Algebra}\index{Lindenbaum-Tarski-Algebra} des
jeweiligen logischen Systems. Während bei der klassischen Aussagenlogik
tatsächlich eine boolesche Algebra herauskommt, führt die intuitionistische
zu einer Heyting"=Algebra.

Es spielt keine wesentliche Rolle, ob man bezüglich $\equiv$ mit den
Formeln, oder bezüglich $=$ mit den Klassen rechnet. Es bleibt trotzdem
von Bedeutung, auf welche Formeln wir uns beziehen. Zwei unterschiedliche
äquivalente Formeln entsprechen dabei direkt zwei unterschiedlichen
Darstellungen derselben Klasse. Beim üblichen Rechnen mit Termen
werden zwei unterschiedliche Darstellungen desselben Elements ja auch
auseinandergehalten.

Der Modus ponens nimmt nun bzgl. $a:=[A]$ und $b:=[B]$ die Form%
\[\dfrac{(a\cond b)=1\quad\;\; a=1}{b=1}\]
an. Gesucht ist eine Herleitung der Konklusion aus den beiden Prämissen,
bei der ausschließlich boolesche Algebra zur Anwendung kommt. Keine
schwierige Aufgabe. Mit $a=1$ gilt die Umformung%
\[1 = (a\cond b) = (\lnot a\lor b) = (\lnot 1\lor b) = (0\lor b) = b.\]
Die Konklusion $b=1$ konnte also tatsächlich \emph{ausgerechnet} werden.
Diese Vorgehensweise müssen wir noch auf die allgemeine Form des Modus
ponens und der restlichen Schlussregeln übertragen.

Die Abhängigkeit von Annahmen lässt sich folgendermaßen beschreiben.
Wir fassen die Annahmen $H_1,\ldots,H_n$ zunächst zur Hypothese
$H:=(H_1\land\ldots\land H_n)$ zusammen. Die Folgerung
$H_1,\ldots,H_n\models A$ ist daraufhin bzgl. $h:=[H]$ mit $h\le a$
gleichbedeutend. Die allgemeine Form des Modus ponens nimmt
diesbezüglich die Gestalt
\[\dfrac{h\le (a\cond b)\quad\;\; h\le a}{h\le b}\]
an. Mit der Einsicht, dass $h\le a$ zu $(h\land a)=h$ äquivalent ist,
findet sich zunächst die Umformung%
\[(h\land\lnot a) = (h\land a\land\lnot a) = (h\land 0) = 0.\]
Weil auch $h\le (a\cond b)$ zu $(h\land (a\cond b))=h$ äquivalent ist,
gelingt daraufhin die Umformung%
\begin{align*}
h &= (h\land (a\cond b)) = (h\land (\lnot a\lor b))
= ((h\land\lnot a)\lor (h\land b))\\
&= (0\lor (h\land b)) = (h\land b).
\end{align*}
Und die Gleichung $h = (h\land b)$ ist schließlich wieder zur
gesuchten Konklusion $h\le b$ äquivalent.

Die Subjunktionseinführung nimmt die Gestalt
\[\dfrac{(h\land a)\le b}{h\le (a\cond b)}\]
an. Zu ihrer Bestätigung wendet man dieses Mal die allgemeine Äquivalenz
von $x\le y$ und $(x\lor y) = y$ auf die Prämisse an, was zu
$((h\land a)\lor b) = b$ führt. Hiermit findet sich die
Umformung%
\begin{align*}
(h\land (a\cond b)) &= (h\land (\lnot a\lor b))
= (h\land (\lnot a\lor (h\land a)\lor b))\\
&= ((h\land\lnot a)\lor (h\land h\land a)\lor (h\land b))\\
&= (h\land (\lnot a\lor a)\lor (h\land b))
= (h\lor (h\land b)) = h.
\end{align*}
Und diese Gleichung ist zur Konklusion äquivalent. Ich denke, die Bestätigung
der restlichen Schlussregeln brauche ich nicht weiter ausführen, -- dies
stellt nunmehr eine Übung dar. Bemerkt werden darf aber noch, dass $h\le a$
auch zu $(\lnot h\lor a)=1$ äquivalent ist. Anders als bei $h\land a=h$
tauchen $h,a$ darin nur einmal auf. Jedenfalls lassen sich alle Prämissen
und Konklusionen auch immer als Gleichungen darstellen. Und so ergibt sich,
dass das natürliche Schließen unter den gemachten Beziehungen allein aus
dem Rechnen mit Gleichungen vermittels der Gesetze der booleschen
Algebra entspringt.

\newpage
\subsection{Vollständigkeit des natürlichen Schließens}

Der Beweis der Korrektheit eines vorgelegten Kalküls versichert,
dass jede Schlussregel und jedes Axiom die inhaltliche Wahrheit
bewahrt. Eine Aussage als wahr zu befinden bezieht sich aber immer
auf eine bestimmte, zuvor festgelete Semantik. Woher wissen wir denn
eigentlich, dass diese Semantik die Logik adäquat beschreibt? Es wäre
eine Semantik denkbar, die viel mehr Aussagen als wahr befindet, als
aus den Regeln und Axiomen folgen, was Zweifel an der Semantik sähte.
Was verschafft uns also die Sicherheit, die »richtige« Semantik in den
Händen zu halten? Oder aber die Semantik erscheint nur zu weitgehend,
weil nicht bemerkt wurde, dass der Kalkül schlicht zu schwach ist.
Tappen wir im Dunkeln?

Idealerweise sollte jede wahre Aussage auch ableitbar sein, was als
\emph{Vollständigkeit} des Kalküls bezeichnet wird. Das heißt, aus
$\Gamma\models A$ sollte $\Gamma\vdash_S A$ folgen. Zusammen mit der
Korrektheit stellt sich damit nämlich eine Charakterisierung des
Kalküls bzw. der Logik her, dergestalt dass $\Gamma\vdash_S A$ als zu
$\Gamma\models A$ äquivalent befunden wird. Zumindest weiß man dann,
dass ein bestimmter Kalkül und eine bestimmte Semantik zusammen gehören
bzw. einander entsprechen. Man mag daraufhin weiter argwöhnisch fragen,
ob das gefundene Duo aus Kalkül und Semantik die wirklichen Umstände
adäquat erfasst, wie auch immer diese Wirklichkeit genau beschaffen
sei. In der Tat existieren unterschiedliche Logiken, und welche davon
geeignet ist, wird letztlich Gegenstand von pragmatischen, über
beweistheoretische, bishin zu epistemiologischen Erwägungen, siehe bspw.
\cite{George}.

Im Folgenden wird die Vollständigkeit des natürlichen Schließens für
die klassische Aussagenlogik bezüglich der klassischen Semantik gezeigt.
Wir benötigen ein paar Hilfssätze.

Eine Formel $A$ wollen wir \emph{erfüllbar} nennen, wenn eine
Interpretation $I$ mit $I\models A$ existiert, andernfalls
\emph{unerfüllbar}. Entsprechend heiße eine Formelmenge $\Gamma$
\emph{erfüllbar}, wenn $I$ mit $I\models\Gamma$ existiert.

\begin{Satz}\label{gueltig-unerf}
Gilt $\Gamma\models A$, so ist $\Gamma\cup\{\lnot A\}$ unerfüllbar.
\end{Satz}
\begin{Beweis}
Es gelte $\Gamma\models A$. Also folgt $I\models A$ aus $I\models\Gamma$
für jede Interpretation $I$. Angenommen, es existierte eine
Interpretation $J$ mit $J\models\Gamma\cup\{\lnot A\}$, also
$J\models\Gamma$ und $J\models\lnot A$. Mit der Spezialisierung $I:=J$
folgern wir $J\models A$ aus $J\models\Gamma$. Nun bedeutet
$J\models\lnot A$ definitionsgemäß $\lnot(J\models A)$, was im
Widerspruch zu $J\models A$ steht.\,\qedsymbol
\end{Beweis}

\noindent
Die Umkehrung dieses Sachverhaltes bestätigt sich ebenfalls leicht,
wird im Weiteren aber nicht benötigt.

Das syntaktische Analogon der Erfüllbarkeit ist die Konsistenz.
Eine Formelmenge $\Gamma$ heiße \emph{inkonsistent}, wenn
$\Gamma\vdash_S\bot$, andernfalls \emph{konsistent}. Zumindest sollte
speziell $\Gamma:=\emptyset$ konsistent sein, was nichts anderes als die
Widerspruchsfreiheit des Kalküls bedeutet. Andernfalls ließe sich
vermittels ex falso quodlibet jede beliebige Aussage ableiten, was den
Kalkül sinnlos machte.

\begin{Satz}\label{inkon-ableitbar}
Ist $\Gamma\cup\{\lnot A\}$ inkonsistent, so gilt $\Gamma\vdash_{\mathrm C} A$.
\end{Satz}
\begin{Beweis}
Es gelte $\Gamma\cup\{\lnot A\}\vdash_{\mathrm C}\bot$. Somit folgt
$\Gamma\vdash_{\mathrm C}\lnot\lnot A$ via Implikationseinführung beim
natürlichen Schließen bzw. Deduktionstheorem beim Hilbertkalkül. Mit der
Beseitigung der Doppelnegation ergibt sich daraufhin
$\Gamma\vdash_{\mathrm C} A$.\,\qedsymbol
\end{Beweis}

\noindent
Man beachte, dass der Beweis die klassische Logik als Gegenstand
voraussetzt. Zur intuitionistischen Logik ist dieser Beweis nicht
durchführbar.

Der nächste Satz stellt eher ein allgemeines Resultat der Mengenlehre
dar. Der Leser mag den Beweis gerne zunächst überspringen, und
nach Bewältigung des Kapitels über Mengenlehre hierher zurückkommen.

\begin{Satz}\label{endliche-Teilmenge-von-lim-sup}
Die Mengenfolge $(\Gamma_i)_{i\in\N}$ sei monoton steigend.
Sei $\Gamma':=\bigcup_{i\in\N}\Gamma_i$. Ist $\Delta\subseteq\Gamma'$
endlich, dann existiert $j\in\N$ mit $\Delta\subseteq\Gamma_j$.
\end{Satz}
\begin{Beweis}
Für $\Delta=\emptyset$ ist der Beweis trivial. Wir setzen im Weiteren
also $\Delta\ne\emptyset$ voraus. Aufgrund der Endlichkeit
besitzt $\Delta$ die Form $\Delta = \{A_1,\ldots,A_n\}$. Wegen
$\Delta\subseteq\Gamma'$ gilt $A_k\in\Gamma'$ zu jedem der $k$. Folglich
existiert jeweils ein Index $i(k)$ mit $A_k\in\Gamma_{i(k)}$, denn
andernfalls läge $A_k$ nicht in der Vereinigung $\Gamma'$. Sei $j$ das
Maximum der Indizes $i(1)$ bis $i(n)$. Sei nun $A\in\Delta$ fest,
aber beliebig. Also ist $A=A_k$ für ein $k$, und somit
$A\in\Gamma_{i(k)}$. Wir haben $i(k)\le j$, daher $\Gamma_{i(k)}\subseteq\Gamma_j$
aufgrund der Monotonie, womit $A\in\Gamma_j$. Ergo haben wir ein $j$
mit $\Delta\subseteq\Gamma_j$ konstruiert.\,\qedsymbol
\end{Beweis}

\begin{Satz}\label{Formeln-abzaehlbar}
Die Menge der logischen Formeln ist abzählbar.
\end{Satz}
\begin{Beweis}
Jede Formel lässt sich auf die übliche Art als endliche Zeichenkette
darstellen. Wir ordnen daher die Zeichenketten nach ihrer Länge. Zu
jeder Länge gibt es nur endlich viele Zeichenketten, aus denen die
wohlgeformten Formeln extrahiert und lexikographisch geordnet werden.
Somit erhält man eine Abzählung aller Formeln, die darüber hinaus
bijektiv ist.\,\qedsymbol
\end{Beweis}

\begin{Definition}
Eine Formelmenge $\Gamma$ heißt \emph{maximalkonsistent}, wenn sie
konsistent ist, aber $\Gamma\cup\{A\}$ inkonsistent für jede Formel
$A\notin\Gamma$.
\end{Definition}

\noindent
Der nächste Satz versichert, dass sich eine konsistente Formelmenge
stets zu einer maximalkonsistenten erweitern lässt.

\begin{Satz}[Lemma von Lindenbaum]\label{Lemma-Lindenbaum}\newlinefirst
Zu jedem konsistenten $\Gamma$ existiert ein maximalkonsitentes
$\Gamma'$ mit $\Gamma\subseteq\Gamma'$.
\end{Satz}
\begin{Beweis}
Laut Satz \ref{Formeln-abzaehlbar} ist die Menge sämtlicher Formeln
abzählbar. Somit liegt eine Abzählung $(A_i)_{i\in\N}$ vor, die jede
Formel mindestens einmal aufzählt. Die Folge $(\Gamma_i)_{i\in\N}$
legen wir rekursiv fest als
\[\Gamma_0 := \Gamma,\qquad \Gamma_{i+1} := \begin{cases}
\Gamma_i\cup\{A_i\}, & \text{falls konsistent},\\
\Gamma_i, & \text{sonst}.
\end{cases}\]
Sei nun $\Gamma':=\bigcup_{i\in\N}\Gamma_i$. Zunächst ist jedes
$\Gamma_i$ konsistent, weil die Folge genau auf diese Weise erklärt
wurde. Angenommen, $\Gamma'$ wäre inkonsistent, dann existierte aufgrund
der Kompaktheit eine endliche Teilmenge $\Delta\subseteq\Gamma'$, die
bereits inkonsistent ist. Zu dieser existiert laut Satz
\ref{endliche-Teilmenge-von-lim-sup} ein $j$ mit $\Delta\subseteq\Gamma_j$,
was aber absurd ist, weil folglich $\Gamma_j$ erst recht inkonsistent wäre.
Ergo muss $\Gamma'$ konsistent sein.

Nun zur Maximalität von $\Gamma'$. Es gelte $A\notin\Gamma'$. Da
mindestens ein $i$ mit $A=A_i$ existiert, muss $A\notin\Gamma_{i+1}$
gelten, denn $\Gamma_{i+1}\subseteq\Gamma'$. Das geht aber nur, wenn
$\Gamma_{i+1}=\Gamma_i$, also wenn $\Gamma_i\cup\{A\}$ inkonsistent ist.
Wegen $\Gamma_i\cup\{A\}\subseteq\Gamma'\cup\{A\}$ ist dann
$\Gamma'\cup\{A\}$ erst recht inkonsistent.\,\qedsymbol
\end{Beweis}

\begin{Satz}\label{maxkon-abgeschlossen}
Eine maximalkonsistente Menge $\Gamma$ ist deduktiv abgeschlossen.
Das heißt, es gilt $\Cn(\Gamma)=\Gamma$, also
\[(\Gamma\vdash_{\mathrm C} A)\cond A\in\Gamma.\]
\end{Satz}
\begin{Beweis}
Es gelte $\Gamma\vdash_{\mathrm C} A$. Angenommen, $\Gamma\cup\{A\}$
wäre inkonsistent. Dann gölte $\Gamma\vdash_{\mathrm C}\lnot A$
via Negationseinführung, ergo $\Gamma\vdash_{\mathrm C}\bot$ via
Negationsbeseitigung. Also wäre $\Gamma$ inkonsistent, was aber der
Konsistenz von $\Gamma$ widerspricht. Ergo muss $\Gamma\cup\{A\}$
konsistent sein. Wäre $A\notin\Gamma$, wäre $\Gamma\cup\{A\}$
inkonsistent aufgrund der Maximalkonsistenz von $\Gamma$, was der
gerade gemachten Feststellung widerspricht.
Ergo gilt $A\in\Gamma$.\,\qedsymbol
\end{Beweis}

\noindent
Der nächste Satz zeigt auf, dass eine maximalkonsistente Formelmenge
die Rekursionsschritte in der Erklärung der Erfüllungsrelation
reflektiert.

\begin{Satz}\label{maxkon-refl-sat}
Eine maximalkonsistente Menge $\Gamma$ besitzt die Eigenschaften:
\[\begin{array}{@{}ll}
1. & \bot\notin\Gamma.\\
2. & (\lnot A)\in\Gamma\bicond A\notin\Gamma,\\
3. & (A\cond B)\in\Gamma\bicond A\notin\Gamma\lor B\in\Gamma,\\
4. & (A\land B)\in\Gamma\bicond A\in\Gamma\land B\in\Gamma,\\
5. & (A\lor B)\in\Gamma\bicond A\in\Gamma\lor B\in\Gamma,\\
6. & (A\bicond B)\in\Gamma\bicond
  (A\cond B)\in\Gamma\land (B\cond A)\in\Gamma.
\end{array}\]
\end{Satz}
\begin{Beweis}
Zu 1. Dies gilt, weil $\Gamma$ konsistent ist.

Zu 2. Es gelte $\lnot A\in\Gamma$, ergo $\Gamma_{\mathrm C} A$.
Angenommen, $A\in\Gamma$. Dann $\Gamma_{\mathrm C} A$, also
$\Gamma_{\mathrm C}\bot$ via Negationsbeseitigung, was aber der
Konsistenz von $\Gamma$ widerspricht. Ergo $A\notin\Gamma$.
Umgekehrt gelte $A\notin\Gamma$. Aufgrund der Maximalkonsistenz
von $\Gamma$ bedeutet dies $\Gamma\cup\{A\}\vdash_{\mathrm C}\bot$,
also $\Gamma\vdash_{\mathrm C}\lnot A$ via Negationseinführung,
ergo $\lnot A\in\Gamma$ vermittels Satz \ref{maxkon-abgeschlossen}.

Zu 3. Es gelte $A\cond B\in\Gamma$, ergo $\Gamma\vdash_{\mathrm C}A\cond B$.
Fallunterscheidung. Im Fall $A\notin\Gamma$ folgt die Disjunktion
unmittelbar. Im Fall $A\in\Gamma$ gilt $\Gamma\vdash_{\mathrm C} A$,
also $\Gamma\vdash_{\mathrm C}B$ via Implikationsbeseitigung.
Und somit $B\in\Gamma$ vermittels Satz \ref{maxkon-abgeschlossen}.
Ergo die Disjunktion. Umgekehrt gelte $A\notin\Gamma\lor B\in\Gamma$.
Fallunterscheidung. Im Fall $A\notin\Gamma$ gilt $(\lnot A)\in\Gamma$
gemäß 2. Ergo $\Gamma\vdash_{\mathrm C}\lnot A$. Zusammen mit
$A\vdash_{\mathrm C} A$ folgt $\Gamma,A\vdash_{\mathrm C}\bot$
via Negationsbeseitigung, also $\Gamma,A\vdash_{\mathrm C}B$ via
ex falso quodlibet. Im Fall $B\in\Gamma$ gilt $\Gamma\vdash_{\mathrm C}B$,
womit man via Abschwächung zu $\Gamma,A\vdash_{\mathrm C} B$ gelangt.
Es folgt $\Gamma\vdash_{\mathrm C}A\cond B$ via Implikationseinführung,
ergo $(A\cond B)\in\Gamma$ vermittels Satz \ref{maxkon-abgeschlossen}.

Zu 4. Es gelte $(A\land B)\in\Gamma$, ergo $\Gamma\vdash_{\mathrm C}A\land B$.
Via Konjunktionsbeseitigung folgt $\Gamma\vdash_{\mathrm C}A$ sowie
$\Gamma\vdash_{\mathrm C}B$. Vermittels Satz \ref{maxkon-abgeschlossen}
folgt daraus $A\in\Gamma$ sowie $B\in\Gamma$. Umgekehrt gelte
$A\in\Gamma$ und $B\in\Gamma$, ergo $\Gamma\vdash_{\mathrm C}A$
und $\Gamma\vdash_{\mathrm C}B$, also $\Gamma\vdash_{\mathrm C}A\land B$
via Konjunktionseinführung. Ergo $A\land B\in\Gamma$ vermittels Satz
\ref{maxkon-abgeschlossen}. Ergo folgt die Disjunktion. Umgekehrt
gelte diese Disjunktion $A\in\Gamma\lor B\in\Gamma$. Fallunterscheidung.
Im Fall $A\in\Gamma$ gilt $\Gamma\vdash_{\mathrm C} A$, also
$\Gamma\vdash_{\mathrm C} A\lor B$ via Disjunktionseinführung.
Im Fall $B\in\Gamma$ gelangt man analog zu $\Gamma\vdash_{\mathrm C} A\lor B$.
Ergo folgt $(A\lor B)\in\Gamma$ vermittels Satz \ref{maxkon-abgeschlossen}.

Zu 5. Es gelte $(A\lor B)\in\Gamma$, ergo
$\Gamma\vdash_{\mathrm C}A\lor B$. Fallunterscheidung. Im Fall
$A\in\Gamma$ folgt die Disjunktion unmittelbar. Im Fall $A\notin\Gamma$
ist $\Gamma\cup\{A\}$ inkonsistent, ergo $\Gamma\cup\{A\}\vdash_{\mathrm C} B$
via ex falso quodlibet. Weiterhin folgt $\Gamma\cup\{B\}\vdash_{\mathrm C} B$
via Abschwächung aus $B\vdash_{\mathrm C} B$. Via Disjunktionsbeseitigung
erhält man somit $\Gamma\vdash_{\mathrm C} B$, ergo $B\in\Gamma$
vermittels Satz \ref{maxkon-abgeschlossen}.

Der Beweis von 6. verläuft analog zu dem von 4.\,\qedsymbol
\end{Beweis}

\begin{Satz}\label{maxkon-sat}
Eine maximalkonsistente Formelmenge ist stets erfüllbar.
\end{Satz}
\begin{Beweis}
Sei $\Gamma$ maximalkonsistent. Die Interpretation $I$ sei definiert
als $I(P):=1$, wenn $P\in\Gamma$, sonst $I(P):=0$. Es gelte $F\in\Gamma$.
Zu zeigen ist $I\models F$. Induktion über den Aufbau von $F$. Ist $F$
atomar, gilt die gemäß der gerade gemachten Definition. In allen anderen
Fällen wird der jeweilige Bestandteil von Satz \ref{maxkon-refl-sat}
verwendet. Beispiel $F=(A\land B)$. Die Induktionsannahmen sind
Mit $A\land B\in\Gamma$ haben wir $A\in\Gamma$ sowie $B\in\Gamma$
gemäß Sat \ref{maxkon-refl-sat}. Daraus folgt $I\models A$ sowie
$I\models B$ vermittels den Induktionsannahmen. Also
$I\models A\land B$ laut Definition der Erfüllbarkeit. Die anderen
Fälle verlaufen analog.\,\qedsymbol
\end{Beweis}

Im Abschluss des Beweisen kommen nun die Hilfssätze der Reihe nach
zur Anwendung.

\begin{Satz}\label{kon-sat}
Eine konsistente Formelmenge ist stets erfüllbar.
\end{Satz}
\begin{Beweis}
Sei $\Gamma$ konsistent. Laut Satz \ref{Lemma-Lindenbaum} existiert ein
maximalkonsistentes $\Gamma'$ mit $\Gamma\subseteq\Gamma'$. Laut Satz
\ref{maxkon-sat} ist $\Gamma'$ erfüllbar, folglich $\Gamma$ als
Teilmenge erst recht erfüllbar.\,\qedsymbol
\end{Beweis}

\begin{Satz}[Vollständigkeit des natürlichen Schließens]\newlinefirst
Aus $\Gamma\models A$ folgt $\Gamma\vdash_{\mathrm C} A$.
\end{Satz}
\begin{Beweis}
Es gelte $\Gamma\models A$. Laut Satz \ref{gueltig-unerf} ist
$\Gamma\cup\{\lnot A\}$ also unerfüllbar, laut Kontraposition von Satz
\ref{kon-sat} somit inkonsistent. Laut Satz \ref{inkon-ableitbar} folgt
daher $\Gamma\vdash_{\mathrm C} A$.\,\qedsymbol
\end{Beweis}

\noindent
Gleichwohl wurden nur die ersten Schritte unternommen. Die Argumentation
wäre im Weiteren auf die Logik erster Stufe zu übertragen; das Resultat
ist als der \emph{gödelsche Vollständigkeitssatz} bekannt. Zudem besteht
der Wunsch, auch die Vollständigkeit der jeweiligen Kalküle zu anderen
Logiken -- wie bspw. der intuitionistischen Logik -- zu versichern.

\newpage
\section{Die klassische Semantik der Logik erster Stufe}

\subsection{Strukturen}

Wie in der Aussagenlogik soll den Formeln eine Bedeutung zukommen.
Eine Formel der Prädikatenlogik ist allerdings komplizierter als eine
Formel der Aussagenlogik. Statt nur atomare Variablen mit
Wahrheitswerten zu belegen, sind nun unterschiedliche Arten von
Symbolen vorhanden, die mit einer Bedeutung versehen werden müssen.

Es gibt erstens die logischen Verknüpfungen, denen eine für die
Logik spezifische feste Funktion zukommt. Ebenso ist die Gleichheit
zweier Terme fest definiert. Zweitens können in einer Formel Symbole
für Funktionen und Relationen vorkommen, darunter fallen auch
Konstanten. Sie werden je nach gewählter Struktur anders mit
Funktionen und Relationen belegt. Drittens darf eine Formel
freie Variablen enthalten. Sie werden durch eine extra definierte
Belegungsfunktion mit Werten belegt. Bei der Interpretation einer
quantifizierten Teilformel wird die Belegungsfunktion modifiziert.

Je nach Festlegung des logisch"=mathematischen Systems und dessen
inhaltlicher Deutung haben die mathematischen Terme Werte
in einer bestimmten Menge $U$, die wir \emph{Universum},
\emph{Domäne}, \emph{Träger}, \emph{Grundbereich} oder
\emph{Diskursuniversum} nennen.  Die Elemente von $U$ nennen wir
\emph{Individuen} oder \emph{Objekte}. Ich möchte noch einmal daran
erinnern, dass der beschriebene Kalkül der Prädikatenlogik die
Menge $U$ als nichtleer fordert.

Eine \emph{Struktur} $M$ sei das Universum $U$ zusammen mit einer
Interpretationsfunktion, die jedem Funktionssymbol $f$ von $n$ Argumenten
eine Funktion $f^M\colon U^n \to U$ und jedem Relationssymbol $R$ von
$n$ Argumenten eine Relation $R^M\colon U^n\to\{0,1\}$ zuordnet.
Konstanten deuten wir als Funktionen von null Argumenten.

Man kann eine Struktur auf unterschiedliche Art kodieren, wobei
unter Umständen die Hilfsbegriffe \emph{Signatur} und
\emph{Ariätsfunktion} gebraucht werden. Die Signatur beschreibt hierbei
die Menge der Funktionssymbole und die Menge Relationssymbole. Die
Aritätsfunktion ordnet jedem Funktionssymbol und jedem Relationssymbol
seine Stelligkeit zu, das ist die Anzahl der Argumente.

\strong{Beispiel 1. Ein Kalkül für Halbordnungen.} Zusätzlich zu
den Regeln und Axiomen der Prädikatenlogik sollen als mathematische Axiome
demnach noch die Axiome für Halbordnungen gelten. Die Signatur enthalte
dazu lediglich das Relationssymbol $\le$. Eine mögliche Struktur $M$
ist die Menge $U:=\mathcal P(G)$ zu irgendeiner Grundmenge $G$, zusammen
mit der Setzung ${\le^M} := {\subseteq}$, so dass je zwei Mengen $A,B\in U$
ein Wahrheitswert $A\subseteq B$ zugeordnet wird. Man notiert so
eine Struktur üblicherweise einfach schludrig als $(U,\le^M)$,
konkret $(\mathcal P(G),\subseteq)$.

\strong{Beispiel 2. Ein Kalkül für Gruppen.} Zusätzlich
zu den Regeln und Axiomen der Prädikatenlogik sollen als mathematische Axiome
demnach noch die Axiome für Gruppen gelten. Die Signatur enthalte
dazu ein Konstantensymbol für das neutrale Element, die Verknüpfung von
Elementen als zweistelliges Funktionssymbol und die Invertierungsfunktion
als einstelliges Funktionssymbol. Eine Struktur notiert man dann als
$(G,\cdot,e,\mathrm{inv})$. Eine mögliche Struktur ist $(\R,+,0,\mathrm{inv})$
mit $\mathrm{inv}(x):=-x$, die Gruppe der reellen Zahlen bezüglich
Addition. Eine andere mögliche Struktur ist $(S(X),\circ,\id,\mathrm{inv})$
mit $\mathrm{inv}(f):=f^{-1}$, die Gruppe der bijektiven Selbstabbildungen
auf der Menge $X$ bezüglich der Verkettung, auch symmetrische Gruppe
genannt. Es ist hier $\id$ die identische Abbildung und $f^{-1}$
die Umkehrabbildung von $f$.

\strong{Beispiel 3. Ein Kalkül für die Mengenlehre.} Der unerfahrene Leser mag
diesbezüglich zunächst das Kapitel über Mengenlehre studieren und später
hierher zurückkommen. Das Diskursuniversum soll das Mengenuniversum sein.
Nun verhält es sich allerdings so, dass die Zusammenfassung aller denkbaren
Mengen eine echte Klasse bildet, die keine Menge mehr sein kann. Man
könnte die Definition des Diskursuniversums so modifizieren, dass es eine
beliebige Klasse sein darf, was aber dazu führen würde, dass dieses
nicht mehr ohne Ausnahme als Element eines Mengensystems auftreten darf.
Wir gehen dieser Komplikation durch die Betrachtung des
Grothendieckuniversums $U=V_\kappa$ für ZFC oder $U=V_{\kappa+1}$ für
die Morse"=Kelly"=Mengenlehre aus dem Weg, wobei mit den $V_\alpha$ die
kumulative Hierarchie und mit $\kappa$ eine als Ordinalzahl betrachtete
unerreichbare Kardinalzahl gemeint ist. Die Einzelheiten erläutert
Shulman \cite{Shulman}. Einziges Relationssymbol ist nun  ${\in}$ für
die Elementrelation, und Funktionssymbole gibt es gar keine. Das mag
verwirrend erscheinen, wo die Mengenlehre doch von Relationen und
Funktionen durchdrungen ist. Diese stellen bei näherer Betrachtung aber
Objekte des Universums dar.

Aber ist die leere Menge nicht ein Konstantensymbol, die Inklusion
${\subseteq}$ nicht ein weiteres Relationssymbol und Operatoren wie
$\cap$ und $\cup$ Funktionssymbole? An sich stimmt das. Das
formale System der Mengenlehre kann um diese Symbole erweitert werden,
die minimalistische Fassung des Systems kommt jedoch allein mit
${\in}$ aus. Alle weiteren Symbole werden durch ihre Definition
ersetzt, obgleich dadurch unter Umständen exorbitant lange Formeln
entstehen. Zum Beispiel wird die Formel $A\cap B\subseteq A\cup B$
expandiert zu
\[\forall x\colon x\in A\land x\in B\cond x\in A\lor x\in B.\]

\subsection{Interpretationen}

Zur Interpretation eines Terms $t$ mit freien Variablen betrachtet
man außerdem noch eine \emph{Belegung} $\beta$, die jeder Variablen $x$
einen Wert $\beta(x)\in U$ zuordnet. Geläufig ist ebenfalls die
Bezeichnung \emph{Zuweisung}.

Man nennt nun $I=(M,\beta)$ eine \emph{Interpretation} der Formel $A$.
Wir definieren $I(t)$ für einen Term $t$ rekursiv über den Termaufbau
als%
\begin{gather*}
I(x) := \beta(x),\\
I(f(t_1,\ldots,t_n)) := f^M(I(t_1),\ldots,I(t_n)),
\end{gather*}
wobei $x$ eine Variable und $f$ ein $n$-stelliges Funktionssymbol
ist. Die Belegung mit einer modifiziert belegten Variable wird
diesbezüglich definiert als%
\begin{gather*}
\beta[y{:=}u](x) := \begin{cases}
u, & \text{wenn}\;x=y,\\
\beta(x), & \text{sonst}.
\end{cases}
\end{gather*}
Für $I=(M,\beta)$ sei $I[y{:=}u] := (M,\beta[y{:=}u])$.

Um Missverständnissen aus dem Weg zu gehen, sollte ich erwähnen, dass
die Begriffe \emph{Variable} und \emph{Atom} in der Prädikatenlogik
anders verstanden werden als in der Aussagenlogik. Die atomaren
Variablen der Aussagenlogik entsprechen nullstelligen Relationssymbolen,
denen bei einer Interpretation Relationen von null Argumenten zukommt,
das sind schlicht die konstanten Wahrheitswerte. Dagegen meint Variable
in der Prädikatenlogik üblicherweise eine Individuenvariable. Eine
atomare Formel ist in der Prädikatenlogik nicht nur ein nullstelliges
Relationssymbol, sondern jede Applikation eines Relationssymbols und
jede Gleichung.

Die Erfüllungsrelation definiert man analog zur Aussagenlogik, also
\[(I\models A\land B)\,:\bicond\,(I\models A)\land (I\models B)\]
und so weiter. Zusätzlich setzt man%
\begin{gather*}
(I\models t_1 = t_2) :\Leftrightarrow I(t_1) = I(t_2),\\
(I\models R(t_1,\ldots,t_n)):\Leftrightarrow R^M(I(t_1),\ldots,I(t_n)),\\
(I\models\forall x\colon A) :\Leftrightarrow\forall u\in U\colon I[x{:=}u]\models A,\\
(I\models\exists x\colon A) :\Leftrightarrow\exists u\in U\colon I[x{:=}u]\models A.
\end{gather*}
Man definiert hiermit die semantische Folgerung analog zur Aussagenlogik als
\[(\Gamma\models A)\,:\bicond\,\forall I\colon (I\models\Gamma)\cond (I\models A).\]
Wir nennen $A$ \emph{allgemeingültig}, falls $\emptyset\models A$, was
wir wieder durch $\models A$ abkürzen wollen.

Sofern die Formel $A$ geschlossen ist, das heißt, keine freien
Variablen besitzt, spielt die Belegung der Variablen noch keine Rolle, sie
entsteht eigentlich erst bei der rekursiven Bestimmung von $I\models A$
mit dem Einstieg in eine quantifizierte Formel. Bei einer einer
geschlossenen Formel $A$ darf man insofern $M\models A$ statt
$M,\beta\models A$ schreiben. Gilt $M\models A$, nennt man $M$ ein
\emph{Modell für} $A$.

Ich will erklären, warum die Interpretation der Quantoren so durchgeführt
wurde wie beschrieben. Intuitiv erfüllt $I$ die Formel
$\forall x\colon A$ genau dann, wenn $I$ die Formel $A[x:=u]$ für jedes
$u\in U$ erfüllt. Allerdings ist $u$ kein Term der logischen Sprache,
womit $A[x:=u]$ keine Formel mehr wäre. Wir könnten nicht einmal für
jedes $u$ eine Konstante zur Sprache hinzufügen, weil $U$ überabzählbar
sein darf, während eine gewöhnliche Sprache nur abzählbar viele
Symbole in Form von Zeichenketten umfassen kann. Die gewählte Ausweg
aus der Problematik besteht darin, dass die Belegung $x:=u$ erst einmal
in der Interpretation gespeichert wird, statt sie direkt
als Substitution auf die Formel anzuwenden.

\subsection{Korrektheit des natürlichen Schließens}

Wir müssen nun wieder aufzeigen, dass der Kalkül des natürlichen
Schließens bezüglich der definierten Semantik korrekt ist. Dies
geschieht wieder per struktureller Induktion über den Aufbau des
Beweises. Für die bereits in der Aussagenlogik vorhandenen Regeln
können wir die dort gemachte Argumentation unverändert übernehmen.
Zu bestätigen verbleibt die Korrektheit der Regeln zur Einführung und
Beseitigung der Quantoren, und ggf. noch die Korrektheit der Regeln zum
Umgang mit Gleichheiten.

\begin{Satz}\label{interpretation-lemma}
Es ist $I[x:=I(t)]\models A$ äquivalent zu $I\models A[x:=t]$.
\end{Satz}
\begin{Beweis}
Man unternimmt hierzu eine strukturelle Induktion über den Formelaufbau
von $A$. Die Induktionsvoraussetzungen sind jeweils unschwer zu verarbeiten.
Im Wesentlichen verbleibt nun%
\[I[x:=I(t)](v) = I(v[x:=t])\]
für jede Individuenvariable $v$ zu zeigen. Man nimmt dazu die
Fallunterscheidung zwischen $x=v$ und $x\ne v$ vor. Im ersten Fall gilt%
\[I[x:=I(t)](v) = I(t) = I(x[x:=t]) = I(v[x:=t]),\]
und im zweiten Fall gilt
\[I[x:=I(t)](v) = I(v) = I(v[x:=t]).\,\qedsymbol\]
\end{Beweis}

\begin{Satz}
Aus $\Gamma\models\forall x\colon A$ folgt $\Gamma\models A[x:=t]$.
\end{Satz}
\begin{Beweis}
Laut der Prämisse muss zu jedem $I$ mit $I\models\Gamma$
auch $I[x:=u]\models A$ für jedes $u\in U$ gelten. Weitere Prämisse
ist ein $I$ mit $I\models\Gamma$. Die Allaussage wird mit diesem $I$
und mit $u:=I(t)$ spezialisiert. Zu zeigen verbleibt%
\[(I[x:=I(t)]\models A)\cond (I\models A[x:=t]).\]
Diese gilt gemäß Satz \ref{interpretation-lemma}.\,\qedsymbol
\end{Beweis}

\begin{Satz}
Aus $\Gamma\models A[x:=t]$ folgt $\Gamma\models\exists x\colon A$.
\end{Satz}
\begin{Beweis}
Laut der Prämisse muss zu jedem $I$ mit $I\models\Gamma$ auch
$I\models A[x:=t]$ gelten. Weitere Prämisse ist ein $I$ mit
$I\models\Gamma$. Die Allaussage wird mit diesem $I$ spezialisiert.
Zu zeigen verbleibt%
\[(I\models A[x:=t])\cond \exists u\in U\colon (I[x:=u]\models A).\]
Gewählt wird $u:=I(t)$. Nun greift wieder Satz
\ref{interpretation-lemma}.\,\qedsymbol
\end{Beweis}

\begin{Satz}
Aus $\Gamma\models A$ und $x\notin\mathrm{FV}(\Gamma)$
folgt $\Gamma\models\forall x\colon A$.
\end{Satz}
\begin{Beweis}
Laut der Prämisse muss zu jedem $I$ mit $I\models\Gamma$ auch
$I\models A$ gelten. Weitere Prämisse ist ein $I$ mit
$I\models\Gamma$. Wegen $x\notin\mathrm{FV}(\Gamma)$ ist
$I\models\Gamma$ äquivalent zu $I[x:=u]\models\Gamma$. Demnach
lässt sich die Allaussage mit $I[x:=u]$ spezialisieren, womit
man $I[x:=u]\models A$ erhält. Ergo gilt%
\[\forall u\in U\colon (I[x:=u]\models A),\;\text{also}\;
I\models\forall x\colon A.\,\qedsymbol\]
\end{Beweis}

\begin{Satz}
Aus $\Gamma\models\exists x\colon A$ und $\Gamma\cup\{A\}\models B$
mit $x\notin\mathrm{FV}(\Gamma\cup\{B\})$ folgt $\Gamma\models B$.
\end{Satz}
\begin{Beweis}
Laut der ersten Prämisse existiert zu jedem $I$ mit $I\models\Gamma$
ein $u\in U$ mit $I[x:=u]\models A$. Laut der zweiten Prämisse muss
zu jedem $I$ mit $I\models\Gamma$ und $I\models A$ auch $I\models B$
gelten. Wegen $x\notin\mathrm{FV}(\Gamma\cup\{B\})$ ist $I\models\Gamma$
zu $I[x:=u]\models\Gamma$ sowie $I\models B$ zu $I[x:=u]\models B$
äquivalent. Dritte Prämisse ist ein $I$ mit $I\models\Gamma$. Die erste
Allaussage wird mit $I$ spezialisiert, die zweite mit $I[x:=u]$. 
Mit den vorliegenden Aussagen gelingt der Schluss auf $I[x:=u]\models B$.
Ergo gilt $I\models B$, wie gewünscht.\,\qedsymbol
\end{Beweis}

\section{Semantik der Modallogik}

\subsection{Die relationale Semantik der Modallogik}

Mit der Entwicklung der Modallogiken ging die Frage einher, wie
modallogische Formeln semantisch interpretiert werden können.
Besonders wichtig sind die nach Saul Kripke benannten
\emph{Kripke"=Semantiken}, auch \emph{relationale Semantiken} genannt.
Die Interpretation der atomaren Aussagen wird hierbei durch eine Welt
parametrisiert. Eine Aussage kann also in einer bestimmten Welt wahr
sein, während sie in einer anderen falsch ist. Als Logik der Metasprache
soll dabei die klassische Prädikatenlogik dienen.

Eine \emph{Kripke"=Struktur} sei ein Tripel $M=(W,R,V)$. Hierbei
ist $W$ eine Menge von Welten, $R$ eine zweistellige Relation auf
$W$, und $V$ eine Wertung, die jeder atomaren Variable $a$
parametrisiert durch die Welt $w\in W$ einen Wahrheitswert
$V(w,a)\in\{0,1\}$ zuordnet. Man nennt $R$ die
\emph{Zugänglichkeitsrelation}, wobei $R(w,w')$ gedeutet wird als
»$w'$ ist von $w$ aus zugänglich«. Die Relation muss nicht unbedingt
symmetrisch sein, womit man unter Umständen nicht mehr von
$w'$ aus in die ursprüngliche Welt $w$ zurückgelangt.

Das Paar $I=(M,w)$, bestehend aus einer Struktur und einer Welt, betrachten
wir als Interpretation. Man legt die Erfüllung einer atomaren Variable
$a$ bezüglich einer Interpretation fest als%
\[(M,w\models a) :\bicond V(w,a).\]
Die Erfüllung der Junktoren wird analog zur klassischen Semantik
definiert, also%
\[(I\models A\land B) :\bicond (I\models A)\land (I\models B)\]
und so weiter, wobei die rechte Seite gemäß den herkömmlichen
Wahrheitstafeln ausgewertet wird. Für die beiden modalen Junktoren legt
man nun fest%
\begin{align*}
(M,w\models\lnec A) &:\bicond\forall w'\in W\colon R(w,w')\cond (M,w'\models A),\\
(M,w\models\lpos A) &:\bicond\exists w'\in W\colon R(w,w')\land (M,w'\models A). 
\end{align*}
In Worten ist die Aussage $\lnec A$ genau dann wahr in der Welt $w$, wenn
$A$ in jeder von $w$ aus zugänglichen möglichen Welt wahr ist.
Es ist $\lpos A$ genau dann wahr in der Welt $w$, wenn $A$ in mindestens einer
von $w$ aus zugänglichen möglichen Welt wahr ist.

Die semantische Folgerung wird wieder in der herkömmlichen Form%
\[(\Gamma\models A)\,:\bicond\,\forall I\colon (I\models\Gamma)\cond (I\models A)\]
definiert. Ausgeschrieben lautet sie demnach
\[(\Gamma\models A)\,:\bicond\,\forall M\colon\forall w\in W(M)\colon
(M,w\models\Gamma)\cond (M,w\models A).\]
Man spricht gelegentlich auch von der \emph{lokalen} semantischen Folgerung,
um sie von einer weiteren, der selten diskutierten \emph{globalen} semantischen
Folgerung abgrenzen zu können. Zur Unterscheidung will ich die globale
Folgerung als $\models^g$ notieren. Man definiert sie als%
\[(\Gamma\models^g A)\,:\bicond\,\forall M\colon (M\models\Gamma)\cond (M\models A),\]
wobei $M\models A$ für $\forall w\in W(M)\colon (M,w\models A)$ stehen soll.

Ist $\Gamma$ die leere Menge, koinzidieren die lokale und die globale
Folgerung. Das heißt, $\emptyset\models A$ gilt genau dann, wenn
$\emptyset\models^g A$. Man schreibt wieder $\models A$ als Kurzform
von $\emptyset\models A$ und sagt dazu, die Formel $A$ sei
\emph{allgemeingültig}.

Bei der tieferen Untersuchung, welche Beziehungen zwischen Kalkül und
Semantik bestehen, zerlegt man eine Kripke"=Struktur $M=(W,R,V)$ in
die zwei Teile $F=(W,R)$ und $V$. Man nennt $F$ den \emph{Rahmen}.

Gezeigt werden muss nun die Korrektheit des natürlichen Schließens
für die Modallogik. Dies geschieht wieder per struktureller Induktion
über den Formelaufbau. Bei den herkömmlichen Schlussregeln verläuft der
Beweis identisch wie in der klassischen Semantik der Aussagenlogik.
Zu bestätigen verbleibt insofern lediglich noch die Korrektheit der
Nezessisierungsregel und die Allgemeingültigkeit des Axiomenschemas K.

\begin{Satz}
Aus $\models A$ folgt $\models\lnec A$.
\end{Satz}
\begin{Beweis}
Gemäß der Prämisse gilt $M,w'\models A$ für jedes $w'$. Demnach gilt
erst recht $R(w,w')\cond (M,w'\models A)$. Ergo gilt
$\models\lnec A$. Quod erat demonstrandum.\,\qedsymbol
\end{Beweis}

\begin{Satz}
Das Schema $\lnec(A\cond B)\cond (\lnec A\cond\lnec B)$ ist allgemeingültig.
\end{Satz}
\begin{Beweis}
Seien $M,w$ fest, aber beliebig. Zu zeigen ist
\[M,w\models\lnec(A\cond B)\cond (\lnec A\cond\lnec B),\]
was gemäß der Definition der Erfüllung äquivalent umgeformt wird zu%
\[(M,w\models\lnec(A\cond B))\cond (M,w\models\lnec A)\cond (M,w\models\lnec B).\]
Zu zeigen ist daher $M,w'\models B$ unter Annahme von $R(w,w')$ und%
\begin{gather*}
\forall w'\colon R(w,w')\cond (M,w'\models A\cond B),\\
\forall w'\colon R(w,w')\cond (M,w'\models A).
\end{gather*}
Die beiden Allaussagen werden mit $w'$ und $R(w,w')$ spezialisiert.
Laut der ersten Aussage $M,w'\models A\cond B$ folgt
$M,w'\models B$ aus $M,w'\models A$. Laut zweiten gilt $M,w'\models A$.
Ergo gilt $M,w'\models B$. Quod erat demonstrandum.\,\qedsymbol
\end{Beweis}

\subsection{Die Standardübersetzung}

Die relationale Semantik der diskutierten Modallogiken deutet bereits
an, dass Modallogik etwas mit der Logik erster Stufe zu tun haben müsste.
Es scheint, dass in gewisser Weise $\lnec A$ der Allquantifizierung,
und $\lpos A$ der Existenzquantifizierung entspricht. Die im Folgenden
beschriebene \emph{Standardübersetzung} modallogischer Formeln in die
Logik erster Stufe präzisiert diesen Gedankengang. Mithin stellen sich
die diskutierten Modallogiken als Spezialfälle der Logik erster Stufe
heraus. Auf diese Weise erhält man eine Möglichkeit, modallogische
Theoreme mit Beweisassistenten zu formulieren, die lediglich die
herkömmliche Logik verstehen.

Die Standardübersetzung ist rekursiv definiert als
\begin{align*}
&\mathrm{ST}_x(a) := P_a(x),
&& \mathrm{ST}_x(A\cond B) := \mathrm{ST}_x(A)\cond\mathrm{ST}_x(B),\\
&\mathrm{ST}_x(\lnot A) := \lnot\mathrm{ST}_x(A),
&& \mathrm{ST}_x(A\bicond B) := \mathrm{ST}_x(A)\bicond\mathrm{ST}_x(B),\\
&\mathrm{ST}_x(A\land B) := \mathrm{ST}_x(A)\land\mathrm{ST}_x(B),
&& \mathrm{ST}_x(\lnec A) := \forall y\colon R(x,y)\cond\mathrm{ST}_y(A),\\
& \mathrm{ST}_x(A\lor B) := \mathrm{ST}_x(A)\lor\mathrm{ST}_x(B),
&& \mathrm{ST}_x(\lpos A) := \forall y\colon R(x,y)\land\mathrm{ST}_y(A).
\end{align*}
Die Formeln $\bot,\top$ bleiben fix, also
$\mathrm{ST}_x(\bot) := \bot$ und $\mathrm{ST}_x(\top) := \top$.

Hierbei ist $R$ ein zweistelliges Prädikatsymbol, und $P_a$ zu
jedem $a$ ein einstelliges Prädikatsymbol. Eine Interpretation
$(M,w)$ mit Kripke"=Struktur $M=(W,R,V)$ und Welt $w\in W$
lässt sich als Interpretation $(M,\beta[x:=w])$ für die übersetzten
Formeln deuten, dergestalt dass das Symbol $P_a$ mit $w\mapsto V(w,a)$
interpretiert wird, und das Symbol $R$ mit der Zugänglichkeitsrelation
$R$. Demzufolge nimmt $W$ die Rolle des Diskursuniversums ein.
Die Belegung $\beta$ spielt keine Rolle, da $x$ die einzige freie
Variable ist. Man gelangt nun zum folgenden Resultat.
\begin{Satz}
Es gilt $M,w\models A$ genau dann, wenn $M,\beta[x:=w]\models\mathrm{ST}_x(A)$.
\end{Satz}
\begin{Beweis}
Per struktureller Induktion über den Formelaufbau von $A$.
Es sei $\beta\tfrac{w}{x}$ Kurzschreibweise für $\beta[x:=w]$. Den
Induktionsanfang liefert die Umformung
\[(M,w\models a) \iff V(w,a) \iff (M,\beta\tfrac{w}{x}\models P_a(x)).\]
Zur Formel $A\land B$ liegt die Äquivalenz von $M,w\models C$
und $M,\beta\tfrac{w}{x}\models\mathrm{ST}_x(C)$ jeweils für $C=A$
und $C=B$ als Induktionsvoraussetzung vor. Es findet sich hiermit
die äquivalente Umformung
\begin{align*}
(M,w\models A\land B) &\iff (M,w\models A)\land (M,w\models B)\\
&\iff (M,\beta\tfrac{w}{x}\models\mathrm{ST}_x(A))\land
(M,\beta\tfrac{w}{x}\models\mathrm{ST}_x(B))\\
&\iff M,\beta\tfrac{w}{x}\models\mathrm{ST}_x(A)\land\mathrm{ST}_x(B)\\
&\iff M,\beta\tfrac{w}{x}\models\mathrm{ST}_x(A\land B).
\end{align*}
Bei den restlichen herkömmlichen Junktoren verläuft die Argumentation
analog. Zur Formel $\lnec A$ liegt die Äquivalenz von $M,w'\models A$
und $M,\beta\tfrac{w'}{y}\models\mathrm{ST}_y(A)$ als
Induktionsvoraussetzung vor. Es findet sich hiermit die äquivalente
Umformung
\begin{align*}
(M,w\models\lnec A) &\iff \forall w'\colon R(w,w')\cond (M,w'\models A)\\
&\iff\forall w'\colon R(w,w')\cond (M,\beta\tfrac{w'}{y}\models\mathrm{ST}_y(A))\\
&\iff\forall w'\colon (M,\beta\tfrac{w}{x}\tfrac{w'}{y}\models R(x,y))
\cond (M,\beta\tfrac{w}{x}\tfrac{w'}{y}\models\mathrm{ST}_y(A))\\
&\iff (M,\beta\tfrac{w}{x}\models\forall y\colon R(x,y)\cond\mathrm{ST}_y(A))\\
&\iff (M,\beta\tfrac{w}{x}\models\mathrm{ST}_x(\lnec A)).
\end{align*}
Bei $\lpos A$ verläuft die Argumentation analog.\,\qedsymbol
\end{Beweis}

\noindent
Mithin gilt in unmittelbarer Konsequenz der
\begin{Satz}
Es gilt $\Gamma\models A$ genau dann, wenn
$\mathrm{ST}_x(\Gamma)\models\mathrm{ST}_x(A)$.
\end{Satz}
Vermittels dieser Äquivalenz erhält man sogleich den
\begin{Satz}
Aus $\vdash\mathrm{ST}_x(A)$ folgt $\models A$.
\end{Satz}
Kraft der Standardübersetzung in die Logik erster Stufe ergibt sich
insofern wie gewünscht ein alternativer korrekter Kalkül für die
diskutierten Modallogiken. Eine tiefergehende Untersuchung zur
Standardübersetzung findet man in \cite{Blackburn}.

\section{Semantik der intuitionistischen Logik}

\subsection{Die relationale Semantik der intuitionistischen Logik}

Wir erinnern uns, dass sich die intuitionistische Logik ergibt, indem
beim natürlichen Schließen auf die Beseitigung der Doppelnegation
verzichtet wird. Dies verdeutlicht, wie ein Beweis eines intuitionistischen
Theorems beschaffen sein muss. Schwierig bleibt jedoch der Nachweis, dass
eine Formel kein solches Theorem sein kann. Hier leistet die Semantik
Aufschluss.

Mit der Zeit fand man recht viele verschiedene Semantiken der
intuitionistischen Logik.\cite{Artemov} Im Rahmen der algebraischen
Semantik wird sie durch Heyting"=Algebren statt boolescher Algebren
beschrieben. Weiterhin gibt es die topologische Semantik, die
Realisierer nach Kleene, die Beth"=Semantik, die Kripke"=Semantik
und die kategorielle Semantik. Streng genommen sollten wir all diese
Semantiken aufführen, und daraufhin mit einem Ringschluss aufzeigen,
dass sie in Bezug auf die Allgemeingültigkeit einer Formel alle äquivalent
sind. Für eine solche Unternehmung bleibt hier aber leider nicht
genügend Zeit und Raum. Ich werde mich insofern auf Grundzüge beschränken.

Zur intuitionistischen Logik gehört eine Kripke"=Semantik,
deren Erfüllung auf eine eigensinnige Art beschaffen ist.\cite{Kripke-IL}
Eine \emph{intuitionistische Kripke"=Struktur} ist ein Tripel $(W,\le,V)$,
wobei $W$ wieder eine Menge von Welten, $\le$ eine Quasiordnung auf ihr und
$V$ eine monotone Wertung ist, das heißt,
\[\forall w\in W\colon\forall w'\in W\colon w\le w'\land V(w,P)\Rightarrow V(w',P),\]
wobei $V$ jeder Welt $w$ und jeder atomaren Variable $P$ einen
Wahrheitswert $V(w,P)\in\{0,1\}$ zuordnet.

Das Paar $I=(M,w)$, bestehen aus einer Struktur und einer Welt, betrachten
wir wieder als Interpretation. Die Erfüllung einer atomaren Variable $P$
ist wieder erklärt als
\[(M,w\models P)\;:\Leftrightarrow\; V(w,P).\]
Die Erfüllung der Junktoren ist daraufhin rekursiv erklärt als
\begin{align*}
(M,w\models A\land B) &\;:\Leftrightarrow\; (M,w\models A)\land (M,w\models B),\\
(M,w\models A\lor B) &\;:\Leftrightarrow\; (M,w\models A)\lor (M,w\models B),\\
(M,w\models A\cond B) &\;:\Leftrightarrow\;
  \forall w'{\in}W\colon w\le w'\land (M,w'\models A)\cond (M,w'\models B).
\end{align*}
Der jeweilige Junktor auf der rechten Seite ist hierbei abermals
als die boolesche Funktion bezüglich ihrer gewöhnlichen Wahrheitstafel
eklärt. Weiterhin setzt man $(M,w\models\bot) :\Leftrightarrow 0$, wie
gewöhnlich. Die Negation $\lnot A$ erhält die Semantik von
$A\cond\bot$. Die Bijunktion $A\bicond B$ erhält die Semantik
von $(A\cond B)\land (B\cond A)$.

Die Gültigkeit von $A$ unter $\Gamma$ ist wieder erklärt als
\[(\Gamma\models A)\;:\Leftrightarrow\; \forall M\colon\forall w\in W(M)\colon
(M,w\models \Gamma)\Rightarrow (M,w\models A).\]
Der Beweis Korrektheit verläuft fast analog wie bei der klassischen
Semantik. Wir müssen lediglich bei der Subjunktion bzw. Negation aufpassen,
die eine veränderte Erfüllung besitzt. Zu bestätigen wäre also die Gültigkeit der
Einführung und Beseitigung der Subjunktion. Zunächst zur Beseitigung
\[(\Gamma\models A\cond B)\land (\Gamma\models A)\Rightarrow (\Gamma\models B).\]
Aus $M,w\models\Gamma$ ist $M,w\models B$ abzuleiten. Vermittels
$\Gamma\models A$ erhält man zunächst $M,w\models A$. Nun entfaltet sich
$\Gamma\models A\cond B$ zu
\[(M,w\models\Gamma)\Rightarrow\forall w'\colon w\le w'\land
(M,w'\models A)\Rightarrow (M,w'\models B).\]
Wir spezialisieren $w':=w$, denn $w\le w$ gilt aufgrund der Reflexivität
von $\le$.

Nun zur Einführung
\[(\Gamma, A\models B)\Rightarrow (\Gamma\models A\cond B).\]
Aus $M,w\models\Gamma$ ist $M,w\models A\cond B$ abzleiten, also
$M,w'\models B$ für jedes $w'$ mit $w\le w'$ und $M,w'\models A$.
Via Spezialisierung von $\Gamma,A\models B$ mit $w'$ erhält man
$M,w'\models B$ aus $M,w'\models\Gamma$ und $M,w'\models A$. Zu zeigen
verbleibt also $M,w'\models\Gamma$. Um den Beweis abzuschließen,
müssen wir $M,w'\models\Gamma$ aus $w\le w'$ und $M,w\models\Gamma$
ableiten. Hierfür ist allerdings ein Hilfssatz erforderlich, laut dem
sich die Monotonie von der Wertung auf die Erfüllung vererbt.

\begin{Satz}
Aus $w\le w'$ und $M,w\models A$ folgt $M,w'\models A$ zu jeder Formel
$A$, jeder intuitionistischen Struktur $M=(W,\le,V)$ und allen Welten $w,w'\in W$.
\end{Satz}
\begin{Beweis}[Beweisskizze]
Strukturelle Induktion über den Aufbau von $A$.

Im Induktionsanfang in der atomaren Variable $P$ gilt dies aufgrund
der Monotonie von $V$. Der Induktionsanfang in $\bot$ ist trivial.

Der jeweilige Induktionsschritt zu $A\land B$, $A\lor B$, $A\cond B$
verläuft unkompliziert unter Anwendung der beiden Induktionsvoraussetzungen,
dass $M,w'\models A$ aus $M,w\models A$ sowie $M,w'\models B$ aus
$M,w\models B$ für $w\le w'$ folgt. Bei $A\cond B$ ist
allerdings zusätzlich die Transitivität von $\le$ erforderlich.\,\qedsymbol
\end{Beweis}

\noindent
Für die Gültigkeit macht es keinen Unterschied, wenn die Erfüllung
lediglich für Halb- statt Quasiordnungen gefordert wird. Dies erklärt,
warum die entsprechend modifizierte Definition der Semantik ebenfalls
verbreitet ist. Wir wollen dies noch kurz verifizieren. Den Hauptteil
übernimmt der folgende Hilfssatz.

\begin{Satz}\label{Modell-quasi-halb}
Zu jedem intuitionistischen Kripke"=Modell $M=(W,R,V)$ existiert ein
intuitionistisches Kripke"=Modell $M' = (W',R',V')$ mit Halbordnung
$R'$, das genau dieselben Formeln erfüllt.
\end{Satz}
\begin{Beweis}
Mit $x\sim y:\Leftrightarrow R(x,y)\land R(y,x)$ ist eine
Äquivalenzrelation festgelegt, wie man unschwer nachrechnet.
Die Quotientenmenge $W' := W/{\sim}$ wird das \emph{Skelett} von $W$
genannt, ihre Elemente heißen \emph{Cluster}. Auf $W'$ definieren wir
\[R'([x],[y]) :\Leftrightarrow R(x,y),\quad\; V'([x],P):\Leftrightarrow V(x,P).\]
Man rechnet unschwer nach, dass $R',V'$ wohldefiniert sind, $R'$ eine
Halbordnung ist, und $V'$ monoton ist. Für jede Formel $A$ gilt nun
\[\forall x\in W\colon (M,x\models A)\bicond (M',[x]\models A),\]
was per struktureller Induktion über den Aufbau von $A$ bestätigt wird.

Zum Induktionsanfang in der atomaren Variable $P$ findet sich die äquivalente Umformung
\[(x\models P)\iff V(x,P)\iff V([x],P)\iff ([x]\models P).\]
Der Induktionsanfang in $\bot$ ist trivial.

Im Induktionsschritt zu $A\land B$ liegen die Induktionsvoraussetzungen
$\forall x\colon (x\models A)\Leftrightarrow ([x]\models A)$ und
$\forall x\colon (x\models B)\Leftrightarrow ([x]\models B)$ vor.
Mit diesen rechnet man
\[(x\models A\land B)\iff (x\models A)\land (x\models B)
\iff ([x]\models A)\land ([x]\models B)
\iff ([x]\models A\land B).\]
Der Schritt zu $A\lor B$ verläuft analog. Im Schritt zu $A\cond B$ rechnet man
\begin{align*}
(x\models A\cond B)&\iff \big(\forall y\in W\colon R(x,y)\land (y\models A)\cond (y\models B)\big)\\
&\iff \big(\forall y\in W\colon R'([x],[y])\land ([y]\models A)\cond ([y]\models B)\big)\\
&\iff \big(\forall y'\in W'\colon R'([x],y')\land (y'\models A)\cond (y'\models B)\big)\\
&\iff ([x]\models A\cond B),
\end{align*}
insofern dank der Surjektivität der kanonischen Projektion $y\mapsto [y]$
zu jedem $y'$ ein $y$ mit $y'=[y]$ existiert.\,\qedsymbol
\end{Beweis}

\noindent
Mit der geleisteten Vorarbeit ergibt sich der

\begin{Satz}\label{Gueltigkeit-halb-quasi}
Aus $\Gamma\models' A$ folgt $\Gamma\models A$, wobei Ersteres die
Gültigkeit in Bezug auf lediglich alle Halb- statt Quasiordnungen meine.
\end{Satz}
\begin{Beweis}
Aus $M,w\models\Gamma$ ist $M,w\models A$ zu folgern. Zu $M$ erhält man
$M'=(W',R',V')$ mit Halbordnung $R'$ vermittels Satz
\ref{Modell-quasi-halb}. Laut der Äquivalenz von Satz \ref{Modell-quasi-halb}
folgt $M',[w]\models\Gamma$ aus $M,w\models\Gamma$. Daraufhin folgt
$M',[w]\models A$ laut Voraussetzung $\Gamma\models' A$. Unter nochmaliger
Anwendung der Äquivalenz folgt schließlich $M,w\models A$.\,\qedsymbol
\end{Beweis}

\noindent
Die Umkehrung von Satz \ref{Gueltigkeit-halb-quasi} stellt natürlich
eine Trivialität dar, weil jede Halb- auch immer eine Quasiordnung ist.

Die Übersetzung nach Gödel-McKinsey-Tarski bettet eine Formel
der intuitionistischen Logik in die Modallogik S4 ein. Sie ist
rekursiv erklärt als
\begin{align*}
T(P) &:= \lnec P, & T(A\land B) &:= T(A)\land T(B), & T(A\cond B) &:= \lnec T(A\cond B),\\
T(\bot) &:= \bot, & T(A\lor B) &:= T(A)\lor T(B), & T(\lnot A) &:= \lnec\lnot T(A).
\end{align*}
Es soll $A$ genau dann ein Theorem der intuitionistischen Logik sein,
wenn $T(A)$ ein Theorem der Modallogik S4 ist. Im Folgen wollen wir
dies als Übung semantisch bestätigen.

Bei Kripke"=Semantiken ist für die Erfüllung auch die Schreibweise
$w\Vdash A$ statt $w\models A$ geläufig. Zur Unterscheidung notieren wir
-- von dieser Freiheit pragmatisch Gebrauch machend -- im Folgenden
$w\Vdash A$ für die Erfüllung in der intuitionistischen Logik und
$w\models T(A)$ für die Erfüllung in der S4, jeweils im Bezug auf die
jeweilige Kripke"=Semantik.

Der Beweis umfasst die Nutzung zweier Hilfssätze.

\begin{Satz}\label{Erf-Modell-S4-aus-Int}
Es liege ein inutitionistisches Kripke"=Modell $(W,\le,V)$ vor.
Dieses bietet sich unverändert als S4"=Kripke"=Modell an. Zu jeder Formel $A$
und jeder Welt $w$ ist dann $w\Vdash A$ zu $w\models T(A)$ äquivalent.
\end{Satz}
\begin{Beweis}[Beweisskizze]
Strukturelle Induktion über den Aufbau von $A$.

Zum Induktionsanfang in der atomaren Variable $P$. Es gelte $w\Vdash P$,
also $V(w,P)$. Zu zeigen ist $w\models\lnec P$, also $V(w',P)$
für jedes $w'$ mit $w\le w'$, was aber mit $V(w,P)$ aus der Monotonie
von $V$ folgt. Umgekehrt gelte $w\models\lnec P$. Zu zeigen ist
$V(w,P)$, was sofort aus $w\models\lnec P$ mit $w\le w$ folgt,
denn $\le$ ist reflexiv.

Der Induktionsanfang in $\bot$ ist trivial.

Die Induktionsschritte zu $A\land B$, $A\lor B$ und $A\cond B$ verlaufen
geradlinig unter Anwendung der beiden Induktionsannahmen, dass $w\Vdash A$
zu $w\models T(A)$ sowie $w\Vdash B$ zu $w\models T(B)$ äquivalent ist.
Der Umstand $T(A\cond B)=\lnec(A\cond B)$ schafft dabei
erstaunlicherweise keine sonderliche Komplikation, die den Gebrauch einer
der anderen Voraussetzungen erfordern würde.\,\qedsymbol
\end{Beweis}

\begin{Satz}\label{Erf-Modell-Int-aus-S4}
Es liege ein S4"=Kripke"=Modell $(W,\le,V)$ vor. Zu diesem ist
$(W,\le,V_i)$ mit $V_i(w,P):\Leftrightarrow (w\models\lnec P)$ ein
intuitionistisches Kripke"=Modell, so dass zu jeder Formel $A$ und jeder
Welt $w$ die Erfüllung $w\Vdash A$ zu $w\models T(A)$ äquivalent ist.
\end{Satz}
\begin{Beweis}[Beweisskizze]
Zunächst ist die Monotonie von $V_i$ zu bestätigen. Es seien $x,y,P$
beliebig, und es gelte $x\le y$ und $V_i(x,P)$. Zu zeigen ist
$V_i(y,P)$, also $V(z,P)$ für jedes $z$ mit $y\le z$. Zunächst
folgt $x\le z$, da $\le$ transitiv ist. Entfaltung von $V_i(x,P)$
bringt $x\models\lnec P$, also $V(w,P)$ für jedes $w$ mit $x\le w$.
Ergo erhält man $V(z,P)$ mit $w:=z$.

Nun wieder strukturelle Induktion über den Aufbau von $A$.

Der Beweis des Induktionsanfangs in der atomaren Variable $P$ verläuft
unkompliziert.

Der Beweis des Induktionsanfangs in $\bot$ ist trivial.

Die Induktionsschritte zu $A\land B$, $A\lor B$ und $A\cond B$ verlaufen
abermals geradlinig, jeweils unter Anwendung der beiden
Induktionsannahmen.\,\qedsymbol
\end{Beweis}

\begin{Satz}
Es gilt genau dann $\models_\text{Int} A$, wenn $\models_\text{S4} T(A)$.
Das heißt, eine Formel $A$ ist genau dann bezüglich intuitionistischer
Kripke"=Semantik gültig, wenn $T(A)$ bezüglich der Kripke"=Semantik von
S4 gültig ist.
\end{Satz}
\begin{Beweis}
Es gelte $\models_\text{Int} A$. Es sei $M=(W,\le,V)$ ein beliebiges
S4"=Kripke"=Modell. Es sei $w$ eine beliebige Welt. Nach Satz
\ref{Erf-Modell-Int-aus-S4}
hat man ein intuitionistisches Modell $M_i:=(W,\le,V_i)$, so dass
$M_i,w\Vdash A$ laut Prämisse, und somit $M,w\models T(A)$. Und dies
bedeutet $\models_\text{S4} T(A)$, da $M,w$ beliebig sind.

Umgekehrt gelte $\models_\text{S4} T(A)$. Es sei $M_i=(W,\le,V)$ ein
beliebiges intuitionistisches Kripke"=Modell. Nach Satz
\ref{Erf-Modell-S4-aus-Int} hat man
ein S4"=Kripke"=Modell $M:=M_i$, so dass $M,w\models T(A)$ laut
Prämisse, und somit $M,w\Vdash A$. Und dies bedeutet
$\models_\text{Int} A$, da $M_i,w$ beliebig sind.\,\qedsymbol
\end{Beweis}

\noindent
Infolge muss $A$ genau dann ein intuitionistisches Theorem sein,
wenn $T(A)$ ein Theorem in S4 ist. Die Brücke dazu ergibt sich
unter Hinzuziehung von Korrektheit und Vollständigkeit des jeweiligen
Kalküls. Die Übersetzung $T$ leistet also wirklich, was postuliert
wurde.

\subsection{Zur Auffindung von Kontramodellen}

Mit der Widerlegung der Allgemeingültigkeit einer Formel verhält es sich
in der intuitionistischen Logik komplizierter als in der klassischen,
da ein Gefüge von Welten gefunden werden muss, auf denen die Wertung
unterschiedlich ausfällt. Zurück zur klassischen Logik bzw.
Semantik kommt man bei Begrenzung dieses Gefüges auf eine einzige Welt.
Damit ist auch schon klar, wie eine Formel widerlegt wird, die bereits
in der klassischen Logik nicht allgemeingültig ist. Im Folgenden
interessieren uns also ausschließlich klassische Tautologien, die
vermutlich keine intuitionistischen sind.

Wir notieren kurz $w\refutes A$ für $\lnot (w\models A$).
Definitionsgemäß gilt $w\models A\land B$ genau dann, wenn
$w\models A$ und $w\models B$. Als Kontraposition ergibt sich also,
dass $w\refutes A\land B$ genau dann, wenn $w\refutes A$ oder
$w\refutes B$. Demnach erhalten wir die beiden Regeln
\[\dfrac{w\refutes A}{w\refutes A\land B},\quad \dfrac{w\refutes B}{w\refutes A\land B}.\]
Analog findet sich zur Widerlegung der Disjunktion die Regel
\[\dfrac{w\refutes A\quad\; w\refutes B}{w\refutes A\lor B}.\]
Bei der Subjunktion verhält es sich wieder ein wenig komplizierter.
Definitionsgemäß gilt $w\models A\cond B$ genau dann, wenn
$\forall w'\ge w\colon (w'\models A)\Rightarrow (w'\models B)$. In
Kontraposition gilt $w\refutes A\cond B$ genau dann,
wenn $\exists w'\ge w\colon (w'\models A)\land (w'\refutes B)$. Dies
führt uns also zur Regel
\[\frac{w\le w'\quad\; w'\models A\quad\; w'\refutes B}{w\refutes A\cond B}.\]
Die Negation $\lnot A$ ist hier als Spezialfall $B:=\bot$ enthalten.
Beachtet man, dass $w\refutes\bot$ allgemeingültig ist, vereinfacht sich
die Regel zu
\[\frac{w\le w'\quad\; w'\models A}{w\refutes\lnot A}.\]
Die Regeln erleichtern die Auffindung zielführender Welten und
Wertung.

\strong{Beispiel 1.} Die Allgemeingültigkeit von $A\lor\lnot A$, dem
Satz vom ausgeschlossenen Dritten, ist zu widerlegen. Wir machen dies
bei der Instanz zu $A:=p$. Unter Rückwärtsschließen vermittels der Regeln
ergibt sich die Ableitung:
\[\begin{prooftree}
    \hypo{V(w,p)=0}
  \infer1{w\refutes p}
    \hypo{w\le w'}
      \hypo{V(w',p)=1}
    \infer1{w'\models p}
  \infer2{w\refutes\lnot p}
\infer2{w\refutes p\lor\lnot p}
\end{prooftree}\]
Zu den drei Hypothesen findet sich das Kontramodell
\[W:=\{w,w'\},\;R := \{(w,w),(w,w'),(w',w')\},\;V(w,p):=0,\;V(w',p):=1\]
mit $w\le w'$ als Notation für $R(w,w')$. Dieses widerspricht nicht der
Monotonieforderung an $V$, die besagt, dass zu je zwei Welten $w,w'$
mit $w\le w'$ und $V(w,p)$ auch $V(w',p)$ gilt, bzw. $\lnot V(w,p)$
oder $V(w',p)$ gelten muss.

\strong{Beispiel 2.} Die Allgemeingültigkeit von $\lnot\lnot A\cond A$,
der Beseitigung der Doppelnegation, ist zu widerlegen. Hierfür genügt die
Instanz zu $A:=p$. Kunstloses Rückwärtsschließen führt uns zur Ableitung:
\[\begin{prooftree}
  \hypo{w_1\le w_2}
      \hypo{w_2\le x}
        \hypo{x\le w_3}
          \hypo{V(w_3,p)=1}
        \infer1{w_3\models p}
      \infer2{x\refutes\lnot p}
    \infer2{\forall x\ge w_2\colon x\refutes\lnot p}
  \infer1{w_2\models\lnot\lnot p}
    \hypo{V(w_2,p)=0}
  \infer1{w_2\refutes p}
\infer3{w_1\refutes\lnot\lnot p\cond p}
\end{prooftree}\]
Zu den fünf Hypothesen findet sich nochmals das Kontramodell
\[W:=\{w,w'\},\;R := \{(w,w),(w,w'),(w',w')\},\;V(w,p):=0,\;V(w',p):=1\]
bei Vornahme der Setzungen $w_1:=w$, $w_2:=w$, $w_3:=w'$.

Die Auffindung eines Kontramodells braucht man bei diesen kleinen
Formeln nicht unbedingt händisch rechnen. Insofern die Modelle selbst
klein sind, schafft es der Computer Brute"=Force durch kombinatorische
Aufzählung der Modelle, wenngleich die kombinatorische Komplexität
rasch zunimmt. Da $W$ endlich ist, werden die ominösen
Allquantifizierungen bei der Berechnung der Erfüllung selbst endlich.
In der obigen Ableitung ist bspw. $x\in W$, womit nach Festlegung von $W$
nur noch $x=w$ und $x=w'$ übrig bleiben. Dass jede widerlegbare Formel
ein endliches Kontramodell besitzt, engl. \emph{finite model property},
müsste man streng genommen beweisen. Gerne würde man auch eine Schranke
zur Komplexität in Bezug auf die, noch zu definierende, Komplexität der
Formel haben.

\strong{Beispiel 3.}
Die Allgemeingültigkeit von $(A\cond B\lor C)\cond (A\cond B)\lor (A\cond C)$
ist zu widerlegen. Hierfür genügt die Instanz $A:=p$, $B:=q$, $C:=r$.
Es findet sich die Ableitung:
\[\begin{prooftree}
  \hypo{w_1\le w_2}
    \hypo{\forall x\ge w_2\colon (x\models p)\Rightarrow (x\models q\lor r)}
  \infer1{w_2\models p\cond q\lor r}
  \hypo{w_2\refutes (p\cond q)\lor (p\cond r)}
\infer3{w_1\refutes (p\cond q\lor r)\cond (p\cond q)\lor (p\cond r)}
\end{prooftree}\]
mit dem Teilbaum
\[\begin{prooftree}
    \hypo{w_2\le w_3}
    \hypo{w_3\models p}
    \hypo{w_3\refutes q}
  \infer3{w_2\refutes (p\cond q)}
    \hypo{w_2\le w_4}
    \hypo{w_4\models p}
    \hypo{w_4\refutes r}
  \infer3{w_2\refutes (p\cond r)}
\infer2{w_2\refutes (p\cond q)\lor (p\cond r)}
\end{prooftree}\]
Man beachte hier nun die Allaussage. Näher betrachtet schließt diese
den Umstand $w_3=w_4$ aus, denn wegen $w_3\models p$ zieht sie
$w_3\models q\lor r$ nach sich, aber wir hätten sowohl $w_3\refutes p$
als auch $w_3\refutes r$. Die Lösung besteht darin, die Wertung
bezüglich $q,r$ wechselhaft erfüllt zu haben. Dies erfordert insgesamt
drei unterschiedliche Welten. Zur Verkürzung notieren wir $w\in V(P)$
statt $V(w,P)=1$. Es findet sich das Kontramodell
\begin{gather*}
W:=\{w,w',w''\},\; R := \{(w,w'),(w,w'')\}\cup\text{Reflexivität}(W),\\
V(p):=\{w',w''\},\; V(q):=\{w'\},\; V(r):=\{w''\}
\end{gather*}
bei Vornahme der Setzungen $w_1:=w$, $w_2:=w$, $w_3:=w''$, $w_4:=w'$.

Weitere widerlegbare klassische Tautologien sind
\begin{align*}
& ((A\cond B)\cond A)\cond A, && \text{(Peirce's law)}\\
& (\lnot A\cond A)\cond A, && \text{(Consequentia mirabilis)}\\
& A\lor (A\cond B), && \text{(LEM bei $B:=\bot$)}\\
& (A\cond B)\cond B\lor\lnot A, && \text{(LEM bei $B:=A$)}\\
& \lnot A\lor\lnot\lnot A, && \text{(weak LEM)}\\
& (A\cond B)\lor (B\cond A), && \text{(DPG: Dirk Gently's Principle)}\\
& (\lnot B\cond\lnot A)\cond (A\cond B), && \text{(klassische Kontraposition)}\\
& (\lnot B\cond A)\cond (\lnot A\cond B),\\
& \lnot (A\land B)\cond\lnot A\lor\lnot B.
\end{align*}


\chapter{Mengenlehre}

\section{Grundbegriffe}

\subsection{Der Mengenbegriff}

Eine \emph{Menge}\index{Menge} darf man sich wie einen Beutel
vorstellen, der einzelne Objekte enthält. Die Objekte heißen
\emph{Elemente}\index{Element} der Menge. Jedoch gilt es hierbei zu
beachten, dass es sich mit einer Menge nicht gänzlich wie mit einem
Beutel verhält, in dem dasselbe Objekt mehrmals zu finden sein kann.
\begin{quote}
»\emph{Unter einer Menge verstehen wir jede Zusammenfassung von
bestimmten wohlunterschiedenen Objekten unserer Anschauung oder
unseres Denkens zu einem Ganzen.}«\\
--- Georg Cantor, 1895 (redigiert aus \cite{Cantor})
\end{quote}
Obwohl blumig anmutend, fassen diese Worte das Konzept recht gut
auf den Punkt. Wichtig ist hier das Wort \emph{wohlunterschieden},
das uns zu verstehen gibt, dass ein Element nicht mehrmals in einer
Menge enthalten sein kann. \emph{Eine Menge ist genau dadurch festgelegt,
welche Elemente sie enthält. Sie enthält Elemente weder mehrmals, noch
in einer bestimmten Reihenfolge.}

Es gibt die \emph{leere Menge}\index{leere Menge}, notiert als
$\emptyset$ oder $\{\}$. Man darf sie sich wie einen leeren Beutel
vorstellen. Dagegen enthält die Menge $\{\emptyset\}$ genau ein Element.
Es ist ein Beutel,  der den leeren Beutel enthält.

Wir schreiben kurz $x\in A$ für »$x$ ist ein Element von $A$«,
auch »$x$ gehört zu $A$« oder »$x$ liegt in $A$«.
Es steht $x\notin A$ für $\lnot x\in A$.

Für eine endliche Menge definiert man
\[x\in\{x_1,\ldots,x_n\}\defiff x=x_1\lor\ldots\lor x=x_n.\]

\subsection{Gleichheit von Mengen}

Es wurde gesagt, eine Menge ist dadurch charakterisiert, welche
Elemente sie enthält. Um diesen Gedanke näher zu erfassen, sollten
wir klären, wie es sich mit der Gleichheit von Mengen verhält.
Insofern Mengen durch ihre Elemente bestimmt sind, darf man doch sagen,
zwei Mengen $A,B$ sind genau dann gleich, falls $A,B$ ein Objekt
gemeinsam enthalten oder gemeinsam nicht enthalten. Das heißt, betrachtet
man ein beliebiges Objekt $x$, so ist $x$ genau dann in $A$ enthalten,
wenn $x$ in $B$ enthalten ist.
\begin{Definition}[Gleichheit von Mengen]\label{def:Mengen-Gleichheit}%
\newlinefirst
Für zwei Mengen $A,B$ definiert man
\[A = B \defiff (\forall x\colon x\in A\bicond x\in B).\]
\end{Definition}
Der Leser wird mühelos bestätigen, dass die Gleichheit die Axiome
einer Äquivalenzrelation erfüllt. Wobei der Begriff der Relation
noch zu definieren wäre. Weil die Gesamtheit aller Mengen eine
sogenannte echte Klasse ist, handelt es sich nicht um eine
Äquivalenzrelation im engeren Sinne.

Es ist zulässig, ein Element bei der aufzählenden Angabe einer Menge
mehrmals aufzuführen. Dies ändert allerdings nichts daran, dass ein
Element stets nur einmal in einer Menge enthalten ist. Zum Beispiel
gilt $\{\emptyset,\emptyset\}=\{\emptyset\}$. Mit der Definition der
aufzählenden Angabe und dem Idempotenzgesetz
der Aussagenlogik findet sich nämlich die äquivalente Umformung%
\[x\in\{\emptyset,\emptyset\}\iff x=\emptyset\lor x=\emptyset
\iff x=\emptyset\iff x\in\{\emptyset\}.\]

\subsection{Beschränkte Quantifizierung}

In der Mathematik erstreckt sich die Quantifizierung meist nicht
über das gesamte Diskursuniversum, sondern bleibt auf eine bestimmte
Menge beschränkt. Eine extra Notation macht dies ergonomisch, wobei
eine Erweiterung der logischen Sprache hierfür nicht nötig ist. Die
beschränkte Quantifizierung wird logisch auf eine unbeschränkte
zurückgeführt.

\begin{Definition}[Beschränkte Quantifizierung]\newlinefirst
Für jede Menge $M$ und jede Aussageform $A(x)$ setzt man
\begin{gather*}
(\forall x\in M\colon A(x))\defiff (\forall x\colon x\in M\cond A(x)),\\
(\exists x\in M\colon A(x))\defiff (\exists x\colon x\in M\land A(x)).
\end{gather*}
\end{Definition}
Die Aussage $\forall x\in\emptyset\colon A(x)$ ist allgemeingültig, man
spricht von der \emph{leeren Wahrheit}\index{leere Wahrheit}, engl.
\emph{vacuous truth}\index{vacuous truth}. Via Ex falso quodlibet
erhält man nämlich:
\[
\infer{\vdash\forall x\colon x\in\emptyset\cond A(x)}{
  \infer[\infernote{EFQ}]{x\in\emptyset\vdash A(x)}{
    \infer{x\in\emptyset\vdash\bot}{
      \infer{\vdash\lnot x\in\emptyset}{}
    & \infer{x\in\emptyset\vdash x\in\emptyset}{}}}}
\]
Viele Regeln zur beschränkten Quantifizierung sind analog zu den
Regeln der unbeschränkten. Beispielsweise gilt%
\[(\forall x\in M\colon A(x)\land B(x)) \iff (\forall x\in M\colon A(x))
\land (\forall x\in M\colon B(x)).\]
Man muss allerdings Vorsicht walten lassen. Nicht bei jeder Äquivalenz
liegt eine direkte Analogie vor. Zwar besteht für eine Formel $A$,
in der $x$ nicht frei vorkommt, die Äquivalenz%
\[(\exists x\colon A) \iff A.\]
Die Analogie ist jedoch von der ein klein wenig intrikateren Form
\[(\exists x\in M\colon A) \iff M\ne\emptyset\land A.\]
Diese Beziehung erklärt sich durch die Umformung
\[(\exists x\in M\colon A)\bicond (\exists x\colon x\in M\land A)
\bicond (\exists x\colon x\in M)\land A \bicond M\ne\emptyset\land A.\]
Die letzte Umformung gilt, weil $M\ne\emptyset$ gleichbedeutend mit
$\exists x\colon x\in M$ ist.

\subsection{Komprehension}

Es wird $\{x\mid A(x)\}$ gelesen als »die Klasse der $x$, für die
$A(x)$ gilt« oder »die Menge der $x$, für die $A(x)$ gilt«.
\begin{Definition}[Komprehension]\newlinefirst
Zu einer Aussageform $A(x)$ definiert man die Klasse
$\{x\mid A(x)\}$ gemäß
\[a\in\{x\mid A(x)\}\defiff A(a).\]
\end{Definition}
Man muss mit der Komprehension ein wenig vorsichtig umgehen, denn nicht
jede Klasse ist eine Menge. Die russellsche Klasse $R := \{x\mid x\notin x\}$
ist das klassische Beispiel. Angenommen, $R$ wäre eine
Menge. Dann dürfte man eine Aussage wie $R\in M$ bezüglich
einer Menge $M$ formulieren, also speziell $R\in R$. Gemäß Definition
von $R$ findet sich die folgende Ableitung:
\[
\infer{\vdash R\in R\bicond R\notin R}{
  \infer[\infernote{laut Def.}]{R\in R\vdash R\notin R}{
    \infer{R\in R\vdash R\in R}{}}
& \infer[\infernote{laut Def.}]{R\notin R\vdash R\in R}{
    \infer{R\notin R\vdash R\notin R}{}}}
\]
Für jede Formel $A$ gilt allerdings das Theorem
\[\vdash (A\bicond\lnot A)\cond\bot.\]
Insgesamt ergibt sich so ein Beweis der Kontradiktion. Irgendetwas kann
also nicht gut sein. Diese von Bertrand Russell im Jahre 1901 entdeckte
Verwicklung, die seit jeher den Namen \emph{russellsche Antinomie}%
\index{russellsche Antinomie} trägt, brachte Gottlob Freges
logizistisches Programm in unangenehme Schwierigkeiten. Frege versuchte,
eine Reduktion der Mathematik auf die Logik zu unternehmen, wurde
daraufhin aber von Russell brieflich in Kenntnis gesetzt, dass sein
Werk \emph{Grundgesetze der Arithmetik} in wesentlicher Weise von der
Antinomie unterhöhlt wird. Russell führte das Programm fort, und
veröffentlichte schließlich die \emph{Principia Mathematica}, die der
Antinomie mithilfe einer Typentheorie aus dem Weg geht.

Man begegnet der Problematik, indem man $R$ als eine echte
Klasse ansieht. Für sie darf die Aussage $R\in M$ nicht formuliert
werden. Man definiert weiterhin eine weniger allgemeine Form der
Komprehension, die \emph{Aussonderung}\index{Aussonderung}. Sie verhält
sich gutartig, da sie nicht mehr ermöglicht, als das Ausfiltern von
Elementen aus einer gegebenen Menge.

\begin{Definition}[Aussonderung]\label{def:Aussonderung}\newlinefirst
Zu einer Menge $M$ und einer Aussageform $A(x)$ definiert man
\[a\in\{x\in M\mid A(x)\}\defiff a\in M\land A(a).\]
\end{Definition}

\subsection{Teilmengen}

Gehört jedes Element einer Menge $A$ auch zu einer Menge $B$, nennt
man $A$ eine \emph{Teilmenge}\index{Teilmenge} von $B$. Man sagt auch,
$B$ umfasse $A$, oder $B$ sei eine Obermenge von $A$. Eine \emph{echte
Teilmenge} sei $A$ dann, wenn zusätzlich $A\ne B$ gilt. Die Menge der
geraden Zahlen ist eine echte Teilmenge der ganzen Zahlen. Jede Menge
ist eine Teilmenge von sich selbst, jedoch keine echte.

Die Menge der Quadrate ist eine Teilmenge der Vierecke, genauer eine
Teilmenge der Rechtecke und auch eine Teilmenge der Rhomben.
Weder ist die Menge der Rechtecke eine Teilmenge der Rhomben, noch
ist die Menge der Rhomben eine Teilmenge der Rechtecke. 
Allerdings ist sowohl die Menge der Rechtecke als auch die der Rhomben
eine Teilmenge der Parallelogramme.

\begin{Definition}[Teilmengenbeziehung]\newlinefirst
Man definiert $A\subseteq B$, gelesen »$A$ ist eine Teilmenge von $B$«, als
\[A\subseteq B\defiff (\forall x\colon x\in A\cond x\in B).\]
\end{Definition}
Unschwer bestätigt sich die Äquivalenz
\[A = B \iff A\subseteq B\land B\subseteq A.\]

\begin{Definition}[Potenzmenge]\newlinefirst
Die Potenzmenge\index{Potenzmenge} einer Menge $M$ ist definiert als
\[\mathcal P(M) := \{A\mid A\subseteq M\}.\]
\end{Definition}
Zum Beispiel ist $\mathcal P(\emptyset) = \{\emptyset\}$ und
$\mathcal P(\{0\}) = \{\emptyset, \{0\}\}$. Des Weiteren
\begin{gather*}
\mathcal P(\{0,1\}) = \{\emptyset, \{0\}, \{1\}, \{0,1\}\},\\
\mathcal P(\{0,1,2\}) = \{\emptyset, \{0\}, \{1\}, \{2\}, \{0,1\}, \{0,2\}, \{1,2\}, \{0,1,2\}\}.
\end{gather*}
Die Teilmengenbeziehung darf als eine Art Ordnung zwischen Mengen
betrachtet werden, jedoch nicht als eine Totalordnung, das heißt,
bei einigen Mengen $A,B$ gilt weder $A\subseteq B$ noch $B\subseteq A$.
So ist wie gesagt weder die Menge der Rechtecke eine Teilmenge der
Rhomben, noch umgekehrt.

Die Potenzmenge einer Grundmenge $G$ bildet mit der Teilmengenbeziehung
eine halbgeordnete Menge, engl. \emph{poset} für
\emph{partially ordered set}. Das heißt, alle Mengen
$A,B,C\in\mathcal P(G)$ erfüllen die drei Axiome
\begin{align*}
& A\subseteq A, &&\text{(Reflexivität)}\\
& A\subseteq B\land B\subseteq A\cond A = B, &&\text{(Antisymmetrie)}\\
& A\subseteq B\land B\subseteq C\cond A\subseteq C. &&\text{(Transitivität)}
\end{align*}


\subsection{Mengenoperationen}

\begin{Definition}[Schnitt, Vereinigung, Differenz]\newlinefirst
Zu zwei Mengen $A,B$ definiert man
\begin{align*}
A\cap B &:= \{x\mid x\in A\land x\in B\}, && \text{(Schnittmenge)}\\
A\cup B &:= \{x\mid x\in A\lor x\in B\}, && \text{(Vereinigungsmenge)}\\
A\setminus B &:= \{x\mid x\in A\land x\notin B\}. && \text{(Differenzmenge)}
\end{align*}
\end{Definition}
Sind $A,B$ Teilmengen einer Grundmenge $G$, so sind auch
$A\cap B$, $A\cup B$ und $A\setminus B$ Teilmengen der Grundmenge.
Der Beweis zu $A\cap B\subseteq G$ ist:
\[
\infer{A\subseteq G\vdash A\cap B\subseteq G}{
  \infer{A\subseteq G\vdash\forall x\colon x\in A\cap B\cond x\in G}{
    \infer{A\subseteq G, x\in A\cap B\vdash x\in G}{
      \infer{x\in A\cap B\vdash x\in A}{
        \infer{x\in A\cap B\vdash x\in A\land x\in B}{
          \infer{x\in A\cap B\vdash x\in A\cap B}{}}}
    & \infer{A\subseteq G\vdash x\in A\cond x\in G}{
        \infer{A\subseteq G\vdash \forall x\colon x\in A\cond x\in G}{
          \infer{A\subseteq G\vdash A\subseteq G}{}}}}}}
\]
In so pedantischer Ausführlichkeit findet man Beweise in Büchern
nicht vor. Erstens wird man stillschweigend zulässige Schlussregeln zur
Verkürzung aufstellen. So nimmt der Baum die konzise Form
\[
\begin{array}{@{}l@{\qquad\quad}l}
\infer{A\subseteq G\vdash A\cap B\subseteq G}{
  \infer{A\subseteq G,x\in A\cap B\vdash x\in G}{
    \infer{x\in A\cap B\vdash x\in A}{
      \infer{x\in A\cap B\vdash x\in A\cap B}{}}
  & \infer{A\subseteq G\vdash A\subseteq G}{}}}
&
\infer[\infernote{$\sim$1}]{A\cap B\subseteq G}{
  \infer{x\in G}{
    \infer{x\in A}{
      \infer[\infernote{1}]{x\in A\cap B}{}}
  & A\subseteq G}}
\end{array}
\]
an. Zweitens formuliert der Mathematiker den Beweis meist in Worten:
Um $A\cap B\subseteq G$ zu zeigen, muss $x\in G$ aus $x\in A\cap B$
abgeleitet werden. Mit $x\in A\cap B$ gilt erst recht $x\in A$.
Wegen $A\subseteq G$ ist somit $x\in G$, was zu zeigen war.\,\qedsymbol

\begin{Definition}[Komplement]\newlinefirst
Bezüglich einer Grundmenge $G$ heißt $A^\compc := G\setminus A$
Komplementärmenge.
\end{Definition}
Es zeigt sich elementar, dass die Potenzmenge einer Grundmenge $G$
mit den Operationen $A\cap B$, $A\cup B$ die Axiome
einer booleschen Algebra\index{boolesche Algebra} erfüllt, wobei
$\emptyset$ das Nullelement und $G$ selbst das Einselement ist. Es
gelten somit analoge Regeln wie in der klassischen Aussagenlogik. Wie
die Notation suggeriert, entspricht der Schnitt der Konjunktion, die
Vereinigung der Disjunktion und das Komplement der Negation.

Wie in jedem Verband ist in einer booleschen Algebra
$(M,\wedge,\vee)$ für $a,b\in M$ eine Halbordnung $a\le b$
definiert, indem $a\le b$ und $a\wedge b = a$ als äquivalent
angesehen werden. Bei der Mengenalgebra entpuppt sie sich
als die Teilmengenrelation. Das heißt, es gilt%
\[A\subseteq B\iff A\cap B = A\iff A\cup B = B.\]
Für mehrere Mengen definiert man
\[\bigcap_{i=1}^n A_i := A_1\cap A_2\cap\ldots\cap A_n,\qquad
\bigcup_{i=1}^n A_i := A_1\cup A_2\cup\ldots\cup A_n.\]
Pedantiker mögen die Schreibweise mit den Auslassungspunkten als
ungenau empfinden. Zufriedenstellend ist für sie die Erklärung, dass
es sich um die rekursive Festlegung%
\[\bigcap_{i=1}^1 A_i := A_1,\qquad
\bigcap_{i=1}^n A_i := A_n\cap\bigcap_{i=1}^{n-1} A_i\]
handelt. Regeln wie $B\cup\bigcap_{i=1}^n A_i
= \bigcap_{i=1}^n (B\cup A_i)$ kann man nun per Induktion über $n$
beweisen. Es geht fast trivial vonstatten, sobald man das Prinzip
verstanden hat. Im Wesentlichen weiten sich hier die Regeln der
booleschen Algebra von den zweistelligen auf die mehrstelligen
Operationen aus.

\begin{Definition}[Allgemeine Vereinigung]\newlinefirst
Sei $M$ eine Menge von Mengen. Die Vereinigung der $A\in M$ ist
\[\bigcup M = \bigcup_{A\in M} A := \{x\mid\exists A\in M\colon x\in A\}.\]
\end{Definition}
Für $M=\emptyset$ ist $\bigcup M = \emptyset$. Die Disjunktion findet
ihre Entsprechung genau in der Vereinigung
von zwei Mengen. Dazu passend findet der Existenzquantor seine
Entsprechung genau in der Vereinigung beliebig vieler Mengen.
Aus diesem Grund weiten sich die Regeln der booleschen Algebra
auf die allgemeine Vereinigung aus. Zum Beispiel lautet das
Distributivgesetz für Mengen%
\[B\cap\bigcup_{A\in M} A = \bigcup_{A\in M}(B\cap A).\]
Entfaltung der Definition führt nämlich zur logischen Äquivalenz
\[x\in B\land(\exists A\in M\colon x\in A) \iff (\exists A\in M\colon x\in B\land x\in A).\]
Ihr Beweis gelingt mühelos.

\begin{Definition}[Allgemeiner Schnitt]\newlinefirst
Sei $M$ eine nichtleere Menge von Mengen. Der Schnitt der $A\in M$ ist%
\[\bigcap M = \bigcap_{A\in M} A := \{x\mid\forall A\in M\colon x\in A\}.\]
\end{Definition}
Im Gegensatz zur Vereinigung wurde der Schnitt $\bigcap M$ 
für $M=\emptyset$ undefiniert gelassen. Hier gibt es zwei Möglichkeiten.
Zum einen könnte man die Bedingung $M\ne\emptyset$ einfach fallen
lassen, dann ergibt im allgemeinen Mengenuniversum beim leeren Schnitt
die Allklasse $\{x\mid\top\}$, die jedoch keine Menge ist.

Aus diesen Grund gibt es noch die alternative Definition
\[\bigcap M :=
\{x\in G\mid\forall A\in M\colon x\in A\}.\]
Hierzu ist eine Grundmenge $G$ festzulegen, so dass $M\subseteq\mathcal P(G)$
gilt, oder man setzt $G:=\bigcup M$, wobei sich da die Frage
nach der Nützlichkeit stellt.

Eine Menge von Mengen nennt man ein \emph{Mengensystem}\index{Mengensystem},
wobei aber einige Autoren diese Begrifflichkeit für eine Familie von
Mengen benutzen, die von einer Menge von Mengen zu unterscheiden ist.
Eine Familie stellt eine Verallgemeinerung einer Folge von Mengen dar.
In ihr darf dieselbe Menge mehrmals vorkommen. Man kann Schnitt und
Vereinigung auch für Familien definieren, was aber eigentlich keine
wesentliche Verallgemeinerung zu den obigen Festlegungen darstellt,
wie ich im Folgenden diskutieren möchte.

Eine \emph{Familie}\index{Familie} $(A_i)$ von Mengen $A_i$ mit $i\in I$
ist eine Abbildung $A\colon I\to Z$, wobei $Z$ eine Zielmenge ist,
welche die $A_i$ als Elemente enthält. Die Menge $I$ wird in diesem
Zusammenhang auch \emph{Indexmenge}\index{Indexmenge} genannt. Man
definiert%
\[\bigcup_{i\in I} A_i := \bigcup A(I)
= \bigcup\{X\mid\exists i\in I\colon X=A_i\} = \{x\mid\exists i\in I\colon x\in A_i\},\]
wobei mit $A(I)$ das Bild von $I$ unter $A$ gemeint ist. Man bekommt%
\begin{align*}
\smash{\bigcup_{i\in I} A_i}
&= \{x\mid\exists X\colon X\in \{X\mid\exists i\in I\colon X=A_i\}\land x\in X\}\\
&= \{x\mid\exists X\colon (\exists i\in I\colon X=A_i)\land x\in X\}\\
&= \{x\mid\exists X\colon \exists i\in I\colon X=A_i\land x\in X\}
= \{x\mid\exists i\in I\colon x\in A_i\}.
\end{align*}
Für $I\ne\emptyset$ definiert man entsprechend
\[\bigcap_{i\in I} A_i := \bigcap A(I) = \{x\mid\forall i\in I\colon x\in A_i\}.\]
Die Operation über eine Familie $(A_i)_{i\in I}$ kann also
auf die jeweilige Operation über das System $A(I)$ zurückgeführt
werden.

Später nützlich ist der

\begin{Satz}\label{Index-in-Schnitt}
Es gilt $\bigcup_{i\in I\cap J} A_i = (\bigcup_{i\in I} A_i)\cap
(\bigcup_{i\in J} A_j)$.
\end{Satz}
\begin{Beweis}
Es findet sich die äquivalente Umformung
\begin{align*}\textstyle
x\in\bigcup_{i\in I\cap J} A_i &\iff (\exists i\colon i\in I\cap J\land x\in A_i)\\
&\iff(\exists i\colon (i\in I\land x\in A_i)\land (i\in J\land x\in A_i))\\
&\iff(\exists i\colon i\in I\land x\in A_i)\land (\exists i\colon i\in J\land x\in A_i)\\
&\iff\textstyle x\in\bigcup_{i\in I} A_i\land x\in\bigcup_{i\in J} A_i\\
&\iff\textstyle x\in(\bigcup_{i\in I} A_i)\cap (\bigcup_{i\in J} A_i).\,\qedsymbol
\end{align*}
\end{Beweis}

\noindent
Gelegentlich hat man es mit einer \emph{disjunkten Vereinigung}%
\index{disjunkte Vereinigung} zu tun. Sie ist bedeutsam in der
Theorie der Kardinalzahlen und in der Informatik bei algebraischen
Datentypen sowie deren Bezug zur Beweistheorie. Genauer gesagt hantiert
man in der Informatik nicht mit Mengen, sondern mit Typen, die sich
in gewissen Zügen analog zu Mengen verhalten. Die disjunkte Vereinigung
zweier Mengen kennzeichnet jedes Element vor der Vereinigung mit einem
Tag, das die Information liefert, aus welcher der Mengen es entstammt.
Man setzt%
\[A\sqcup B := \{(0,a)\mid a\in A\}\cup\{(1,b)\mid b\in B\}.\]
Die Zahlen $0,1$ sind hier die \emph{Tags} oder \emph{Diskriminatoren}.
Anstelle der Zahlen könnten genauso gut zwei beliebige unterschiedliche
Elemente als Tags verwendet werden. Beispielsweise ginge auch%
\[A\sqcup B = \{(\text{Grün},a)\mid a\in A\}\cup\{(\text{Blau},b)\mid b\in B\}.\]
In der Informatik verwendet man gern left, right als Tags.

Mit jeder disjunkten Vereinigung ist eine Fallunterscheidung verbunden.
Liegt ein $x\in A\sqcup B$ vor, so muss entweder $x=(0,a)$ für ein
$a\in A$ oder $x=(1,b)$ für ein $b\in B$ sein. Bei einer gewöhnlichen
Vereinigung besteht dagegen kein ausschließendes Oder.

Wir fassen zwei Objekte $x,y$ zu einem \emph{geordneten Paar}%
\index{Paar}\index{geordnetes Paar} $(x,y)$ zusammen. Die beiden
Objekte müssen nicht unbedingt etwas miteinander zu tun haben, sie
dürfen völlig verschiedener Art sein. Im Unterschied zu einer Menge
spielt bei Paaren die Reihenfolge eine Rolle, auch darf dasselbe Objekt
zweimal vorkommen. Im Paar $t=(x,y)$ ist genau die Information über
$x,y$ enthalten. Das heißt, es lassen sich $x,y$ aus dem Paar
extrahieren. Man schreibt dafür $t_1=x$ und $t_2=y$. Die Schreibweise
$t_i$ heißt \emph{Indizierung}, es ist darin $i$ der
\emph{Index}\index{Index}.

Zwei Paare seien definitionsgemäß genau dann gleich, wenn sie
komponentenweise gleich sind,%
\[(x,y) = (x',y')\defiff x=x'\land y=y'.\]
Die genannten Eigenschaften charakterisieren den Begriff Paar
im Wesentlichen, mehr müssen wir nicht wissen. Man hat sich trotzdem
auch mal überlegt, wie Paare in der reinen Mengenlehre dargestellt
werden können, wo alle Objekte Mengen sein sollen. Nach Kuratowski
sind Paare kodierbar als%
\[(x,y) := \{\{x\},\{x,y\}\}.\]
Unter dem allgemeineren Begriff \emph{Tupel}\index{Tupel} fassen
wir eine beliebige endliche Zahl von Objekten in einer bestimmten
Reihenfolge zusammen. Tupel aus drei Objekten heißen \emph{Tripel},
die aus vier heißen \emph{Quadrupel}. Analog zu den Paaren ist ihre
Gleichheit definiert als%
\[(x_1,\ldots,x_n) = (x_1',\ldots,x_n')
\defiff x_1=x_1'\land\ldots\land x_n=x_n'.\]
Zu zwei Mengen $X,Y$ kann man nun die Menge aller Paare betrachten,
deren erste Komponente $X$ entstammt, und deren zweite
Komponente $Y$ entstammt. Sie heißt \emph{Produktmenge}\index{Produktmenge}
oder \emph{kartesisches Produkt}\index{kartesisches Produkt}
der Mengen $X,Y$.
\begin{Definition}[Produktmenge]\newlinefirst
Das kartesische Produkt zweier Mengen $X,Y$ ist die Menge
\[X\times Y := \{(x,y)\mid x\in X\land y\in Y\}.\]
\end{Definition}
Genau genommen handelt es sich hierbei um eine Bildmenge, was bedeutet,
dass sich im rechten Term Existenzquantoren verstecken. Ausführlich
ausgeschrieben lautet der Term%
\begin{align*}
X\times Y &= \{t\mid\exists x\in X\colon\exists y\in Y\colon t=(x,y)\}\\
&= \{t\mid\exists x\colon\exists y\colon x\in X\land y\in Y\land t=(x,y)\}.
\end{align*}
Wie jede Bildmenge ist das Produkt als Vereinigung darstellbar. Es ist%
\[X\times Y = \bigcup_{x\in X}\bigcup_{y\in Y}\{(x,y)\}.\]
Die beiden kurzen Identitäten $X\times\emptyset = \emptyset$ und
$\emptyset\times Y = \emptyset$ gehen unmittelbar aus der Definition
hervor. Es verhält sich analog wie mit der Multiplikation einer
Zahl mit null. Eine sinnvolle Sichtweise, wie die Theorie der
Kardinalzahlen lehrt.

\begin{Satz}
Ist $A\subseteq X$ und $B\subseteq Y$, dann ist
$A\times B\subseteq X\times Y$.
\end{Satz}
\begin{Beweis}
Es liege $t$ in $A\times B$. Laut Definition existieren mithin
$a\in A$ und $b\in B$, so dass $t=(a,b)$. Wegen $A\subseteq X$ ist aber
auch $a\in X$ und wegen $B\subseteq Y$ ist auch $b\in Y$. Daher existieren
$a\in X$ und $b\in Y$, so dass $t=(a,b)$. Gemäß Definition heißt das,
$t\in X\times Y$. Gemäß Definition ist $A\times B$ somit eine Teilmenge
von $X\times Y$.\;\qedsymbol
\end{Beweis}

\begin{Satz}
Es gilt $X\times (A\cap B) = (X\times A)\cap (X\times B)$.
\end{Satz}
\begin{Beweis}
Es ginge zu Fuß. Aber Satz \ref{Index-in-Schnitt} ermöglicht die
kurze Termumformung%
\begin{align*}
& X\times (A\cap B) = \bigcup_{x\in X}\bigcup_{y\in A\cap B} \{(x,y)\}
= \bigcup_{x\in X}\Big(\Big(\bigcup_{y\in A}\{(x,y)\}\Big)\cap \Big(\bigcup_{y\in B}\{(x,y)\}\Big)\Big)\\
&\quad = \Big(\bigcup_{x\in X}\bigcup_{y\in A}\{(x,y)\}\Big)\cap\Big(\bigcup_{x\in X}\bigcup_{y\in B}\{(x,y)\}\Big)
= (X\times A)\cap (X\times B).\,\qedsymbol
\end{align*}
\end{Beweis}

\newpage
\section{Abbildungen}

\subsection{Der Abbildungsbegriff}

Unter einer \emph{Abbildung}\index{Abbildung}, auch \emph{Funktion}%
\index{Funktion} genannt, verstehen wir ganz allgemein eine Zuordnung von
Elementen einer Definitionsmenge zu Elementen einer Zielmenge, bei der
zu \emph{jedem} Element der Definitionsmenge \emph{genau ein} Element
der Zielmenge gehört. Es hat allerdings ein wenig gedauert, bis
diese Vorstellung als die Förderliche erkannt wurde.

In der historischen Entwicklung ging ihr die speziellere Vorstellung
voran, dass eine Größe, etwa eine physikalische Größe, in einer
rechnerischen Abhängigkeit von einer anderen Größe steht. So steht
die Periodendauer der Schwingung eines Pendels in einer bestimmten
rechnerischen Abhängigkeit von der Pendellänge.

Eine recht pragmatische Vorstellung von einer Funktion vermittelt das Modell
der Black Box, das soll eine Rechenmaschine sein, deren innere Mechaniken
bzw. Elektroniken unbekannt bleiben. Speist man ein Argument $x$ in die
Black Box ein, spuckt sie daraufhin einen Wert $f(x)$ aus. Speist man
abermals dasselbe Argument ein, spuckt sie abermals denselben Wert aus.

Der Begriff der Abbildung ist für die Mathematik zentral.

\begin{Definition}[Abbildung]\newlinefirst
Eine Abbildung $f\colon X\to Y$ ist eine Relation $f=(X,Y,G)$ mit
$G\subseteq X\times Y$, die die beiden Eigenschaften
\begin{gather*}
\forall x\in X\colon\exists y\in Y\colon (x,y)\in G,\\
\forall x\in X\colon\forall y,y'\in Y\colon (x,y)\in G\land (x,y')\in G\cond y=y'
\end{gather*}
erfüllt. Man nennt $G$ den Graph, $X$ die Definitions- und
$Y$ die Zielmenge.
\end{Definition}
Genau genommen möchte man die Abbildung $f$ eigentlich nicht als mit
dem Tripel $(X,Y,G)$ identisch sehen, sondern als dasjenige Objekt,
dessen innere Information aus diesem Tripel besteht.

Zuweilen wird der Graph von $f$ auch einfach als $f$ bezeichnet.
Sollte dies verwirrend sein, schreibt man ausführlicher $G_f$ oder
$\mathrm{Graph}(f)$.

Statt $(x,y)\in G_f$ schreibt man üblicherweise $y=f(x)$. Es
wird $f(x)$ gelesen als »$f$ von $x$« oder »das Bild von $x$ unter
$f$«. Es wird $f\colon X\to Y$ gelesen als »$f$ ist eine Abbildung
von $X$ nach $Y$«. Besitzt ein $y\in Y$ ein $x\in X$ mit $y=f(x)$,
nennt man $x$ ein Urbildelement zu $y$.

\newpage
\subsection{Bild, Urbild}

Wird ein jedes Element einer Teilmenge der Definitionsmenge in eine
Abbildung geschickt, formen die Werte eine neue Menge, die
\emph{Bildmenge}\index{Bildmenge} der Teilmenge unter der Abbildung.

\begin{Definition}[Bild]\newlinefirst
Die Bildmenge einer Menge $A\subseteq X$ unter der Abbildung
$f\colon X\to Y$ ist
\[f(A) := \{f(x)\mid x\in A\} = \{y\mid\exists x\in A\colon y=f(x)\}.\]
\end{Definition}
Insofern $y=f(x)$ als äquivalent zu $y\in\{f(x)\}$ befunden wird, ergibt
sich auch die Darstellung $f(A) = \bigcup_{x\in A} \{f(x)\}$.
Insbesondere in Programmiersprachen treten endliche Mengen als Objekte
auf. Die Berechnung verläuft mithin gemäß der rekursiven Festlegung%
\[f(\emptyset) := \emptyset,\qquad f(\{x_1,\ldots,x_n\}) :=
f(\{x_1,\ldots,x_{n-1}\})\cup \{f(x_n)\},\]
für die auch die Bezeichnung $\mathrm{map}(f,A)$ gebräuchlich ist.

Die analytische Geometrie sieht Figuren als Punktmengen. Abbildungen
transformieren die Punktmengen, woraus neue Figuren hervorgehen.
Ein einfaches Beispiel bietet die Abbildung%
\[f\colon\R\to\R^2,\quad f(t) :=
\begin{pmatrix}p_x + v_x t\\ p_y + v_y t\end{pmatrix}.\]
In der Hinsicht, dass die reellen Zahlen $\R$ als eine Zahlengerade
aufgefasst werden, formt ihr Bild $f(\R)$ eine Gerade der euklidischen
Ebene. Die passende Wahl der Parameter $(p_x,p_y)$ als Basispunkt
und $(v_x,v_y)$ als Geschwindigkeitsvektor gestattet es, jede
beliebige Gerade der Ebene zu beschreiben. Man nennt in der üblichen
Terminologie aber auch $t$ den Parameter und $f$ eine Parametergerade,
wobei das Bild $f(\R)$ auch Spur genannt wird. Diese Sprechweisen
entspringen der Sichtweise, dass $f(t)$ ein durch die Zeit $t$
parametrisierter Punkt ist, der mit dem Lauf der Zeit eine Spur zieht.
Der Geschwindigkeitsvektor charakterisiert die Bewegung des Punktes,
für die Spur ist dagegen nur dessen Richtung von Bedeutung.

Das \emph{Urbild}\index{Urbild} eines Wertes gibt Auskunft, welche
Elemente der Definitionsmenge unter der Abbildung zu dem Wert führen.
Allgemeiner gibt das Urbild einer Menge von Werten Auskunft, welche
Elemente der Definitionsmenge in die Menge führen.

\begin{Definition}[Urbild]\newlinefirst
Das Urbild einer Menge $B$ unter $f\colon X\to Y$ ist
\[f^{-1}(B) := \{x\in X\mid f(x)\in B\} = \{x\in X\mid \exists y\in B\colon y=f(x)\}.\]
\end{Definition}
Dem Wesen des Abbildungsbegriffs entspringend, zeichnet sich die
Urbildoperation durch Verträglichkeit mit den Mengenoperationen aus,
was bei der Bildoperation nur zum Teil stimmt.

\begin{Satz}
Für jede Abbildung $f$ und beliebige Mengen $A,B,A_i$ gilt
\begin{align*}
f^{-1}(A\cap B) &= f^{-1}(A)\cap f^{-1}(B), &&\textstyle
f^{-1}(\bigcap_{i\in I} A_i) = \bigcap_{i\in I} f^{-1}(A_i),\\
f^{-1}(A\cup B) &= f^{-1}(A)\cup f^{-1}(B), &&\textstyle
f^{-1}(\bigcup_{i\in I} A_i) = \bigcup_{i\in I} f^{-1}(A_i).
\end{align*}
\end{Satz}
\begin{Beweis}
Den Beweis verschafft die äquivalente Umformung
\begin{align*}
&\textstyle x\in f^{-1}(\bigcap_{i\in I} A_i)
\iff f(x)\in\bigcap_{i\in I} A_i
\iff (\forall i\in I\colon f(x)\in A_i)\\
&\textstyle\quad\iff (\forall i\in I\colon x\in f^{-1}(A_i))
\iff x\in\bigcap_{i\in I} f^{-1}(A_i).
\end{align*}
Bei der Vereinigung verläuft die Umformung analog.\,\qedsymbol
\end{Beweis}

\noindent
Der Satz lehrt die wichtige Folgerung, dass die Urbilder zweier
disjunkter Mengen ebenfalls disjunkt sind. Allgemein überführt die
Urbildoperation eine disjunkte Zerlegung einer Menge in disjunkte
Urbilder. Ist die Zielmenge am feinsten zerlegt, besteht sie also
aus einelementigen Mengen, nennt man die Urbilder dieser Mengen
auch Fasern. Man beachte, dass ein Urbild oder eine Faser leer
sein kann.

Ist die Zielmenge bezüglich einer Ordnungsrelation angeordnet,
nennt man die Fasern auch \emph{Niveaumengen}\index{Niveaumenge}.
Diese Begrifflichkeit betrifft vor allem die Funktionen $f\colon X\to\R$
mit $X\subseteq\R^n$. Eine Faser $f^{-1}(\{c\})$, auch als $f(x)=c$
beschrieben, nennt man auch eine \emph{implizite Funktion}. Sie sind
in der analytischen Geometrie und der mehrdimensionalen Analysis von
großer Bedeutung.

\subsection{Komposition}

\begin{Definition}[Komposition]\newlinefirst
Sei $f\colon X\to Y$ und $g\colon Y\to Z$. Ihre Verkettung,
gelesen »$g$ nach $f$«, ist%
\[(g\circ f)\colon X\to Z,\quad (g\circ f)(x):=g(f(x)).\]
\end{Definition}

\noindent
Oft liegt die Situation $f\colon X\to Y$ und $g\colon Y'\to Z$
mit $Y\subseteq Y'$ vor. Das ist aber nicht weiter schlimm. Es
darf dann $g\circ f := g|_Y\circ f$ gesetzt werden, wobei mit $g|_Y$ die
Einschränkung des Definitionsbereichs von $g$ auf $Y$ gemeint ist.

\begin{Definition}[Einschränkung]\newlinefirst
Sei $f\colon X\to Y$ und $A\subseteq X$. Die Einschränkung von
$f$ auf $A$ ist%
\[f|_A\colon A\to Y,\quad f|_A(x):=f(x).\]
\end{Definition}

\noindent
Die Abbildung $\id_X\colon X\to X$ mit $\id_X(x):=x$ heißt
\emph{identische Abbildung}\index{identische Abbildung}. Sie verhält
sich bei der Komposition neutral, das heißt, bezüglich
$f\colon X\to Y$ gilt $f\circ\id_X = f$ und $\id_Y\circ f=f$.

\subsection{Injektionen, Surjektionen, Bijektionen}

\begin{Definition}[Injektion]\newlinefirst
Eine Abbildung $f\colon X\to Y$ heißt injektiv\index{injektiv}, wenn%
\[\forall x,x'\in X\colon f(x)=f(x')\cond x=x'.\]
\end{Definition}

\noindent
Erinnern wir uns an den Abschnitt \emph{Logik mit Gleichheit}, fällt
auf, dass die Definition der Injektivität als Umkehrung der
Ersetzungsregel betrachtbar ist. Das heißt, für eine auf den Werten
der Terme $t,t'$ definierte Injektion $f$ gilt die Äquivalenz%
\[t=t' \iff f(t)=f(t').\]
Die Injektionen vermitteln somit genau die \emph{Äquivalenzumformungen}%
\index{Aequivalenzumformung@Äquivalenzumforumung} von Gleichungen.
Wie sich aus Def. \ref{def:Mengen-Gleichheit} in Verbindung mit
Def. \ref{def:Aussonderung} ergibt, erfährt die Lösungsmenge einer
Bestimmungsgleichung durch sie keine Veränderung. Man notiert%
\[L := \{x\in G\mid t=t'\} = \{x\in G\mid f(t)=f(t')\}\]
für die Lösungsmenge $L$ der Gleichung $t=t'$ in der Variable $x$,
die die Werte der Grundmenge $G$ durchläuft. Aufgrund dessen schaffen
sie ein wesentliches Werkzeug zum Lösen von Bestimmungsgleichungen.

Es muss $X$ bei $f\colon X\to Y$ nicht notwendigerweise die Grundmenge
sein. Wichtig ist allein, dass die Terme $t,t'$ ausschließlich Werte in
$X$ annehmen können. Beispielsweise ist das Quadrieren auf den reellen
Zahlen zwar nicht injektiv, bei Einschränkung der Definitionsmenge auf
die nichtnegativen reellen Zahlen allerdings schon. So initiiert es
die äquivalente Umformung%
\[|x| = |x-2| \bicond x^2 = (x-2)^2 = x^2-4x+4 \bicond
0 = -4x+4 \bicond x = 1.\]
Außerdem darf die Variable der Bestimmungsgleichung in einer
Äquivalenzumformung als Parameter auftauchen. Die injektive Funktion%
\[f_a\colon\R\to\R,\quad f_a(\xi) := \xi + a\]
beschreibt zum Beispiel den Sachverhalt, dass eine Zahl $a$ zu beiden
Seiten einer Gleichung addiert werden darf,%
\[t = t' \iff t + a = t' + a.\]
Hier darf insbesondere auch $a:=x$ als Parameter eingesetzt werden. Bei%
\[f_a\colon\R\to\R,\quad f_a(\xi) := a\xi\]
gilt es allerdings $a\ne 0$ zu berücksichtigen. Das heißt, die Setzung
$a:=x$ liefert nur dann eine Äquivalenzumformung, wenn die Grundmenge,
die $x$ durchläuft, eine Teilmenge von $\R\setminus\{0\}$ ist. Andernfalls
müsste man diesen Umstand durch eine Fallunterscheidung künstlich herstellen.

Erwähnenswert ist weiterhin, dass die gemachten Begriffe bereits
Gleichungen in mehreren Variablen und Gleichungssysteme%
\index{Gleichungssystem} umfassen. Eine Gleichung in zwei Variablen
liegt vor, wenn die Grundmenge aus Paaren besteht. Ein System von zwei
Gleichungen liegt vor, wenn die die Terme $t,t'$ jeweils ein Paar zum
Wert haben. Diese formale Beschreibung müsste man allerdings als ein
wenig oberflächlich befinden, kämen nicht weitere Erörterungen hinzu.
Tiefersinnig wäre es an sich ohne Pointe, tauchte nur eine der Variablen in
der Gleichung auf. Gleichermaßen mag ein System in zwei Variablen
erst reizvoll sein, wenn die Gleichungen in nichttrivialer Weise
miteinander verwoben sind.

% Stellt eine Verbindung zu den Konzepten der Katengorientheorie her.
\begin{Satz}
Eine Abbildung $g$ mit $g\circ f = \id$ heißt Linksinverse von $f$.
Eine Abbildung $f\colon X\to Y$ mit nichtleerem $X$ ist genau dann
injektiv, wenn sie mindestens eine Linksinverse besitzt.
\end{Satz}
\begin{Beweis}
Sei $g$ eine Linksinverse von $f$. Seien $x,x'$ fest, aber beliebig,
und sei $f(x)=f(x')$. Mithin gilt $g(f(x))=g(f(x'))$. Weil $g$ eine
Linksinverse ist, hat man aber $g(f(x))=x$ und $g(f(x'))=x'$, womit
sich wie gewünscht $x=x'$ ergibt.

Sei $f$ injektiv. Mit $X\ne\emptyset$ liegt irgendein $a\in X$ vor.
Sofern $y\in f(X)$ ist, liegt außerdem ein $x$ mit $y=f(x)$ vor.
Es wird $g$ festgelegt per Fallunterscheidung
\[g(y) := \begin{cases}
x, & \text{wenn}\;y\in f(X),\\
a, & \text{wenn}\;y\notin f(X).
\end{cases}\]
Sie ist verhält sich so, dass $g(f(x))=x$ für jedes $x\in X$ gilt,
denn mit $x\in X$ ist $f(x)\in f(X)$. Somit existiert mit $g$ eine
Linksinverse.\,\qedsymbol
\end{Beweis}

\begin{Definition}[Surjektion]\newlinefirst
Eine Abbildung $f\colon X\to Y$ heißt surjektiv\index{surjektiv},
wenn $f(X)=Y$ ist.
\end{Definition}

\noindent
Weil $f(X)\subseteq Y$ allgemeingültig ist, genügt es generell,
$Y\subseteq f(X)$ zu zeigen.

\begin{Definition}[Bijektion]\newlinefirst
Eine Abbildung heißt bijektiv, wenn sie sowohl injektiv als auch
surjektiv ist.
\end{Definition}

\begin{Satz}
Jede Bijektion $f$ besitzt eine eindeutig bestimmte Abbildung, welche
sowohl ihre einzige Linksinverse als auch ihre einzige Rechtsinverse
ist. Man nennt sie die Umkehrabbildung $f^{-1}$.
\end{Satz}
\begin{Beweis}
Sei $f\colon X\to Y$ bijektiv. Es existiert somit mindestens eine
Linksinverse $g$ und mindestens eine Rechtsinverse $h$. Weil die
Verkettung das Assoziativgesetz erfüllt, darf man rechnen
\[g = g\circ\id_Y = g\circ (f\circ h) = (g\circ f)\circ h
= {\id_X}\circ h = h.\]
Sei nun $g'$ eine weitere Linksinverse und $h'$ eine weitere
Rechtsinverse. Wiederholt man die obige Rechnung abermals, erhält man
$g'=h$ und $g'=h'$. Ergo gilt $g = h = g' = h'$.\,\qedsymbol
\end{Beweis}

\subsection{Allgemeines Mengenprodukt}

Nachdem nun der Begriff der Abbildung bereits erklärt wurde, kann eine
Begriffsverallgemeinerung des kartesischen Produktes erörtert werden,
die Tupel mit Folgen und Funktionen in Beziehung setzt. Eine Folge
darf man einerseits als Funktion betrachten, deren Definitionsbereich
die natürlichen Zahlen sind. Andererseits ist eine Folge indizierbar wie
ein Tupel. Das bringt uns auf die Idee, Tupel wie Folgen als Funktionen
aufzufassen.

Die Ziffern der Zahl 2520 sind in Little"=Endian"=Konvention das Tupel
\[t=(0,2,5,2)\in\N\times\N\times\N\times\N.\]
Die Indexmenge $I:=\{0,1,2,3\}$ enthält die Stellenwerte der Ziffern.
Es verhält sich nun so, dass die Funktion
\[f\colon I\to\N,\quad f(0) := 0,\; f(1) := 2,\; f(2) := 5,\; f(3) := 2\]
genau die Information enthält, die das Tupel charakterisiert. Es ist
$f(i) = t_i$ die Indizierung des Tupels. So gesehen ist jedes Tupel durch
eine Funktion kodiert. Führt man diese Überlegung fort, gelangt man
schließlich zum allgemeinen kartesischen Produkt
\[\prod_{i\in I} X_i := \{f\colon I\to\bigcup_{i\in I} X_i\mid
\forall i\in I\colon f(i)\in X_i\}.\]
Der Formelsalat sagt im Wesentlichen nur, dass eine Abbildung
$f\in\prod_{i\in I} X_i$ einen in $X_i$ liegenden Funktionswert
$f(i)$ besitzt.

Diese allgemeine Produkte von Mengen treten in der Mathematik nur sehr
sporadisch in Erscheinung. Eine Rolle spielen sie in der Theorie der
Kardinalzahlen. In der abhängigen Typentheorie besitzen sie allerdings ein
direktes allgegenwärtiges Analogon, den \emph{Typ abhängiger Funktionen}.

Bei einer Gleichsetzung wie $\prod_{i\in\{1,2\}} X_i\cong X_1\times X_2$ ist
zu beachten, dass es sich eigentlich nicht um eine Gleichung handelt, denn
Abbildungen sind ja nicht dasselbe wie Paare. Allerdings gehört zu jedem
Paar wie gesagt in kanonischer Weise genau eine gleichartige Abbildung.
Es besteht eine Isomorphie, was durch die Tilde über dem Gleichheitszeichen
angedeutet wird. Damit verbunden liegt ein Isomorphismus vor, mit dem
sich ein Objekt der einen Seite in ein Objekt der anderen Seite
überführen lässt. Das ist so ähnlich wie die Übersetzung eines Textes
von der einen in die andere Sprache. Die Gestalt ändert sich, aber der
Inhalt bleibt gleich.

Intuitiv sollte $\prod_{i\in I} X_i$ nichtleer sein, sofern jedes $X_i$
nichtleer ist. Dennoch lässt sich dieser Sachverhalt nicht ohne Weiteres
ableiten. Es handelt sich um ein Axiom der Mengenlehre, das sogenannte
\emph{Auswahlaxiom}\index{Auswahlaxiom}. Es besagt, dass zu jeder
Familie $(X_i)$ nichtleerer Mengen eine Auswahlfunktion $f$ existiert,
die zu jeder Menge $X_i$ ein Element $f(i)\in X_i$ auswählt.
Eine gründliche Untersuchung des Auswahlaxioms zeigt auf, dass es sich
mit den Unendlichkeiten schwierig verhält. Unter bestimmten Umständen
ist das Axiom nicht vonnöten. Außerdem, so lässt sich zeigen, impliziert
es den Satz vom ausgeschlossenen Dritten.

\section{Relationen}

\subsection{Relationen im Allgemeinen}

Beim Erlangen von tieferliegenden Einsichten in Probleme und
Zusammenhänge spielt das Aufspüren und Klären unterschiedlicher
Beziehungen eine wesentliche Rolle. Zwei Objekte können auf
unterschiedliche Art und Weise in Beziehung stehen. In der Mathematik
beschäftigt man sich mit Beziehungen, die zu den zwei Objekten eine
Aussage trifft, der ein Wahrheitsgehalt beigemessen wird. Zwei Zahlen
$x,y$ sind gleich. Eine Zahl $x$ ist kleiner als eine Zahl $y$. Der
Betrag der Differenz zweier Zahlen $x,y$ ist kleiner als eine bstimmte
Konstante. Der Abstand der Punkte $x=(x_1,x_2)$ und $y=(y_1,y_2)$ ist
kleiner als eine bestimmte Konstante. Zwei Stühle $x,y$ eines
Stuhlkreises sind benachbart. Die Schüler $x,y$ gehen in dieselbe
Klasse. Es gibt einen Weg, der vom Ort $x$ zum Ort $y$ führt.

Die eindrückliche Vielfalt möglicher Beziehungen macht es ja
unerlässlich, eine allgemeine Auffassung von ihnen zu erhalten.
Eine zweistellige Relation $R$ schafft ein Beziehungsgefüge zwischen
zwei Mengen $X,Y$. Wir sagen, ein Objekt $x\in X$ steht bezüglich
$R$ zu einem Objekt $y\in Y$ in Beziehung, wenn die Aussage $R(x,y)$
eine wahre ist. Ein Beispiel wäre $R(x,y):=(x<y)$, wobei $X=Y$ die Menge
der ganzen Zahlen sei.

\begin{Definition}[Zweistellige Relation]\newlinefirst
Es seien $X,Y$ zwei Mengen. Jede Teilmenge $R\subseteq X\times Y$ heißt
Relation zwischen $X$ und $Y$. Zur Relation soll die Kenntnis von
$X,Y$ dazugehören. Insofern ist das durch das Tripel $(X,Y,R)$ kodierte
Objekt die Relation und $R$ ihr Graph.
\end{Definition}

\noindent
Eine Relation ist auch interpretierbar als wahrheitswertige Funktion
\[1_R\colon X\times Y\to\{0,1\},\quad 1_R(x,y):=[(x,y)\in R].\]
Sie ist die Indikatorfunktion der Teilmenge $R$ bezüglich der
Grundmenge $X\times Y$. Anstelle $(x,y)\in R$ oder $1_R(x,y)=1$
schreiben wir schlicht $R(x,y)$ für die Aussage, dass $x$ und $y$
bezüglich $R$ in Relation stehen.

Eine Relation mit $X\ne Y$ heißt \emph{heterogen}, eine mit
$X=Y$ heißt \emph{homogen}. Wir werden uns fast ausschließlich mit
homogenen Relationen beschäftigen. Man darf allerdings sagen, dass
auch die heterogenen in der Mathematik weit verbreitet sind, denn
jede Abbildung ist betrachtbar als Relation mit speziellen
Eigenschaften. Genauer ist eine Abbildung eine linkstotale und
rechtseindeutige Relation. Linkstotal heißt, dass jedes $x\in X$
zu mindestens einem $y\in Y$ in Beziehung steht. Rechtseindeutig heißt,
dass ein $y\in Y$ höchstens zu einem $x\in X$ in Beziehung steht. Es
handelt sich also lediglich um eine Sprechweise für die bereits
bekannten definierenden Eigenschaften.

Relationen lassen sich in Pfeildiagrammen darstellen. Man man zweichnet
hierbei einem Pfeil von einem $x\in X$ zu einem $y\in Y$, falls $x,y$
in Relation stehen.

\subsection{Äquivalenzrelationen}

Manchmal interessiert man sich nicht für die gänzliche Gleichheit
zweier Objekte. Die \emph{Äquivalenzrelationen}%
\index{Aequivalenzrelation@Äquivalenzrelation} verallgemeinern den
Gleichheitsbegriff dahingehend, dass zwei Objekte schon dann als
\emph{gleichartig}\index{gleichartig} angesehen werden, wenn sie in
in einer bestimmten Eigenschaft übereinstimmen. Um welche Eigenschaft
es sich dabei handelt, bestimmt die Relation.

\begin{Definition}[Äquivalenzrelation]\newlinefirst
Seien $A$ eine Menge und seien $x,y,z\in A$. Sei $R(x,y):=(x\sim y)$ eine
Relation. Man nennt $R$ Äquivalenzrelation, wenn gilt:
\[\begin{array}{ll}
x\sim x, &\text{(Reflexivität)}\\
x\sim y \implies y\sim x, & \text{(Symmetrie)}\\
x\sim y\land y\sim z\implies x\sim z.\quad & \text{(Transitivität)}
\end{array}\]
\end{Definition}

\begin{Definition}[Äquivalenzklasse]\newlinefirst
Sei $M$ eine Menge und $x\sim y$ eine Äquivalenzrelation für $x,y\in M$.
Die Menge%
\[[a] := \{x\in M\mid x\sim a\}\]
nennt man die Äquivalenzklasse zum Repräsentanten $a\in M$.
\end{Definition}

\begin{Satz}[Äquivalenzrelation induziert Zerlegung]\newlinefirst
Eine Menge wird durch eine Äquivalenzrelation in disjunkte
Äquivalenzklassen zerlegt, lat. partitioniert.
\end{Satz}
\strong{Beweis.} Sei $M$ die Menge und $x\sim y$ die Äquivalenzrelation.
Zu zeigen ist, dass kein Element von $M$ in mehr als einer
Äquivalenzklasse vorkommt. Seien $a,b,c\in M$, sei $c\in [a]$
und $c\in [b]$. Aufgrund von $c\sim a$ sowie $c\sim b$ und der
Transitivität gilt%
\[x\in [a]\iff x\sim a\iff x\sim c\iff x\sim b\iff x\in [b].\]
Man hat also
\[(\forall x\in M\colon x\in [a]\Leftrightarrow x\in [b])\iff [a]=[b].\]
Wenn also $[a]\ne [b]$ ist, kann nicht gleichzeitig $c\in [a]$ und $c\in [b]$
sein.\;\qedsymbol

\begin{Satz}[Zerlegung induziert Äquivalenzrelation]\newlinefirst
Sei $M$ eine Menge. Die Familie $(A_k)$ von Mengen $A_k\subseteq M$
bilde eine Zerlegung von $M$, d.\,h. dass die Vereinigung aller
$A_k$ die Menge $M$ überdeckt und dass paarweise $A_i\cap A_j=\emptyset$
für $i\ne j$ ist. Dann ist%
\[x\sim y\defiff \exists k\colon x\in A_k\land y\in A_k\]
eine Äquivalenzrelation auf $M$.
\end{Satz}
\strong{Beweis.} Da die $A_k$ die Menge $M$ überdecken,
muss es für ein beliebiges $x\in M$ mindestens eine Menge $A_k$
geben, so dass $x\in A_k$. Daher gilt $x\sim x$.

Die Symmetrie ergibt sich trivial.

Zur Transitivität. Voraussetzung ist $x\sim y$ und $y\sim z$.
Es gibt also ein $i$ mit $x\in A_i$ und $y\in A_i$. Außerdem gibt
es ein $j$ mit $y\in A_j$ und $z\in A_j$. Somit gilt%
\[\exists i\colon\exists j\colon x\in A_i\land y\in A_i\land y\in A_j\land z\in A_j.\]
Wegen
\[A_i\cap A_j = \emptyset \iff \forall y\colon (y\in A_i\land y\in A_j\Leftrightarrow 0)\]
für $i\ne j$ kann $y\in A_i\land y\in A_j$ aber nur erfüllt sein,
wenn $i=j$ ist. Daher ergibt sich%
\[\exists i\colon x\in A_i\land z\in A_i,\]
das heißt $x\sim z$.\;\qedsymbol

\begin{Definition}[Quotientenmenge]%
\index{Quotientenmenge}\index{Faktormenge}\newlinefirst
Für eine gegebene Äquivalenzrelation wird die aus allen
Äquivalenzklassen bestehende Menge
\[M/{\sim} := \{[x]\mid x\in M\}\]
als Quotientenmenge oder Faktormenge bezeichnet.
\end{Definition}

\begin{Definition}[Quotientenabbildung]%
\index{Quotientenabbildung}\newlinefirst
Für eine gegebene Äquivalenzrelation ist die Projektion
\[\pi\colon M\to M/{\sim},\quad \pi(x):=[x]\]
surjektiv und wird Quotientenabbildung genannt.
\end{Definition}

\begin{Definition}[Repräsentantensystem]%
\index{Repräsentantensystem}\newlinefirst
Für eine gegebene Äquivalenzrelation auf $M$ nennt man eine
Teilemenge $A\subseteq M$ ein vollständiges Repräsentantensystem,
wenn die Einschränkung $\pi|_A$ bijektiv ist, wobei mit $\pi$
die Quotientenabbildung gemeint ist.
\end{Definition}
Repräsentantensysteme ermöglichen die einfache Handhabung von
Äquivalenzklassen. Möchte man wissen, ob ein Element $x$ in der
Äquivalenklasse $[a]$ enthalten ist, dann braucht man bloß
zu überprüfen, ob $x\sim a$ ist. Außerdem besitzt die
Quotientenabbildung nun eine Darstellung $p\colon M\to A$,
dergestalt dass $\pi = \pi|_{A}\circ p$. Warum sollte das von
Bedeutung sein? Nun, Äquivalenzklassen fallen oft unendlich groß
aus. In der Kombinatorik treten zwar auch endliche Äquivalenzklassen auf,
diese werden trotzdem schnell unzugänglich groß. Die Äquivalenzklassen
und die Quotientenabbildung muss man also als abstrakte mathematische
Objekte betrachten. Abstrakte mathematische Objekte müssen wir erst
über eine Darstellung zugänglich machen, und genau dies ermöglicht
ein Repräsentantensystem.


\begin{Satz}[Charakterisierung von Äquivalenzklassen]\newlinefirst
Sei auf der Menge $M$ eine Äquivalenzrelation gegeben. Eine
Teilmenge $A\subseteq M$ ist genau dann eine Äquivalenzklasse,
wenn%
\begin{align*}
1.\;& A\ne\emptyset,\\
2.\;& x,y\in A\implies x\sim y,\\
3.\;& x\in A\land y\in M\land x\sim y\implies y\in A.
\end{align*}
\end{Satz}
\strong{Beweis.} Angenommen, $A$ ist eine Äquivalenzklasse.
Dann gibt es definitionsgemäß ein $a$ mit $A=[a]$. Daher ist
mindestens $a\in A$ und somit $A\ne\emptyset$. Mit $x,y\in A$ ergibt
sich $A=[x]=[y]$. Aufgrund von%
\[x\sim y \iff [a]=[b]\]
muss somit $x\sim y$ sein. Sei nun $x\in A$ und $y\in M$ mit
$x\sim y$. Es folgt $A=[x]=[y]$. Daher muss $y\in A$ sein.

Umgekehrt angenommen, die drei Eigenschaften sind erfüllt.
Zu zeigen ist, dass es ein $a$ gibt mit $A=[a]$. Da $A$ gemäß 1.
nichtleer ist, enthält es mindestens ein Element, dieses nennen wir
$a$. Für jedes weitere Element $x\in A$ ergibt sich
$x\sim a$, da sonst 2. verletzt sein würde. Schließlich muss man
noch wissen, ob $x\in A$, wenn $x\sim a$ und $x\in M$ ist.
Dies ist aber mit 3. gesichert. Es gibt also
tatsächlich ein $a$ mit $A=\{x\in M\mid x\sim a\}$.\;\qedsymbol

Eine große Fülle von Äquivalenzrelationen lässt sich auf die folgende
einfache Art konstruieren. Hat man eine beliebige Abbildung
$f\colon X\to Y$, dann sind die Urbilder $f^{-1}(\{y_1\})$ und
$f^{-1}(\{y_2\})$ disjunkt, sofern $y_1\ne y_2$, denn%
\[f^{-1}(\{y_1\})\cap f^{-1}(\{y_2\}) = f^{-1}(\{y_1\}\cap\{y_2\})
= f^{-1}(\emptyset) = \emptyset.\]
Demnach ist gemäß
\[Z = X/{\sim} = \{f^{-1}(\{y\})\mid y\in f(X)\}\]
eine Zerlegung des Definitionsbereichs $X$ gegeben und somit auch eine
Äquivalenzrelation. Für $x_1,x_2\in X$ gilt%
\[x_1\sim x_2 \iff f(x_1) = f(x_2).\]
Ist $f$ zudem surjektiv, dann gehört zu jedem Element von $Y$ genau eine
Äquivalenzklasse. Demnach definiert $f$ dann eine verallgemeinerte
Quotientenabbildung, da die Elemente von $Y$ die Äquivalenzklassen
charakterisieren. Die Bijektion $\varphi\colon Z\to Y$ hat dabei
die Eigenschaft $f = \varphi\circ\pi$. Sofern $Y$ für uns zugänglich
ist, resultiert hieraus auch eine verallgemeinerte Darstellung der
Quotientenabbildung, denn
\[f = \varphi\circ\pi = \varphi\circ (\pi|_A\circ p)
= (\varphi\circ\pi|_A)\circ p.\]
Nun ist $\varphi\circ\pi|_A$ auch bijektiv, weil $\varphi$
und $\pi|_A$ es sind. Somit charakterisiert $Y$ ein
vollständiges Repräsentantensystem.

Was bisher erläutert wurde, mag recht abstrakt erscheinen. Wir haben
aber eigentlich ein recht intuitives Verständnis für diese
Begrifflichkeiten. Ein Bilderbuchbeispiel für eine Quotientenmenge
bieten die Klassen einer Schule. Zwei Schüler seien genau dann
äquivalent, wenn sie in dieselbe Klasse gehen. Die Äquivalenzklasse eines
Schülers ist dann schlicht seine Schulklasse. Die Menge der Schüler der
Schule wird in disjunkte Schulklassen zerlegt. Die Menge dieser
Schulklassen bildet die Quotientenmenge. Ein vollständiges
Repräsentantensystem entsteht zum Beispiel durch die Wahl eines
Klassensprechers in jeder Klasse.

Ein weiteres typisches Beispiel für eine Äquivalenzrelation ist die
Kongruenz modulo $m$, die elementar in der Zahlentheorie
und Gruppentheorie vorkommt. Die Äquivalenzklassen sind hier die
Restklassen. Die Reste bilden ein kanonisches vollständiges
Repräsentantensystem. Das Bilden des Restes zu einer Zahl ist
eine Darstellung der Quotientenabbildung.

\subsection{Operationen auf Äquivalenzklassen}

Äquivalenzklassen werden später wichtig sein für die Formulierung von
Konstruktionen. Bei diesen Konstruktionen ist eine Abbildung
zwischen Quotientenmengen erforderlich. Weil die Äquivalenzklassen
dabei über Repräsentanten dargestellt sind, liegt es nahe, auch
die Abbildung über Repräsentanten zu definieren. Dies wirft die
Frage nach der \emph{Wohldefiniertheit}\index{wohldefiniert} auf.
Darunter versteht man, dass die Abbildung auch tatsächlich unabhängig
von den gewählten Repräsentanten ist. Was das genau bedeutet, wird im
folgenden Abschnitt erklärt.

Gegeben seien zwei Quotientenmengen $M/\sim$ und $M'/\sim'$.
Eine vorhandene Abbildung $f\colon M\to M'$ induziert dann
eventuell gemäß%
\[f\colon M/\sim\to M'/\sim',\quad f([a]):=[f(a)]\]
eine Abbildung zwischen den Quotientenmengen. Kommt es dabei nicht
zu einem Widerspruch, liegt also eine Abbildung vor, spricht man
von \emph{Wohldefiniertheit}. Hierfür darf der Funktionswert nicht
vom gewählten Repräsentant abhängen, d.\,h. die Bedingung%
\[\forall x\in [a]\colon f(x)\in [f(a)]\]
muss erfüllt sein. Anders formuliert:
\[x\sim a \implies f(x)\sim' f(a).\]
Für mehrstellige Abbildungen ist das Vorgehen analog. Eine
Abbildung $f\colon M^2\to M'$ induziert%
\[f\colon (M/\sim)^2\to M'/\sim',\quad f([a],[b]):=[f(a,b)],\]
sofern
\[x\sim a\land y\sim b\implies f(x,y)\sim' f(a,b).\]
Bei den Konstruktionen kommen in der Regel zweistellige Abbildungen
(mit $M=M'$) vor, weil die Verknüpfungen von Elementen der algebraischen
Strukturen zweistellig sind. Diese Verknüpfungen werden im nächsten
Abschnitt besprochen.

\subsection{Kongruenzrelationen}

\begin{Definition}[Kongruenzrelation]%
\index{Kongruenzrelation}\newlinefirst
Gegeben sei eine Menge $M$, auf der eine zweistellige
Verknüpfung $*\colon M^2\to M$ definiert ist. Eine Äquivalenzrelation
auf $M$ nennt man Kongruenzrelation, wenn die induzierte Verknüpfung
$[a]*[b]:=[a*b]$ wohldefiniert ist.
\end{Definition}
Bei einer Kongruenzrelation sagt man »$a$ ist kongruent zu $b$«
anstelle von »$a$ ist äquivalent zu $b$« und schreibt $a\equiv b$
anstelle von $a\sim b$. Eigentlich kann man den Begriff für eine
beliebige Stelligkeit definieren. Es besteht jedoch zunächst nur Bedarf
an zweistelligen Verknüpfungen.

Im Folgenden schreiben wir für die Verknüpfung kurz $ab$ anstelle
$a*b$. Das spart ein wenig Schreibaufwand und ist so üblich, solange
keine Verwechslungsgefahr mit einer bereits auf andere Art definierten
Multiplikation besteht.

\begin{Satz}\label{Kongruenz-Quotient}
Sei $M$ eine Struktur aus der Liste Magma, Monoid, Halbgruppe,
Gruppe, kommutatives Monoid, kommutative Gruppe. Sei $\equiv$ eine
Kongruenzrelation auf $M$ und $\varphi$ die zugehörige
Quotientenabbildung. Dann bildet die Quotienmenge $M/\equiv$
bezüglich der induzierten Verknüpfung $\varphi(a)\varphi(b):=\varphi(ab)$
ebenfalls eine Struktur derselben Art und $\varphi$ ist ein
Homomorphismus.
\end{Satz}
\strong{Beweis.} Im Folgenden seien $a',b',c'$ beliebige Elemente
der Quotientenmenge. Weil $\varphi$ surjektiv ist, gibt es immer
$a,b,c\in M$ mit $a'=\varphi(a)$, $b'=\varphi(b)$ und $c'=\varphi(c)$.

Die Verknüpfung auf $M$ sei abgeschlossen. Dann ist
\[a'b' = \varphi(a)\varphi(b) = \varphi(ab)\in M/\equiv.\]
Somit ist die Quotientenmenge bezüglich der induzierten Verknüpfung
abgeschlossen.

Die Verknüpfung auf $M$ erfülle das Assoziativgesetz. Dann gilt
\[(a'b')c' = \varphi(ab)\varphi(c) = \varphi(abc)
= \varphi(a)\varphi(bc) = a'(b'c').\]
Die induzierte Verknüpfung erfüllt somit ebenfalls das
Assoziativgesetz.

Die Verknüpfung auf $M$ habe ein neutrales Element $e$. Dann gilt
\[\varphi(a)=\varphi(ea) = \varphi(e)\varphi(a),\quad
\varphi(a)=\varphi(ae)=\varphi(a)\varphi(e).\]
Demzufolge besitzt $M/\equiv$ mit $e':=\varphi(e)$ ebenfalls
ein neutrales Element.

Zur Verknüpfung auf $M$ gebe es zu jedem Element ein inverses. Dann gilt
\[\varphi(e) = \varphi(aa^{-1}) = \varphi(a)\varphi(a^{-1}),\quad
\varphi(e) = \varphi(a^{-1}a) = \varphi(a^{-1})\varphi(a).\]
Demzufolge gibt es auf der Quotientenstruktur mit
$\varphi(a)^{-1}:=\varphi(a^{-1})$ ebenfalls zu jedem Element
ein inverses.

Die Verknüpfung auf $M$ sei kommutativ. Dann gilt
\[a'b' = \varphi(a)\varphi(b) = \varphi(ab) = \varphi(ba)
= \varphi(b)\varphi(a) = b'a'.\]
Somit ist die Verknüpfung auf der Quotientenstruktur $M/\equiv$ ebenfalls
kommutativ.\,\qedsymbol

\begin{Satz}\label{Kongruenz-Ring-Quotient}
Satz \ref{Kongruenz-Quotient} gilt auch für Ringe, unitäre Ringe,
kommutative Ringe und kommutative unitäre Ringe, sofern die
Relation eine Kongruenzrelation sowohl bezüglich der additiven
als auch der multiplikativen Verknüpfung ist.
\end{Satz}
\strong{Beweis.} Sei $(R,+,\cdot)$ der Ring und $\equiv$ die
Kongruenzrelation. Gemäß Satz \ref{Kongruenz-Quotient} ist
$(R/\equiv,+)$ eine kommutative Gruppe und $(R/\equiv,\cdot)$
eine Halbgruppe. Bei einem unitären Ring ist $(R/\equiv,\cdot)$
ein Monoid, und bei einem kommutativen unitären Ring ein
kommutatives Monoid.

Es verbleiben noch die Distributivgesetze zu prüfen.
Sei $\varphi$ die Quotientenabbildung. Man rechnet
\begin{align*}
a'(b'+c') &= \varphi(a)(\varphi(b)+\varphi(c)) = \varphi(a)(\varphi(b+c))
= \varphi(a(b+c)) = \varphi(ab+ac)\\
&= \varphi(ab)+\varphi(ac) = \varphi(a)\varphi(b)+\varphi(a)\varphi(c)
= a'b'+a'c'.
\end{align*}
Die Rechnung zum Rechtsdistributivgesetz ist analog.

Damit ist der Satz gezeigt, und ferner ist gezeigt dass $\varphi$ ein
Ringhomomorphismus ist. Und für einen unitären Ring ist $\varphi$
Eins-erhaltend, wie bereits aus Satz \ref{Kongruenz-Quotient}
hervorgeht.\,\qedsymbol

% Ordinalzahlen. Eine Dominoreihe verläuft sich in einer Schachtel
% ins Unendliche. Der Stein mit Limesordinal muss extra umfallen,
% zu ihm gibt es keinen Vorgänger.





\chapter{Algebra}
\section{Gruppentheorie}
\subsection{Grundbegriffe}
\begin{definition}[Gruppenhomomorphismus]
Sind $(G,*)$ und $(H,\bullet)$ zwei Gruppen, so
heißt $\varphi\colon G\to H$ \emdef{Gruppenhomomorphismus}%
\index{Gruppenhomomorphismus}, wenn
\begin{gather}
\forall g_1,g_2\in G\colon
  \varphi(g_1*g_2) = \varphi(g_1)\bullet\varphi(g_2)
\end{gather}
gilt. Ein \emdef{Gruppenisomorphismus}\index{Isomorphismus!zwischen Gruppen}
ist ein bijektiver Gruppenhomomorphismus, da die Umkehrabbildung
auch wieder ein Gruppenhomomorphismus ist.
\end{definition}
\begin{definition}[Direktes Produkt]
\emdef{Direktes Produkt}\index{direktes Produkt}:
\begin{gather}
G\times H := \{(g,h)\mid g\in G, h\in H\},\\
(g_1,h_1)*(g_2,h_2) := (g_1*g_2, h_1*h_2).
\end{gather}
\end{definition}
\noindent
\strong{Satz von Lagrange.} Für Gruppen $G,H$ gilt:
\begin{equation}
H\le G\implies |G| = |G/H|\cdot |H|.
\end{equation}

\subsection{Gruppenaktionen}\label{Gruppenaktion}
\begin{definition}[Gruppenaktion]
Eine Funktion $f\colon G\times X\to X$ heißt
\emdef{Gruppenaktion}\index{Gruppenaktion}, wenn
\begin{gather}
\hspace{-1em}\forall g_1,g_2{\in}G, x{\in}X\colon f(g_1,f(g_2,x)) = f(g_1 g_2,x),\\
\hspace{-1em}\forall x\in X\colon f(e,x) = x
\end{gather}
gilt, wobei mit $e$ das neutrale Element von $G$ gemeint ist.
Anstelle von $f(g,x)$ wird üblicherweise kurz $gx$ (oder
$g+x$ bei einer Gruppe $(G,+)$) geschrieben.

Anstelle von \emdef{Linksaktionen} kommen auch \emdef{Rechtsaktionen}
vor, die sich von Linksaktionen in der Reihenfolge unterscheiden.
Eine Rechtsaktion $f\colon X\times G\to X$ genügt den Regeln
\begin{gather}
\hspace{-1em}\forall g_1,g_2{\in}G, x{\in}X\colon f(f(x,g_1),g_2) = f(x,g_1 g_2),\\
\hspace{-1em}\forall x\in X\colon f(x,e) = x.
\end{gather}
\end{definition}

\begin{definition}[Orbit, Stabilisator]
Für ein $x\in X$ wird
\begin{equation}\label{eq:Orbit}
Gx := \{gx\mid g\in G\}
\end{equation}
\emdef{Bahn}\index{Bahn} oder
\emdef{Orbit}\index{Orbit!unter einer Gruppenaktion} genannt.
Die Menge
\begin{equation}
G_x := \{g\in G\mid gx=x\}
\end{equation}
wird \emdef{Fixgruppe}\index{Fixgruppe}
oder \emdef{Stabilisator}\index{Stabilisator} genannt.
Die Menge
\begin{equation}
X^g := \{x\in X\mid gx=x\}
\end{equation}
heißt \emdef{Fixpunktmenge}.
\end{definition}

\noindent
Fixgruppen sind immer Untergruppen:
\begin{equation}
\forall x\colon G_x\le G.
\end{equation}
Bahnen sind Äquivalenzklassen, die Quotientenmenge
\begin{equation}
X/G := \{Gx\mid x\in X\}
\end{equation}
wird \emdef{Bahnenraum}\index{Bahnenraum} genannt.

\strong{Bahnformel.}\index{Bahnformel}
Ist $G$ eine endliche Gruppe, so gilt:
\begin{equation}
|G| = |Gx|\cdot |G_x|.
\end{equation}
\strong{Lemma von Burnside.}\index{Lemma von Burnside}
Für eine endliche Gruppe $G$ gilt:%
\begin{equation}
|X/G| = \frac{1}{|G|}\sum_{g\in G}|X^g|.
\end{equation}

\section{Ringe}\index{Ring}
Ist $(R,+,*)$ ein Ring, so gilt für alle $a,b\in R$:
\begin{align}
0*a &= a*0 = 0,\\
(-a)*b &= a*(-b) = -(a*b),\\
(-a)*(-b) &= -(a*b).
\end{align}
\begin{definition}[Ringhomomorphismus]
Sind $(R,+,*)$ und $(R',+',*')$ Ringe, so wird
$\varphi\colon R\to R'$ als \emdef{Ringhomomorphismus}
bezeichnet, wenn
\begin{align}
\varphi(a+b) &= \varphi(a)+'\varphi(b),\\
\varphi(a*b) &= \varphi(a)*'\varphi(b),
\end{align}
für alle $a,b\in R$ gilt und $\varphi(1)=1$ ist.
\end{definition}

\subsection{Polynome}\index{Polynom}
\begin{definition}[Polynom, Polynomring, Koeffizienten]
Sei $R$ ein kommutativer unitärer Ring.
Mit $R[X]$ bezeichnen wir die Menge der unendlichen Folgen
\begin{equation}
(a_k) = (a_0,a_1,\ldots,a_n,0,0,0,\ldots)
\end{equation}
mit $a_k\in R$, bei denen ab einem Index alle Folgenglieder null sind.

Für zwei Folgen aus $R[X]$ wird nun die Addition
\begin{equation}
(a_k) + (b_k) := (a_k+b_k)
\end{equation}
und die Multiplikation
\begin{equation}\label{eq:Faltung}
(a_i)*(b_j) = \bigg(\sum_{i=0}^k a_i b_{k-i}\bigg)
\end{equation}
erklärt. In der Form \eqref{eq:Faltung} wird die Operation auch
\emdef{Faltung}\index{Faltung!von zwei Folgen}
der Folgen $(a_i)$ und $(b_j)$ genannt.

Die Menge $R[X]$ bildet mit der Addition und Multiplikation
einen kommutativen unitären Ring, den \emdef{Polynomring}
mit Koeffizienten in $R$. Ein Element von $R[X]$ wird
\emdef{Polynom} genannt.

Man definiert nun
\begin{equation}
X:=(0,1,0,0,0,\ldots),
\end{equation}
womit sich jedes Polynom in der Form
\begin{equation}\textstyle
(a_k) = \sum_{k=0}^n a_k X^k
\end{equation}
schreiben lässt. Die $a_k$ nennt man \emdef{Koeffizienten}
des Polynoms.
\end{definition}

\noindent
Die Addition bekommt nun die Form
\begin{equation}
\sum_{k=0}^m a_k X^k + \sum_{k=0}^n b_k X^k
:= \sum_{k=0}^p (a_k+b_k)X^k.
\end{equation}
mit $p=\max(m,n)$. Die Multiplikation lässt sich nun in der Form
\begin{equation}
\bigg(\sum_{i=0}^m a_i X^i\bigg)\bigg(\sum_{j=0}^n b_j X^j\bigg)
:= \sum_{k=0}^{m+n}\bigg(\sum_{i=0}^k a_i b_{k-i}\bigg) X^k.
\end{equation}
schreiben. Die Multiplikation von Polynomen ist das gewöhnlichen
Ausmultiplizieren der Polynome, wobei $X^i X^j=X^{i+j}$ gilt.

Die $X^k$ können als Vektorraumbasis betrachtet
werden und dienen dabei dazu, die $a_k$ auseinanderzuhalten.
Zwei Polynome $\sum_{k=0}^m a_k X^k$ und $\sum_{k=0}^n b_k X^k$
sind genau dann gleich, wenn $a_k=b_k$ für alle $k\le\max(m,n)$ gilt.

Da $R[X]$ wieder ein kommutativer unitärer Ring ist,
ist auch $R[X][Y]$ ein Polynomring. Man definiert
\begin{equation}
R[X,Y] := R[X][Y].
\end{equation}
Polynome aus $R[X,Y]$ lassen sich in der Form
\begin{equation}
\sum_{j=0}^n \bigg(\sum_{i=0}^m a_{ij}X^i\bigg)Y^j
= \sum_{i=0}^m\sum_{j=0}^n a_{ij} X^i Y^j
\end{equation}
mit $a_{ij}\in R$ schreiben.

Allgemein ist die Menge
\begin{equation}
R[X_1,\ldots,X_q] := X[X_1,\ldots,X_{q-1}][X_q]
\end{equation}
ein kommutativer unitärer Ring. Die Polynome lassen sich in der Form
\begin{equation}
\sum_{k\in\N_0^q} a_k X^k\quad (a_k\in R)
\end{equation}
mit
\[k=(k_1,\ldots,k_q)\quad\text{und}\quad X^k:=\prod_{i=1}^q X_i^{k_i}\]
schreiben.

\begin{definition}[Grad]
Für ein Polynom $f=\sum_{k=0}^n a_k X^k$ mit $a_n\ne 0$ wird
$n$ als \emdef{Grad} von $f$ bezeichnet. Man schreibt $n=\deg f$.

Für ein Monom $a_{ij} X^i Y^j$ mit $a_{ij}\ne 0$ heißt $i+j$
\emdef{Totalgrad}. Der \emdef {Grad} eines Polynoms
\begin{equation}
\textstyle\sum_{i=1}^m\sum_{j=1}^n a_{ij} X^i Y^j
\end{equation}
ist der maximale Totalgrad aller Monome mit $a_{ij}\ne 0$.
Für Polynome in mehr als zwei Variablen ist die Definition analog.
\end{definition}

\strong{Regeln.}\\
Für zwei Polynome $f,g\in R[X_1,\ldots,X_q]$ gilt:
\begin{align}
\deg(f+g)&\le \max(\deg f,\deg g),\\
\deg(fg)&\le (\deg f)(\deg g).
\end{align}
Für zwei Polynome $f,g$ mit $\deg f\ne\deg g$ gilt:
\begin{equation}
\deg(f+g) = \max(\deg f,\deg g).
\end{equation}
Ist $R$ ein Integritätsring, so gilt für $f,g\in R[X_1,\ldots,X_q]$:%
\begin{equation}
\deg(fg) = (\deg f)(\deg g).
\end{equation}
Jeder Körper, z.\,B. $\R$ oder $\C$ ist ein Integritätsring.
Auch die ganzen Zahlen $\Z$ bilden einen Integritätsring.
Ein Polynomring ist genau dann ein Integritätsring, wenn die
Koeffizienten aus einem Integritätsring entstammen.

\begin{definition}[Einsetzungshomomorphismus]
Seien $R,R'$ kommutative unitäre Ringe. Sei $R'$ eine Ringerweiterung
von $R$ und sei $r\in R'$. Die Abbildung $\varphi_r\colon R[X]\to R'$
mit
\begin{equation}\textstyle
\varphi_r(p) = p(r) := \sum_{k=0}^n a_k r^k 
\end{equation}
für jedes Polynom
\[\textstyle p = \sum_{k=0}^n a_k X^k\]
ist ein Ringhomomorphismus. Man bezeichnet $p(r)$ als \emdef{Einsetzung}
von $r$ in $p$ und $\varphi_r$ als
\emdef{Einsetzungshomomorphismus}\index{Einsetzungshomomorphimus}.
\end{definition}

Man kann auch $R'=R$ und $r=X$ setzen, dann gilt $p=p(X)$.
Ein Polynom stimmt also mit der Einsetzung seiner eigenen formalen
Variablen überein.

\begin{definition}[Polynomfunktion]
Für ein festes $p\in R[X]$ wird die Funktion
\begin{equation}
f\colon R'\to R',\quad f(x):=p(x)
\end{equation}
als \emdef{Polynomfunktion} bezeichnet.
\end{definition}

\noindent
In einigen Ringen können unterschiedliche Polynome zur selben
Polynomfunktion führen. Handelt es sich bei $R$ jedoch um einen
unendlichen Körper, z.\,B. $R=\R$ oder $R=\C$, dann gibt es zu jeder
Polynomfunktion nur ein einziges Polynom.

\section{Körper}
\begin{definition}[Körperhomomorphismus]
Sind $(K,+,\bullet)$ und $(K',+',\bullet')$ Körper, so
wird $\varphi\colon K\to K'$ als \emph{Körperhomomorphismus}
bezeichnet, wenn
\begin{align}
\varphi(a+b) &= \varphi(a)+'\varphi(b),\\
\varphi(a\bullet b) &= \varphi(a)\bullet'\varphi(b)
\end{align}
für alle $a,b\in K$ gilt und $\varphi(1)=1$ ist.
\end{definition}




\chapter{Zahlenbereiche}

\section{Ganze Zahlen}

\subsection{Konstruktion}

\begin{Definition}[Ganze Zahlen]\mbox{}\\*
Auf $\N_0\times\N_0$ wird die folgende Äquivalenzrelation definiert:
\[(x_1,y_1)\sim (x_2,y_2) \defiff x_1+y_2 = x_2+y_1.\]
Die Quotientenmenge $\Q := (\N_0\times\N_0)/{\sim}$ nennt man
die ganzen Zahlen.
\end{Definition}

\begin{Satz}[Ring der ganzen Zahlen]\mbox{}\\*
Die Operationen
\begin{align*}
[(x_1,y_1)]+[(x_2,y_2)] &:= [(x_1+x_2,y_1+y_2)],\\
[(x_1,y_1)]\cdot [(x_2,y_2)] &:= [(x_1x_2+y_1y_2,x_1y_2+x_2y_1)]
\end{align*}
sind auf $\Q$ wohldefiniert und $(\Q,+,\cdot)$ bildet einen
kommutativen unitären Ring.
\end{Satz}
\strong{Beweis.} Wohldefiniert heißt, dass das Ergebnis der Operationen
nicht von den gewählten Repräsentanten der Argumente abhängig ist.
Sei dazu $(x_1,y_1)\sim(a_1,b_1)$ und
$(x_2,y_2)\sim (a_2,b_2)$. Zu zeigen ist nun
\begin{gather*}
(x_1+x_2,y_1+y_2)\sim (a_1+a_2,b_1+b_2)\\
\iff (x_1+x_2)+(b_1+b_2) = (a_1+a_2)+(y_1+y_2).
\end{gather*}
Gemäß Voraussetzung ist $x_1+b_1=a_1+y_1$ und $x_2+b_2=a_2+y_2$.
Man bekommt damit auf der linken Seite
\[x_1+x_2+b_1+b_2 = a_1+y_1+a_2+y_2,\]
was wiederum mit der rechten Seite übereinstimmt.

Mit der Multiplikation verhält es sich etwas komplizierter.
Zu Vereinfachung wird zunächst gezeigt:
\begin{gather*}
[(x_1,y_1)]\cdot [(x_2,y_2)] = [(a_1,b_1)]\cdot [(x_2,y_2)]\\
\iff (x_1x_2+y_1y_2,x_1y_2+y_1x_2)\sim (a_1x_2+b_1y_2,a_1y_2+b_1x_2)\\
\iff x_1x_2+y_1y_2 + a_1y_2+b_1x_2 = a_1x_2+b_1y_2 + x_1y_2+y_1x_2\\
\iff (x_1+b_1)x_2 + (a_1+y_1)y_2 = (a_1+y_1)x_2 + (x_1+b_1)y_2.
\end{gather*}
Diese Gleichung ist gemäß Voraussetzung $(x_1,y_1)\sim (a_1,b_1)$
bzw. $x_1+b_1=a_1+y_1$ erfüllt.

Analog bestätigt man
\[[(a_1,b_1)]\cdot [(x_2,y_2)] = [(a_1,b_1)]\cdot [(a_2,b_2)].\]
Gemäß Transitivität ergibt sich somit
\[[(x_1,y_1)]\cdot [(x_2,y_2)] = [(a_1,b_1)]\cdot [(a_2,b_2)].\]
Es ist nun zu bestätigen, dass $(\Q,+)$ eine kommutative Gruppe ist.
Das Assoziativgesetz:
\begin{gather*}
([(x_1,y_1)]+[(x_2,y_2)])+[(x_3,y_3)]
= [(x_1+x_2,y_1+y_2)] + [(x_3,y_3)]\\
= [(x_1+x_2+x_3,y_1+y_2+y_3)]
= [(x_1,y_1)]+[(x_2+x_3,y_2+y_3)]\\
= [(x_1,y_1)]+([(x_2,y_2)]+[(x_3,y_3)]).
\end{gather*}
Das neutrale Element ist $[(0,0)]$:
\[[(x,y)]+[(0,0)] = [(x+0,y+0)] = [(x,y)].\]
Das inverse Element zu $[(x,y)]$ ist $[(y,x)]$, denn es gilt
\begin{gather*}
[(x,y)]+[(y,x)] = [(x+y,y+x)] = [(0,0)]\\
\iff (x+y,y+x)\sim (0,0)\iff x+y+0 = y+x+0.
\end{gather*}
Das Kommutativgesetz:
\begin{gather*}
[(x_1,y_1)]+[(x_2,y_2)] = [(x_1+x_2,y_1+y_2)]
= [(x_2+x_1,y_2+y_1)]\\
= [(x_2,y_2)]+[(x_1,y_1)].
\end{gather*}
Es ist nun zu bestätigen, dass $(\Q,\cdot)$ ein kommutatives
Monoid bildet. Das Assoziativgesetz:
\begin{gather*}
([(x_1,y_1)]\cdot [(x_2,y_2)])\cdot [(x_3,y_3)]
= [(x_1x_2+y_1y_2,x_1y_2+x_2y_1)]\cdot [(x_3,y_3)]\\
= [(x_1x_2x_3+x_3y_1y_2+x_1y_2y_3+x_2y_1y_3,\;
x_1x_2y_3+y_1y_2y_3+x_1x_3y_2+x_2x_3y_1)]\\
= [(x_1,y_1)]\cdot [(x_2x_3+y_2y_3,x_2y_3+x_3y_2)]
= [(x_1,y_1)]\cdot ([(x_2,y_2)]\cdot [(x_3,y_3)]).
\end{gather*}
Das Kommutativgesetz:
\begin{gather*}
[(x_1,y_1)]\cdot [(x_2,y_2)] = [(x_1x_2+y_1y_2,\;x_1y_2+y_1x_2)]\\
= [(x_2x_1+y_2y_1,\;x_2y_1+x_1y_2)] = [(x_2,y_2)]\cdot [(x_1,y_1)].
\end{gather*}
Das neutrale Element ist $[(1,0)]$, denn es gilt
\[[(x,y)]\cdot [(1,0)] = [(x\cdot 1+y\cdot 0,\;1\cdot y+x\cdot 0)]
= [(x,y)].\]
Schließlich ist noch das Distributivgesetz zu bestätigen.
Man findet
\begin{gather*}
[(a,b)]\cdot ([(x_1,y_1)]+[(x_2,y_2)])
= [(a,b)]\cdot [(x_1+x_2,y_1+y_2)]\\
= [(ax_1+ax_2+by_1+by_2,\;ay_1+ay_2+bx_1+bx_2)]\\
= [(ax_1+by_1,ay_1+bx_1)]+[(ax_2+by_2,ay_2+bx_2)]\\
= [(a,b)]\cdot [(x_1,y_1)] + [(a,b)]\cdot [(x_2,y_2)].
\end{gather*}
Somit sind alle Axiome bestätigt.\;\qedsymbol

\begin{Definition}[Monoidhomomorphismus]\mbox{}\\*
Seien $(M,+)$ und $(M',+')$ zwei Monoide. Eine Abbildung
$\varphi\colon M\to M'$ heißt Monoidhomomorphismus, wenn
für alle $a,b\in M$ gilt
\[\varphi(a+b) = \varphi(a)+\varphi(b)\]
und $\varphi(0)=0'$ ist.
\end{Definition}
Einen injektiven Homomorphismus nennt man Monomorphismus. Ein
Monomorphismus charakterisiert eine Einbettung einer Struktur als
Unterstruktur einer anderen.

\begin{Satz}[Einbettung der natürlichen Zahlen in die ganzen]%
\label{embedding-nat-int}\mbox{}\\*
Die Abbildung $\varphi\colon\N_0\to\Z$ mit $\varphi(n):=[(n,0)]$
ist ein Monoidmonomorphismus.
\end{Satz}
\strong{Beweis.} Es ergibt sich
\[\varphi(a+b) = [(a+b,0)] = [(a,0)]+[(b,0)] = \varphi(a)+\varphi(b).\]
Außerdem ist $\varphi(0)=[(0,0)]$, und $[(0,0)]$
ist das neutrale Element von $(\Z,+)$.

Schließlich ist noch die Injektivität zu prüfen:
\begin{gather*}
[(a,0)] = \varphi(a) = \varphi(b)  = [(b,0)]\iff (a,0)\sim (b,0)\\
\iff a+0 = b+0 \iff a=b.\;\qedsymbol
\end{gather*}
Anstelle von $\varphi(n)=[(n,0)]$ darf man daher einfach schreiben
$n=[(n,0)]$. Außerdem definiert man $a-b:=a+(-b)$. Daraus
ergibt sich nun
\[[(x,y)] = [(x,0)]+[(0,y)] = [(x,0)] - [(y,0)] = x-y.\]
Die umständliche Schreibweise $[(x,y)]$ wird ab jetzt nicht
mehr benötigt.

\begin{Definition}[Totalordnung der ganzen Zahlen]\mbox{}\\*
Auf $\Z$ wird die folgende strenge Totalordnung definiert:
\[[(x_1,y_1)] < [(x_2,y_2)] \defiff x_1+y_2<x_2+y_1.\]
\end{Definition}

\begin{Satz}[Einbettung der Totalordnung]\mbox{}\\*
Die Abbildung $\varphi$ aus Satz \ref{embedding-nat-int}
genügt der Forderung
\[a<b \implies \varphi(a)<\varphi(b).\]
\end{Satz}
\strong{Beweis.} Nach den Definitionen ist
\[\varphi(a)<\varphi(b)\iff [(a,0)]<[(b,0)]\iff a+0<0+b\iff a<b.\;\qedsymbol\]


\newpage
\section{Rationale Zahlen}

\subsection{Konstruktion}

\begin{Definition}[Rationale Zahlen]\mbox{}\\*
Auf $\Z\times\N_1$ wird die folgende Äquivalenzrelation definiert:
\[(x_1,y_1)\sim (x_2,y_2) \defiff x_1y_2 = x_2y_1.\]
Die Quotientenmenge $\Q := (\Z\times\N_1)/{\sim}$ nennt man
die rationalen Zahlen.
\end{Definition}
Für die Äquivalenzklasse $[(x,y)]$ schreibt man $\frac{x}{y}$.

\begin{Satz}[Körper der rationalen Zahlen]\mbox{}\\*
Die Operationen
\[\frac{x_1}{y_1}+\frac{x_2}{y_2} := \frac{x_1y_2+x_2y_1}{y_1y_2},
\qquad\frac{x_1}{y_1}\cdot \frac{x_2}{y_2} := \frac{x_1x_2}{y_1y_2}\]
sind auf $\Q$ wohldefiniert und $(\Q,+,\cdot)$ bildet einen Körper.
\end{Satz}
\strong{Beweis.} Wohldefiniert bedeutet, dass das Ergebnis der
Operationen nicht von den gewählten Repräsentanten der Argumente
abhängig ist. Sei dazu $(a_1,b_1)\sim (x_1,y_1)$ und
$(a_2,b_2)\sim (x_2,y_2)$. Zu zeigen ist nun
\begin{align*}
&(a_1b_2+a_2b_1,b_1b_2)\sim (x_1y_2+x_2y_1,y_1y_2)\\
&\iff (a_1b_2+a_2b_1)(y_1y_2) = (x_1y_2+x_2y_1)(b_1b_2)\\
&\iff a_1b_2 y_1y_2 + a_2b_1y_1y_2 = x_1y_2b_1b_2+x_2y_1b_1b_2.
\end{align*}
Substituiert man $a_1y_1=x_1b_1$ und $a_2y_2=x_2b_2$ auf
der linken Seite der Gleichung, ergibt sich die rechte Seite.
Zu zeigen ist weiterhin
\begin{align*}
(a_1a_2,b_1b_2)\sim (x_1x_2,y_1y_2)
\iff a_1a_2y_1y_2 = x_1x_2b_1b_2.
\end{align*}
Wieder wird linke Seite der Gleichung über $a_1y_1=x_1b_1$
und $a_2y_2=x_2b_2$ in die rechte Seite überführt.
Die Wohldefiniertheit der Operationen ist damit bestätigt.

Es bleibt zu prüfen, dass $(\Q,+,\cdot)$ allen Körperaxiomen genügt.
Das neutrale Element der Addition ist $0/1$, denn es gilt
\[\frac{x}{y}+\frac{0}{1} = \frac{x\cdot 1+0\cdot y}{y\cdot 1} = \frac{x}{y}.\]
Das neutrale Element der Multiplikation ist $1/1$, denn es gilt
\[\frac{x}{y}\cdot\frac{1}{1} = \frac{x\cdot 1}{y\cdot 1} = \frac{x}{y}.\]
Die Assoziativität der Addition ergibt sich ohne größere Umstände:
\begin{align*}
\bigg(\frac{x_1}{y_1}+\frac{x_2}{y_2}\bigg)+\frac{x_3}{y_3}
&= \frac{x_1y_2+x_2y_1}{y_1y_2} + \frac{x_3}{y_3}
= \frac{x_1y_2y_3+x_2y_1y_3+x_3y_1y_2}{y_1y_2y_3},\\
\frac{x_1}{y_1}+\bigg(\frac{x_2}{y_2}+\frac{x_3}{y_3}\bigg)
&= \frac{x_1}{y_1}+\frac{x_2y_3+x_3y_2}{y_2y_3}
= \frac{x_1y_2y_3+x_2y_1y_3+x_3y_1y_2}{y_1y_2y_3}.
\end{align*}
Die Assoziativität der Multiplikation ist etwas einfacher:
\[\bigg(\frac{x_1}{y_1}\cdot\frac{x_2}{y_2}\bigg)\cdot\frac{x_3}{y_3}
= \frac{x_1x_2}{y_1y_2}\cdot\frac{x_3}{y_3} = \frac{x_1x_2x_3}{y_1y_2y_3}
= \frac{x_1}{y_1}\cdot\frac{x_2x_3}{y_2y_3}
= \frac{x_1}{y_1}\cdot\bigg(\frac{x_2}{y_2}\cdot\frac{x_3}{y_3}\bigg).\]
Das Kommutativgesetz der Addition:
\[\frac{x_1}{y_1}+\frac{x_2}{y_2} = \frac{x_1y_2+x_2y_1}{y_1y_2}
= \frac{x_2y_1+x_1y_2}{y_2y_1}
= \frac{x_2}{y_2}+\frac{x_1}{y_1}.\]
Das Kommutativgesetz der Multiplikation:
\[\frac{x_1}{y_2}\cdot\frac{x_2}{y_2}
= \frac{x_1x_2}{y_1y_2} = \frac{x_2x_1}{y_2y_1}
= \frac{x_2}{y_2}\cdot\frac{x_1}{y_1}.\]
Das additiv inverse Element zu $x/y$ ist $(-x)/y$, denn es gilt
\[\frac{x}{y}+\frac{-x}{y} = \frac{xy+(-x)y}{y^2}
= \frac{0}{y^2} = \frac{0}{1}.\]
Das multiplikativ inverse Element zu $x/y$ mit $x\ne 0$
ist $y/x$, denn es gilt
\[\frac{x}{y}\cdot\frac{y}{x} = \frac{xy}{xy} = \frac{1}{1}.\]
Schließlich findet bestätigt man noch das Distributivgesetz:
\begin{align*}
&\frac{a}{b}\cdot\bigg(\frac{x_1}{y_1}+\frac{x_2}{y_2}\bigg)
= \frac{a}{b}\cdot\frac{x_1y_2+x_2y_1}{y_1y_2}
= \frac{ax_1y_2+ax_2y_1}{by_1y_2},\\
&\frac{ax_1}{by_1}+\frac{ax_2}{by_2}
= \frac{ax_1by_2+ax_2by_1}{by_1by_2}
= \frac{b}{b}\cdot\frac{ax_1y_2+ax_2y_1}{by_1y_2}.
\end{align*}
Hierbei beachtet man, dass $b/b=1/1$ das multiplikativ
neutrale Element ist.\;\qedsymbol

\begin{Definition}[Ringhomomorphismus]\mbox{}\\*
Seien $(R,+,*)$ und $(R',+',*')$ zwei Ringe. Die Abbildung
$\varphi\colon R\to R'$ heißt Ringhomomorphismus, wenn für alle
$a,b\in R$ gilt:
\begin{align*}
\varphi(a+b) = \varphi(a)+'\varphi(b),\qquad
\varphi(a*b) = \varphi(a)*'\varphi(b).
\end{align*}
Besitzt $R$ ein Einselement $1$ und $R'$ ein Einselement $1'$,
dann nennt man $\varphi$ Eins"=erhaltend, wenn $\varphi(1)=1'$ ist.
\end{Definition}
Einen injektiver Homomorphismus wird Monomorphismus genannt. Ein
Monomorphismus charakterisiert eine Einbettung einer Unterstruktur
in eine andere Struktur.

\newpage
\begin{Satz}[Einbettung der ganzen Zahlen in die rationalen]\mbox{}\\*
Sei $\varphi\colon\Z\to\Q$ mit $\varphi(z):=z/1$. Die
Abbildung $\varphi$ ist Eins"=erhaltender Ringmonomorphismus.
\end{Satz}
\strong{Beweis.} Die Erhaltung des Einselements ergibt sich
trivial. Ferner findet man
\[\varphi(a+b) = \frac{a+b}{1} = \frac{a\cdot 1+b\cdot 1}{1\cdot 1}
= \frac{a}{1}+\frac{b}{1} = \varphi(a)+\varphi(b)\]
und
\[\varphi(ab) = \frac{ab}{1} = \frac{ab}{1\cdot 1} = \frac{a}{1}\cdot\frac{b}{1}
= \varphi(a)\cdot\varphi(b).\;\qedsymbol\]
Gemäß der Einbettung können wir die ganze Zahl $z$ ab jetzt
mit der rationalen Zahl $z/1$ identifizieren. D.\,h. man schreibt
einfach $z=z/1$ anstelle von $\varphi(z)=z/1$.

\begin{Definition}[Division rationaler Zahlen]\mbox{}\\*
Wie in jedem Körper ist die Division für $a,b\in\Q$
definiert als $a/b := ab^{-1}$.
\end{Definition}
Die Division ist also gerade die Multiplikation des Kehrwertes
des Nenners:
\[\frac{x_1}{y_1}/\frac{x_2}{y_2} = \frac{x_1}{y_1}\cdot\frac{y_2}{x_2}.\]
Die Division muss natürlich mit der Notation für rationale Zahlen
kompatibel sein, sonst dürfte man nicht die gleiche Schreibweise
verwenden. Zur Unterscheidung schreiben wir Division für einen
Augenblick mit Doppelstrich als $a\doubleslash b$. Man findet
\[\frac{x}{y} = \frac{x}{1}\cdot\frac{1}{y}
= \frac{x}{1}\doubleslash\frac{y}{1} = x\doubleslash y.\]
Tatsächlich führt beides zum gleichen Ergebnis.

Da die rationalen Zahlen einen Körper bilden, gilt $a/a=aa^{-1}=1$
für jede rationale Zahl $a$.

\begin{Satz}[Addition, Subtraktion, Multiplikation von Brüchen]\mbox{}\\*
Seien $a,b,c,d$ rationale Zahlen mit $b\ne 0$ und $d\ne 0$. Es gilt
\[\frac{a}{b}+\frac{c}{d} = \frac{ad+bc}{bd},
\qquad \frac{a}{b}-\frac{c}{d} = \frac{ad-bc}{bd},
\qquad \frac{a}{b}\cdot\frac{c}{d} = \frac{ac}{bd}.\]
\end{Satz}
Der Beweis wird dem Leser überlassen.


\chapter{Kategorientheorie}

\section{Grundbegriffe}

\begin{Definition}[Kategorie]
Eine Kategorie ist ein Tripel $\category C = (\Ob,\Hom,{\circ})$,
sofern die folgenden beiden Axiome erfüllt sind:
\begin{enumerate}
\item Für $f\colon A\to B$, $g\colon B\to C$, $h\colon C\to D$ gilt
das Assoziativgesetz\\
$h\circ (g\circ f) = (h\circ g)\circ f$.
\item Für jedes Objekt $X$ existiert die Identität $\id_X\colon X\to X$,
so dass $f\circ\id_A = \id_B\circ f=f$ für alle Objekte $A,B$
und $f\colon A\to B$.
\end{enumerate}
\end{Definition}
Die Elemente der Klasse $\Ob$ nennt man Objekte. Die
Elemente der Klasse $\Hom$ nennt man Morphismen. Die
Verknüpfung $g\circ f$, sprich $g$ nach $f$, nennt man Verkettung
von $g$ und $f$.

Die Schreibweise ist $f\colon X\to Y$ gleichbedeutend mit
$f\in\Hom(X,Y)$, wobei $X,Y\in\Ob$.
Mit $\Hom(X,Y)$ ist die Teilklasse von
$\Hom$ gemeint, die alle Morphismen von $X$ nach $Y$
enthält. Man schreibt $\dom(f) = X$ und $\cod(f) = Y$.

Nun gut, man macht hier zunächst zwei Beobachtungen. Erstens
erinnern die Axiome an die Monoid"=Axiome, haben aber den Unterschied,
dass die Morphismen kompatibel sein müssen. D.\,h. um $g\circ f$
bilden zu können, muss $\cod(f)=\dom(g)$ sein.

Zweitens erinnern die Axiome an die Regeln für die Verkettung
von Abbildungen. Tatsächlich bilden die Abbildungen eine Kategorie.

\begin{Satz}[Kategorie der Mengen]\mbox{}\\*
Sei $\Omega$ das Mengenuniversum und für $A,B\in\Omega$ sei
$\Hom(A,B):=\Abb(A,B)$. Sei $g\circ f$ die Verkettung
von Abbildungen. Dann bildet $\strong{Set}:=(\Omega,\Hom,\circ)$
eine Kategorie.
\end{Satz}
\strong{Beweis.} Trivial.\;\qedsymbol

\begin{Satz}[Kategorie der Gruppen]\mbox{}\\*
Sei $\Omega$ die Klasse aller Gruppen und für $G,H\in\Omega$ sei
$\Hom(G,H)$ die Klasse der Homomorphismen von $G$ nach $H$.
Sei $g\circ f$ die Verkettung von Homomorphismen.
Dann bildet $\strong{Group}:=(\Omega,\Hom,\circ)$
eine Kategorie.
\end{Satz}
\strong{Beweis.} Homomorphismen sind Abbildungen, die Axiome
daher wie bei der Kategorie der Mengen erfüllt. Die Verkettung
zweier Homomorphismen ist ja auch ein Homomorphismus.\;\qedsymbol

Entsprechend bilden Ringe mit Ringhomomorphismen, Körper mit
Körperhomomorphismen, Vektorräume mit Vektorraumhomomorphismen
usw. Kategorien. Des Weiteren bilden die endlichen Mengen, Gruppen,
Ringe jeweils eine Kategorie.

Nun ist es so, dass Gruppen auch Mengen und Homomorphismen
auch Abbildungen sind. Die Kategorie der Gruppen ist gewissermaßen
in der Kategorie der Mengen enthalten. Um das zu präzisieren,
benötigen wir den Begriff des Vergissfunktors.

\begin{Definition}[Kovarianter Funktor]\mbox{}\\*
Sind $\category C,\category D$ Kategorien, dann nennt man
$F\colon\category C\to\category D$ einen
kovarianten Funktor, wenn jedem Objekt $X\in\Ob(\category C)$ ein Objekt
$F(X)\in\Ob(\category D)$ zugeordnet wird und jedem Morphismus
$f\in\Hom_{\category C}(X,Y)$ ein ein Morphismus
$F(f)\in\Hom_{\category D}(F(X),F(Y))$ zugeordnet wird,
so dass die folgenden beiden Verträglichkeitsaxiome erfüllt sind:%
\begin{gather*}
F(g\circ f) = F(g)\circ F(f),\\
F(\id_X) = \id_{F(X)}.
\end{gather*}
\end{Definition}
\begin{Definition}[Kontravarianter Funktor]\mbox{}\\*
Wie beim kovarianten Funktor, mit dem Unterschied
$F(g\circ f) = F(f)\circ F(g)$.
\end{Definition}
Bemerkung: Die Notation ist überladen. Nämlich ist die Zuordnung
$F\colon\Ob(\category C)\to\Ob(\category D)$ zu unterscheiden
von
\[\tilde F\colon\Hom_{\category C}(X,Y)\to\Hom_{\category D}(F(X),F(Y)).\]
Das Paar $(F,\tilde F)$ kodiert dann eigentlich den Funktor
$C\to D$.

\begin{Satz}[Vergissfunktor]\mbox{}\\*
Sei $F\colon\strong{Group}\to\strong{Set}$ mit $F((G,*,e)):=G$,
und jedem Gruppenhomomorphismus%
\[\varphi\colon (G,*,e)\to (G',*',e')\]
sei die Abbildung $F(\varphi)\colon G\to G'$ mit
$F(\varphi)(x):=\varphi(x)$ zugeordnet. Bei $F$ handelt
es sich um einen kovarianten Funktor.
\end{Satz}
\strong{Beweis.}
Es gilt $F(\id)(x) = \id(x)$, und daher $F(\id)=\id$.
Außerdem gilt%
\[
F(\varphi_2\circ\varphi_1)(x) = (\varphi_2\circ\varphi_1)(x)
= \varphi_2(\varphi_1(x))
= F(\varphi_2)(F(\varphi_1)(x))
= (F(\varphi_2)\circ F(\varphi_1))(x),
\]
und daher $F(\varphi_2\circ\varphi_1)
= F(\varphi_2)\circ F(\varphi_1)$.\;\qedsymbol

\begin{Satz} Sei $P(X)=2^X$ die Potenzmenge von $X$. Dann ist
wie folgt ein kovarianter Funktor gegeben:
\[P\colon\strong{Set}\to\strong{Set},\quad
P(X):=2^X,\quad P(f)(M):=f(M),\]
wobei $f$ eine beliebige Abbildung
und $f(M)$ die Bildmenge von $M$ unter $f$ ist.
\end{Satz}
\strong{Beweis.} Nach Satz \ref{Bild-unter-Komposition} gilt
\[P(g\circ f)(M) = (g\circ f)(M) = g(f(M))
= P(g)(P(f)(M)) = (P(g)\circ P(f))(M).\]
Daher ist $P(g\circ f)=P(g)\circ P(f)$. Außerdem ist
\[P(\id_X)(M) = \id_X(M) = M = \id_{P(X)}(M)\]
und daher $P(\id_X)=\id_{P(X)}$.\;\qedsymbol

Zum Funktor $P$ kommt noch ein weiterer Aspekt hinzu.
Für eine Abbildung $f$ kann man ganz pedantisch das
Bild $f(x)$ von der Bildmenge $f(\{x\})$ unterscheiden.
Aufgrund der Gleichung $f(\{x\})=\{f(x)\}$ verschwimmt diese
Unterscheidung aber gewissermaßen.
Die Abbildungen $f$ und $P(f)$ verhalten sich also
gewissermaßen gleich. Man kann sagen, dass $f$
auf ganz natürliche Art und Weise die Abbildung $P(f)$
zugeordnet ist. Definiert man
\[\eta(X)\colon X\to 2^X,\quad \eta(X)(x):=\{x\},\]
dann kommutiert das folgende Diagramm:
\[\xymatrix{
X \ar[r]^f \ar[d]_{\eta(X)} & Y \ar[d]^{\eta(Y)} \\
2^X \ar[r]_{P(f)} & 2^Y }\]
D.\,h. es gilt $\eta(Y)\circ f = P(f)\circ\eta(X)$.
Die Zuordnung $\eta$ ist eine sogenannte natürliche Transformation.

\begin{Definition}[Natürliche Transformation]
Seien $\category C, \category D$ Kategorien und
$F,G\colon \category C\to \category D$ Funktoren.
Dann schreibt man $\eta\colon F\to G$ und nennt $\eta$ natürliche
Transformation, wenn die folgenden beiden Axiome erfüllt sind:
\begin{enumerate}
\item Jedes Objekt $X\in\Ob(\category C)$ bekommt einen Morphismus
$\eta(X)\colon F(X)\to G(X)$.
\item Für jeden Morphismus $f\colon X\to Y$ gilt
$\eta(Y)\circ F(f)=G(f)\circ\eta(X)$.
\end{enumerate}
\end{Definition}
Die zweite Bedingung lässt sich übersichtlich als kommutierendes Diagramm
darstellen:
\[\xymatrix{
F(X) \ar[r]^{F(f)} \ar[d]_{\eta(X)} & F(Y) \ar[d]^{\eta(Y)} \\
G(X) \ar[r]_{G(f)} & G(Y)}\]
Ein weiteres Beispiel ergibt sich bezüglich Äquivalenzrelationen
in Erinnerung an \eqref{eq:induzierte-Abbildung}.
Eine Abbildung $f\colon M\to M'$ heiße \emph{induzierend}, wenn%
\[x\sim a \implies f(x)\sim' f(a).\]
\begin{Satz}
Die Paare $(M,\sim)$, bestehend aus Menge und Äquivalenzrelation,
bilden mit den induzierenden Abbildungen
als Morphismen bezüglich Verkettung eine Kategorie.
\end{Satz}
\strong{Beweis.}
Die identische Abbildung ist offensichtlich induzierend. Hat man
neben $f\colon M\to M'$ eine weitere induzierende Abbildung $g\colon M'\to M''$, dann
folgt $g(y)\sim'' g(b)$ aus $y\sim' b$. Aus $x\sim a$ folgt
mit $y:=f(x)$ und $b:=f(a)$ somit $g(f(x))\sim'' g(f(x))$.
Daher ist auch $g\circ f$ induzierend.\;\qedsymbol

Genau dann wenn $f$ induzierend ist, existiert eine induzierte
Abbildung
\[I(f)\colon M/\sim\to M'/\sim',\;\text{so dass}\;I(f)\circ\pi = \pi'\circ f,\]
wobei $\pi,\pi'$ jeweils die kanonische Projektion ist.
\begin{Satz}
Bei der Induktion $I$ handelt es sich um einen kovarianten Funktor.
\end{Satz}
\strong{Beweis.} Man betrachte das folgende kommutierende Diagramm:
\[\xymatrix{
M \ar[r]^{f} \ar[d]_{\pi}
& M' \ar[r]^{g} \ar[d]^{\pi'}
& M'' \ar[d]^{\pi''}\\
M/\sim \ar[r]_{I(f)}
& M'/\sim' \ar[r]_{I(g)}
& M''/\sim''}\]
Die Induktion $I$ besitzt die Eigenschaften
\begin{gather*}
I(f)\circ\pi = \pi'\circ f,\\
I(g)\circ\pi' = \pi''\circ g,\\
I(g\circ f)\circ\pi = \pi''\circ (g\circ f).
\end{gather*}
Damit kann man nun rechnen
\begin{equation}
I(g\circ f)\circ\pi = \pi''\circ g\circ f
= I(g)\circ\pi'\circ f = I(g)\circ I(f)\circ\pi.
\end{equation}
Infolge gilt $I(g\circ f)=I(g)\circ I(f)$, da die kanonische
Projektion $\pi$ eine Surjektion ist. Aus der Forderung $I(\id)\circ\pi
= \pi\circ\id = \pi$ ergibt sich $I(\id) = \id$,
da $\pi$ surjektiv ist.\;\qedsymbol

Die Abbildung $\eta((M,\sim)):=\pi$, die jeder Menge mit
Äquivalenzrelation ihre kanonische Projektion zuordnet,
ist eine natürliche Transformation.

\strong{Beispiel zur Vertiefung.}
Ein weiteres Beispiel berührt einen Grundbegriff der linearen Algebra.
Hat man einen Vektorraum $Y$ über dem Körper $K$, ohne dass wir jetzt
genau verstehen müssen was das bedeutet -- man stelle sich $Y:=\R$
und $K:=\R$ vor --, dann bildet für eine beliebige Menge $X\ne\emptyset$
auch $\Abb(X,Y)$ einen Vektorraum über diesem Körper bezüglich den
punktweisen Operationen%
\begin{gather*}
(\lambda f)(x) := \lambda f(x),\\
(f_1+f_2)(x) := f_1(x)+f_2(x),
\end{gather*} 
wobei $\lambda\in K$, $x\in X$, $f,f_1,f_2\in\Abb(X,Y)$. Die lineare
Algebra handelt von \emph{linearen Abbildungen}. Seien $V,W$
Vektorräume über dem Körper $K$, wobei auch $V=W$ sein darf. Eine
Abbildung $\varphi\colon V\to W$ heißt linear, falls%
\begin{gather*}
\forall v_1,v_2\in V\colon\;\varphi(v_1+v_2) = \varphi(v_1)+\varphi(v_2),\\
\forall \lambda\in K,v\in V\colon\;\varphi(\lambda v) = \lambda\varphi(v).
\end{gather*}
Man kann zeigen dass die Vektorräume mit den linearen Abbildungen
als Morphismen auch eine Kategorie bilden, aber darauf will ich an
dieser Stelle nicht hinaus. Hat man nun eine feste, aber beliebige
Abbildung $g\colon X'\to X$, dann ist%
\[\varphi\colon\Abb(X,Y)\to\Abb(X',Y),\quad \varphi(f):=f\circ g\]
eine lineare Abbildung, man spricht auch von einem
\emph{linearen Operator}, dem \emph{Kompositionsoperator}
$C_g=\varphi$. Die Bestätigung ist nicht sonderlich schwer, man
muss bloß blind die Definitionen einsetzen und dem Formalismus
folgen. Es gilt%
\begin{align*}
\varphi(\lambda f)(x) &= ((\lambda f)\circ g)(x)
= (\lambda f)(g(x)) = \lambda f(g(x))\\
&= \lambda (f\circ g)(x) = \lambda (\varphi(f))(x)
= (\lambda\varphi(f))(x),
\end{align*}
kurz $\varphi(\lambda f)=\lambda\varphi(f)$. Und es gilt
\begin{align*}
\varphi(f_1+f_2)(x) &= ((f_1+f_2)\circ g)(x)
= (f_1+f_2)(g(x)) = f_1(g(x)) + f_2(g(x))\\
&= (f_1\circ g)(x) + (f_2\circ g)(x)
= \varphi(f_1)(x)+\varphi(f_2)(x)\\
&= (\varphi(f_1)+\varphi(f_2))(x),
\end{align*}
kurz $\varphi(f_1+f_2)=\varphi(f_1)+\varphi(f_2)$.

Man bemerkt nun, dass diese Rechnungen lediglich auf der Eigenschaft
der Operationen beruhen, punktweise zu sein. Dies soll im Folgenden
präzisiert werden. Die Formulierung wollen wir allgemein für
Operationen beliebiger Stelligkeit haben. Sei also $p\colon Y^n\to Y$
eine $n$-stellige Operation, man stelle sich dabei $Y:=\R$ vor. Man
definiert nun die punktweise Anwendung von $p$ als%
\begin{equation}\label{eq:eta-punktweise}
\eta_p\colon\Abb(X,Y)^n\to\Abb(X,Y),\quad
\eta_p(f)(x) := p(f_1(x),\ldots,f_n(x)),
\end{equation}
wobei $f:=(f_1,\ldots,f_n)$ ein Tupel von Funktionen ist. Sei außerdem%
\begin{equation}
F(\varphi)(f) := (\varphi(f_1),\ldots,\varphi(f_n)).
\end{equation}
Zeigen wollen wir für $\varphi(f):=f\circ g$ nun
\begin{equation}
\varphi(\eta_p(f))(x) = \eta_p(F(\varphi)(f))(x),
\quad\text{kurz}\quad
\varphi\circ\eta_p = \eta_p\circ F(\varphi).
\end{equation}
Bei der Bestätigung folgt man wieder blind den Definitionen und dem
Formalismus. Es gilt%
\begin{align}
\varphi(\eta_p(f))(x)
&= \eta_p(f)(g(x)) = p(f_1(g(x)),\ldots,f_n(g(x)))\\
&= p(\varphi(f_1)(x),\ldots,\varphi(f_n)(x))
= \eta_p(\varphi(f_1),\ldots,\varphi(f_n))(x)\\
&= \label{eq:Ende-eta-punktweise}\eta_p(F(\varphi)(f))(x).
\end{align}
Das bedeutet, dieses Diagramm kommutiert:
\[\xymatrix{
\Abb(X,Y)^n \ar[r]^{F(\varphi)} \ar[d]_{\eta_p}
& \Abb(X',Y)^n \ar[d]^{\eta_p} \\
\Abb(X,Y) \ar[r]_{\varphi} & \Abb(X',Y)}\]
Das legt den Verdacht nahe, dass es sich bei $F$ um einen Funktor
und bei $\eta_p$ um eine natürliche Transformation handelt.

\begin{Satz}[Kategorie mit Kompositionsoperatoren als Morphismen]\mbox{}\\*
Sei $Y\ne\emptyset$. Sei $\Omega:=\{\Abb(X,Y)\mid X\;\text{ist beliebig}\}$.
Sei%
\begin{gather*}\Hom(\Abb(X,Y),\Abb(X',Y))\\
\quad := \{\varphi\mid\exists f{\in}\Abb(X,Y), g{\in}\Abb(X',X)\colon
\varphi = f\circ g\}.
\end{gather*}
Sei $\psi\circ \varphi$ die gewöhnliche Verkettung. Dann bildet
$\strong{Komp}:=(\Omega,\Hom,\circ)$ eine Kategorie.
\end{Satz}
\strong{Beweis.}
Zunächst müssen wir bestätigen, dass $\Hom$ bezüglich
$\circ$ abgeschlossen ist. Sei $\varphi_1(f):=f\circ g_1$ und
$\varphi_2(f):=f\circ g_2$ mit $\dom(\varphi_2)=\cod(\varphi_1)$, so
dass man $\varphi:=\varphi_2\circ\varphi_1$ bilden kann.
Gesucht ist ein $g$, so dass $\varphi = f\circ g$. Nun gilt%
\begin{equation}
(\varphi_2\circ\varphi_1)(f) = (f\circ g_1)\circ g_2
= f\circ g_1\circ g_2 = f\circ (g_1\circ g_2).
\end{equation}
Man kann also $g:=g_1\circ g_2$ setzen. Nun verbleibt bloß noch die
Existenz fester Identitäten $\id$ zu bestätigen. Man definiert dazu
$\id(f):=f\circ\id$. Für $\varphi(f):=f\circ g$ gilt dann%
\begin{gather}
(\varphi\circ\id)(f) = (f\circ \id)\circ g = f\circ g = \varphi(f),\\
(\id\circ\varphi)(f) = (f\circ g)\circ\id = f\circ g = \varphi(f).\;\qedsymbol
\end{gather}
\begin{Satz}
Bei $F(\varphi)(f):=(\varphi(f_1),\ldots,\varphi(f_n))$ für
$f=(f_1,\ldots,f_n)$ handelt es sich um einen kovarianten Funktor.
\end{Satz}
\strong{Beweis.} Es gilt
\begin{align}
F(\psi\circ\varphi)(f) &= (\psi(\varphi(f_1)),\ldots,\psi(\varphi(f_n)))
= F(\psi)(F(\varphi)(f))
= (F(\psi)\circ F(\varphi))(f),
\end{align}
kurz $F(\psi\circ\varphi) = F(\psi)\circ F(\varphi)$. Und es gilt
\begin{equation}
F(\id)(f) = (\id(f_1),\ldots,\id(f_n)) = (f_1,\ldots,f_n) = f = \id(f),
\end{equation}
kurz $F(\id) = \id$.\;\qedsymbol
\begin{Satz}
Die in \eqref{eq:eta-punktweise} definierte Operation $\eta_p$ ist
eine natürliche Transformation.
\end{Satz}
\strong{Beweis.} Wurde in \eqref{eq:eta-punktweise} bis
\eqref{eq:Ende-eta-punktweise} schon gezeigt.\;\qedsymbol

Jetzt haben wir so viele Funktoren kennengelernt, dass allgemeine
Regeln betreffend Funktoren gut motiviert sind. Der folgende Korollar
durchleuchtet die Funktoren ein wenig.
\begin{Korollar}
Wird ein Funktor auf einen Isomorphismus angewendet, ist das Resultat
wieder ein Isomorphismus.
\end{Korollar}
\strong{Beweis.} Dieser ist direkt aus den Definitionen zu erhalten.
Wir betrachten nur einen kovarianten Funktor $F$, weil der Beweis für
einen kontravarianten Funktor analog verläuft.

Sei $f\colon X\to Y$ ein beliebiger Isomorphismus im Definitionsbereich
des Funktors. Laut Definition existert eine Inverse, das heißt, ein
$g\colon Y\to X$ mit $g\circ f = \id_X$ und $f\circ g = \id_Y$. Gemäß
der definierenden Eigenschaft eines Funktors darf man rechnen
\begin{gather*}
\id_{F(X)} = F(\id_X) = F(g\circ f) = F(g)\circ F(f),\\
\id_{F(Y)} = F(\id_Y) = F(f\circ g) = F(f)\circ F(g).
\end{gather*}
Somit ist $F(f)$ ein Isomorphismus mit Inverse $F(g)$.\,\qedsymbol

Wird beispielsweise der Vergissfunktor von Gruppen zu Mengen auf
einen Gruppenisomorphismus angewendet, ist das Resultat zwingend
eine Bijektion, da die Bijektionen die Isomorphismen in der Kategorie
der Mengen sind.

\section{Anfangsobjekte und Endobjekte}

\begin{Definition}[Anfangsobjekt, Endobjekt, Nullobjekt]\mbox{}\\*
Es sei $\category C$ eine Kategorie. Ein
$A\in\Ob(\category C)$ heißt Anfangsobjekt, wenn
es zu jedem Objekt $X\in\Ob(\category C)$ genau einen
Morphismus $A\to X$ gibt.
Ein $E\in\Ob(\category C)$ heißt Endobjekt, wenn
es zu jedem Objekt $X\in\Ob(\category C)$ genau einen
Morphismus $X\to E$ gibt. Ein Objekt heißt Nullobjekt, wenn
es sowohl Anfangsobjekt als auch Endobjekt ist.
\end{Definition}

\noindent
\strong{Mengen.}
Wir untersuchen zunächst die Kategorie der Mengen. Anfangsobjekt
bedeutet hier eine Menge $A$, bei der es zu jeder Menge $X$ genau
eine Abbildung $A\to X$ gibt. Betrachten wir zunächst endliche
Mengen, dann ergibt sich aufgrund von Gleichung \eqref{eq:Anzahl-Abb}
ja die Bedingung $|X|^{|A|}=1$. Das geht nur, wenn $|A|=0$ ist,
und das bedeutet $A=\emptyset$. Die einzige Abbildung in $\Abb(\emptyset,X)$
ist die leere Abbildung, und dies bleibt auch dann richtig, wenn
$X$ gänzlich beliebig ist. Somit haben wir die leere Menge als
einziges Anfangsobjekt identifiziert.

Endobjekt bedeutet eine Menge $E$, so dass es zu jeder Menge $X$
genau eine Abbildung $X\to E$ gibt. Wieder beschränken wir uns
zunächst auf endliche Mengen und nutzen $\eqref{eq:Anzahl-Abb}$.
Wir erhalten die Bedingung $|E|^{|X|}=1$. Das geht nur, wenn
$|E|=1$ ist. Jede Menge mit einem Element ist also Endobjekt,
denn allgemein gibt es dann nur eine einzige Abbildung, nämlich
die konstante Abbildung. Dies bleibt auch dann richtig, wenn
$X$ gänzlich beliebig ist.

Ein Nullobjekt existiert offenbar nicht.

Benutzen wir doch $1:=\{\emptyset\}$ als kanonisches Endobjekt.
Interessant ist, dass man zu einer Menge $X$ jedes Element
$x\in X$ mit der Abbildung $x\colon 1\to X$ identifizieren kann,
für die $x(\emptyset)=x$ gilt. Zu einer Abbildung $f\colon X\to Y$
können wir die Zuordnung $f(x)=y$ bzw. $(x,y)\in f$ nun in der Form
$f\circ x = y$ beschreiben.

\section{Produkt und Koprodukt}

Ein wichtiger Begriff der Theorie ist das Produkt von
Objekten. Weil es sich dabei um eine Verallgemeinerung des
kartesischen Produktes von Mengen handelt, möchte ich die
Zusammenhänge zunächst am vertrauten Schauplatz der Mengen
betrachten.

Zu zwei Mengen $A_1, A_2$ können wir das Produkt
$A_1\times A_2$ bilden. Man definiert die Projektionen
auf die Komponenten als
\begin{align*}
&\pi_1\colon A_1\times A_2\to A_1,\quad\pi_1((x,y)) := x,\\
&\pi_2\colon A_1\times A_2\to A_2,\quad\pi_2((x,y)) := y.
\end{align*}
Nun betrachten wir Abbildungen $f_1\colon X\to A_1$ und
$f_2\colon X\to A_2$. Zunächst sei $X:=\{\emptyset\}$. Die jeweilige
Abbildung pickt dann ein Element aus der jeweiligen Menge heraus,
das sind $a_1:=f_1(\emptyset)$ und $a_2:=f_2(\emptyset)$. Nun
ist eine Abbildung $f$ gesucht, sodass das Diagramm
\[\xymatrix{
& X\ar[dl]_{f_1}\ar[dr]^{f_2}\ar[d]^f & \\
A_1 & A_1\times A_2 \ar[l]^{\pi_1}\ar[r]_{\pi_2} & A_2
}\]
kommutiert. Die Bedingungen an $f$ sind also
$\pi_i\circ f = f_i$ für $i\in\{1,2\}$. Damit ist aber eindeutig
festgelegt, dass $f(\emptyset)=(a_1,a_2)$ sein muss, denn ein
Tupel ist durch die Komponenten festgelegt, und die sind
$\pi_i(f(\emptyset)) = f_i(\emptyset) = a_i$ für $i\in\{1,2\}$.
Somit ist $f$ eindeutig bestimmt.

Die Betrachtung kann man genauso für eine allgemeine Menge $X$ führen,
weil die Argumentation dann jeweils für jedes Element von $X$ gilt.
Wieder ist $f$ eindeutig bestimmt.

Gelegentlich wird $f$ als $f=f_1\times f_2$ notiert.

\begin{Definition}[Produkt]\newlinefirst
Sei $C$ eine Kategorie und seien $Y_1,Y_2\in C$.
Ein Objekt $Y\in\mathcal C$ mit Projektionen $\pi_1\colon Y\to Y_1$ und
$\pi_2\colon Y\to Y_2$ heißt Produkt, wenn zu jedem Objekt
$X\in C$ und allen Morphismen $f_1\colon X\to Y_1$
und $f_2\colon X\to Y_2$ genau ein Morphismus $f\colon X\to Y$
existiert, so dass $f_1=\pi_1\circ f$ und $f_2=\pi_2\circ f$.
\end{Definition}

\noindent
Das kartesische Produkt $Y:=Y_1\times Y_2$ ist ein Produkt in
der Kategorie der Mengen. Wir schreiben $f(x)=y$ mit $y=(y_1,y_2)$.
Nun ist $\pi_1(y)=y_1$ und $\pi_2(y)=y_2$, laut Forderung soll also
$y_1=f_1(x)$ und $y_2=f_2(x)$ sein. Dadurch ist $f$ mit
$f(x):=(f_1(x),f_2(x))$ eindeutig festgelegt.

Zu zwei Mengen können wir weiterhin die disjunkte Vereinigung
$X_1\sqcup X_2$ bilden. Wir rekapitulieren, dass zu ihr die beiden
kanonischen Injektionen
\begin{align*}
& i_1\colon X_1\to X_1\sqcup X_2,\quad i_1(x) := (1,x),\\
& i_2\colon X_2\to X_1\sqcup X_2,\quad i_2(x) := (2,x)
\end{align*}
gehören. Man stellt sich nun die Frage, was das Wesensmerkmal
der disjunkten Vereinigung ist. Das ist doch, dass zu jedem ihrer
Elemente die Information vorliegt, ob es aus der linken oder der
rechten Menge entstammt. Das heißt, es muss eine Abbildung geben,
die auf den Tag projiziert. Betrachten wir dazu die Abbildungen
$f_1\colon X_1\to Y$ und $f_2\colon X_2\to Y$ mit $Y:=\{1,2\}$ und
$f_k(x):=k$. Mit der Abbildung $f_k$ gelangt man von $X_k$ also
direkt zum Tag $k$. Nun ist eine Abbildung $f$ gesucht, so dass das
Diagramm
\[\xymatrix{
X_1\ar[r]^{i_1}\ar[dr]_{f_1} & X_1\sqcup X_2\ar[d]^{f}
& X_2\ar[l]_{i_2}\ar[dl]^{f_2}\\
& Y &
}\]
kommutiert. Das heißt, es soll $f\circ i_k = f_k$ für $k\in\{1,2\}$ sein.
Das macht $f((1,x)) = 1$ und $f((2,x)) = 2$. Dadurch ist $f$ eindeutig
bestimmt. Es ist die gesuchte Projektion auf den Tag.

\begin{Definition}[Koprodukt]\newlinefirst
Sei $C$ eine Kategorie und seien $X_1,X_2\in C$. Ein Objekt $X$
mit Morphismen $i_1\colon X_1\to X$ und $i_2\colon X_2\to X$
heißt Koprodukt, wenn zu jedem Objekt $Y\in C$ und allen Morphismen
$f_1\colon X_1\to Y$ und $f_2\colon X_2\to Y$ genau ein Morphismus
$f\colon X\to Y$ existiert, so dass $f_1 = f\circ i_1$
und $f_2 = f\circ i_2$.
\end{Definition}

\noindent
Die disjunkte Vereinigung $X:=X_1\sqcup X_2$ mit den Injektionen
$i_1(x):=(1,x)$ und $i_2(x):=(2,x)$ ist ein Koprodukt in der
Kategorie der Mengen. Es gilt schon mal
\[f(x) = \strong{match}\; x \begin{cases}
(1,x)\mapsto y_1,\\
(2,x)\mapsto y_2.
\end{cases}\]
Laut Forderung soll außerdem $y_1 = f_1(x)$ und $y_2 = f_2(x)$ sein,
wodurch $f$ eindeutig festgelegt ist.


\chapter{Diskrete Mathematik}

\section{Kombinatorik}

\subsection{Endliche Summen}

Problemstellungen der Kombinatorik führen oftmals zu Summen.
Ein Grund dafür liegt darin, dass die Lösungen bestimmter
Differenzengleichungen als Summen darstellbar sind. In der Analysis
sind die Reihen sehr bedeutsam, das sind Summen von unendlich vielen
Summanden. Der Wert einer Reihe wird hierbei, sofern existent, als der
Grenzwert der Folge der Partialsummen erklärt.
Somit bauen Reihen konzeptuell auf den endlichen Summen auf.
Ein weiteres wichtiges Werkzeug stellen die formalen Potenzreihen bereit,
die ohne den Grenzwertbegriff auskommen, aber analoge Rechenregeln zu
den eigentlichen Potenzreihen aufweisen. In der linearen Algebra dienen
Summen ebenfalls zur allgemeinen Formulierung von Rechenregeln.
Darin einbegriffen ist das Rechnen mit Polynomen. Auch in der
Stochastik sind Summen dienlich, zum Beispiel zur Erklärung des
Erwartungswertes einer Zufallsgröße.

Zur Einführung eine kleine Aufgabe. Man beobachtet
\[1+3 = 2^2,\quad 1+3+5 = 3^2,\quad 1+3+5+7 = 4^2,\quad 1+3+5+7+9 = 5^2.\]
Anscheinend besteht die Gesetzmäßigkeit
\[1 + 3 + \ldots  + (2n-1) = n^2.\]
In Worten ist die Summe der ersten ungeraden Zahlen eine Quadratzahl,
genauer die Anzahl der Summanden zum Quadrat. Ein strenger Beweis dafür
findet sich unschwer. Zuvor will ich aber die allgemeinen Gesetzmäßigkeiten
für das Rechnen mit Summen diskutieren.

Statt Summen mit drei Pünktchen zu beschreiben, verwendet man meist die
unmissverständliche Notation
\[\textstyle\sum_{k=1}^n a_k := a_1 + a_2 + \ldots + a_n.\]
Beispielsweise gilt
\[\textstyle\sum_{k=1}^3 (2k-1) =
(2\cdot 0-1) + (2\cdot 1 -1 ) + (2\cdot 2 - 1)
= 1 + 3 + 5.\]
Die drei Pünktchen treten allerdings noch in der Erklärung der Summe
auf. Wir werden sie wieder gänzlich los, indem die Summe rekursiv
festgelegt wird.

\begin{Definition}[Summe]\newlinefirst
Für eine Folge $a\colon\Z\to\R$ ist die \emph{Summe} der $a_k$ für $k=m$
bis $k=n$ rekursiv definiert gemäß
\[\sum_{k=m}^{m-1} a_k := 0,\quad \sum_{k=m}^n a_k := a_n + \sum_{k=m}^{n-1} a_k.\]
\end{Definition}
Es darf $k$, genannt der \emph{Laufindex}, als gebundene Variable nur
im Term $a_k$ auftreten. Die Zahlen $m,n$ nennt man die \emph{Grenzen}
der Summierung. Ist $a$ lediglich auf $D\subseteq\Z$ definiert, können
wir $a$ schlicht per $a_k := 0$ für $k\notin D$ auf ganz $\Z$ erweitern.
Dies bietet sich an, wenn die Summe stets definiert sein soll,
unabhängig davon wie die Grenzen gewählt werden.

\begin{Satz}\label{Summe-Homogenitaet}
Es gilt $\sum_{k=m}^n ca_k = c\sum_{k=m}^n a_k$ für jede Konstante $c$.
\end{Satz}
\begin{Beweis}
Induktion über $n$. Im Induktionsanfang $n=m-1$ gilt
\[\textstyle\sum\limits_{k=m}^{m-1} ca_k \stackrel{\mathrm{def}}= 0 = c\cdot 0
\stackrel{\mathrm{def}}= c\sum\limits_{k=m}^{m-1} a_k.\]
Der Induktionsschritt ist
\[\textstyle\sum\limits_{k=m}^n ca_k \stackrel{\mathrm{def}}= ca_n+\sum\limits_{k=m}^{n-1} ca_k
\stackrel{\mathrm{IV}}= ca_n + c\sum\limits_{k=m}^{n-1} a_k
= c\Big(a_n + \sum\limits_{k=m}^{n-1} a_k\Big) \stackrel{\mathrm{def}}= c\sum\limits_{k=m}^n a_k.\,\qedsymbol\]
\end{Beweis}

\begin{Satz}\label{Summe-Additivitaet}
Es gilt $\sum_{k=m}^n (a_k+b_k) = \sum_{k=m}^n a_k + \sum_{k=m}^n b_k$.
\end{Satz}
\begin{Beweis}
Induktion über $n$. Im Induktionsanfang $n=m-1$ reduziert sich die
Behauptung zu $0 = 0 + 0$. Der Induktionsschritt ist%
\begin{align*}
\textstyle\sum_{k=m}^n (a_k + b_k) &\textstyle\stackrel{\mathrm{def}}=
a_n + b_n + \sum_{k=m}^{n-1} (a_k + b_k)
\stackrel{\mathrm{IV}}=
a_n + b_n + \sum_{k=m}^{n-1} a_k + \sum_{k=m}^{n-1} b_k\\
&\textstyle\stackrel{\mathrm{def}}= \sum_{k=m}^n a_k + \sum_{k=m}^n b_k.\,\qedsymbol
\end{align*}
\end{Beweis}
Man kann den Raum $\Abb(\Z,\R)$ auch als einen Vektorraum über dem Körper
$\R$ deuten, wobei die Multiplikation mit einem Skalar $c\in\R$ und die
vektorielle Addition punktweise definiert werden, das heißt,%
\[(ca)_k := ca_k,\quad\; (a+b)_k := a_k+b_k.\]
Dann ist $\sum_m^n a := \sum_{k=m}^n a_k$ eine lineare Abbildung,
denn die Homogenität gilt laut Satz \ref{Summe-Homogenitaet}, und
die Additivität laut Satz \ref{Summe-Additivitaet}.

Der Differenzoperator
\[(\Delta a)_n := a_{n+1} - a_n\]
wird die nachfolgende Diskussion übersichtlicher halten. Es handelt
sich ebenfalls um eine lineare Abbildung, das heißt, für jede Konstante
$c$ und alle Folgen $a,b$ gilt $\Delta (ca) = c\Delta a$ und
$\Delta(a+b) = \Delta a + \Delta b$, denn%
\begin{align*}
(\Delta (ca))_n &= ca_{n+1} - ca_n = c(a_{n+1}-a_n) = c(\Delta a)_n,\\
(\Delta(a+b))_n &= (a_{n+1}+b_{n+1}) - (a_n + b_n)
= a_{n+1} - a_n + b_{n+1} - b_n\\
&= (\Delta a)_n + (\Delta b)_n.
\end{align*}

\begin{Satz}[Teleskopsumme]
Es gilt $\sum_{k=m}^{n-1} (\Delta a)_k = a_n - a_m$.
\end{Satz}
\begin{Beweis}
Induktion über $n$. Im Induktionsanfang $n=m$ sieht man trivial,
dass beide Seiten der Gleichung den Wert null haben. Der
Induktionsschritt ist%
\[\sum\limits_{k=m}^n (\Delta a)_k \stackrel{\mathrm{def}}=
(\Delta a)_n + \sum\limits_{k=m}^{n-1} (\Delta a)_k
\stackrel{\mathrm{IV}}= a_{n+1} - a_n + a_n - a_m = a_{n+1} - a_m.\,\qedsymbol\]
\end{Beweis}

\noindent
Das Prinzip der Teleskopsummen bietet zwei wesentliche Vorteile. 
Der erste Vorteil besteht darin, dass man damit bequem nach
Summenformeln fischen kann, wobei man deren Beweis mit geschenkt
bekommt. Für $s_n := (n-1)/n$ gilt zum Beispiel%
\[(\Delta s)_n = \frac{n}{n+1} - \frac{n-1}{n}
= \frac{n^2 - (n+1)(n-1)}{(n+1)n} = \frac{1}{(n+1)n}.\]
Ergo findet sich die Summenformel
\[\sum_{k=1}^n \frac{1}{(k+1)k} = \frac{n}{n+1}.\]
Der zweite Vorteil besteht darin, dass sich der Beweis einer bereits
ermittelten Summenformel dahingehend vereinfacht, dass nicht jedes
mal ein Induktionsbeweis geführt werden muss, womit die unter Umständen
schwierige Auffindung der wesentlichen Umformungen im Induktionsschritt
erleichtert wird. Zum Beweis einer Summenformel der Form%
\[\textstyle\sum_{k=1}^n a_k = s_{n+1} - s_1\]
genügt es nämlich, $\Delta s = a$ zu prüfen. Zur Summenformel%
\[\textstyle\sum_{k=1}^n (2k-1) = n^2\]
gilt beispielsweise $a_n = 2n-1$ und $s_n = (n-1)^2$. Es bestätigt sich
kurzerhand%
\[(\Delta s)_n = n^2 - (n-1)^2 = n^2 - n^2 + 2n - 1  = 2n-1 = a_n.\]
Aus der nun bewiesenen Formel ziehen sich schnell weitere Schlüsse.
Setzt man in die Gleichung%
\[\textstyle n^2 = \sum_{k=1}^n (2k-1) = 2\sum_{k=1}^n k - \sum_{k=1}^n 1\]
die Vereinfachung $\sum_{k=1}^n 1 = n$ ein und stellt sie daraufhin nach
der verbleibenden Summe um, findet sich die Summenformel%
\[\sum_{k=1}^n k = \frac{1}{2}(n^2 + n) = \frac{n}{2}(n+1).\]
Die Werte dieser Summe werden die \emph{Dreieckszahlen} genannt.%
\footnote{Engl. \emph{triangular numbers}. Eintrag A000217 in OEIS.}

\begin{Satz}[Aufteilung von Summen]%
\label{Summe-Aufteilung}\newlinefirst
Für $m\le p\le n$ gilt
$\sum_{k=m}^n a_k = \sum_{k=m}^{p-1} a_k+\sum_{k=p}^n a_k$.
\end{Satz}
\begin{Beweis}
Induktion über $n$. Der Induktionsanfang liege in $n=p$. Die Gleichung
ist in diesem erfüllt, weil $\sum_{k=p}^p a_k=a_p$ gilt.
Der Induktionsschritt ist%
\[\sum\limits_{k=m}^n a_k \stackrel{\mathrm{def}}= a_n+\sum\limits_{k=m}^{n-1} a_k
\stackrel{\mathrm{IV}}= a_n + \sum\limits_{k=m}^{p-1} a_k + \sum\limits_{k=p}^{n-1} a_k
\stackrel{\mathrm{def}}= \sum\limits_{k=m}^{p-1} a_k + \sum\limits_{k=p}^n a_k.\,\qedsymbol\]
\end{Beweis}

\begin{Satz}[Indexverschiebung]%
\label{Summe-Indexshift}\newlinefirst
Für die Indexverschiebung der Distanz $d\in\Z$,
kurz Indexshift, gilt\\
$\sum_{k=m}^n a_k = \sum_{k=m+d}^{n+d} a_{k-d}.$
\end{Satz}
\begin{Beweis}
Induktion über $n$. Im Induktionsanfang $n=m-1$ wird unmittelbar
ersichtlich, dass beide Seiten der Gleichung im Wert null resultieren.
Der Induktionsschritt gilt gemäß der Umformung%
\begin{align*}
\sum_{k=m}^n a_k \stackrel{\mathrm{def}}= a_n+\sum_{k=m}^{n-1} a_k
\stackrel{\mathrm{IV}}= a_{(n+d)-d}+\sum_{k=m+d}^{n+d-1} a_{k-d}
\stackrel{\mathrm{def}}= \sum_{k=m+d}^{n+d} a_{k-d}.\,\qedsymbol
\end{align*}
\end{Beweis}

\noindent
Ich will zusätzlich noch eine alternative Herleitung aufzeigen. Bei
dieser wird der Indexbereich als Ungleichung beschrieben, was den
Vorteil bietet, dass die Ungleichung einer Umformung unterzogen werden
kann. Mit der Substitution $k=k'-d$ gelingt dahingehend die Überlegung%
\begin{align*}
\sum_{k=m}^n a_k = \sum_{m\le k\le n} a_k
= \!\!\sum_{m\le k'-d\le n}\!\! a_{k'-d}
= \!\!\sum_{m+d\le k'\le n+d}\!\! a_{k'-d}
= \sum_{k'=m+d}^{n+d} a_{k'-d}.
\end{align*}

\begin{Satz}[Umkehrung der Reihenfolge]%
\label{Summe-Umkehrung}\newlinefirst
Es gilt $\sum_{k=0}^n a_k = \sum_{k=0}^n a_{n-k}$.
\end{Satz}
\strong{Beweis.}
Der Induktionsanfang bei $n=0$ ist trivial. Beim Induktionsschritt
macht man sich Satz \ref{Summe-Indexshift} (Indexshift) und
Satz \ref{Summe-Aufteilung} (Aufteilung) zunutze,%
\begin{gather*}
\sum_{k=0}^n a_{n-k} = a_{n-n}+\sum_{k=0}^{n-1} a_{n-k}
= a_0+\sum_{k=0}^{n-1} a_{n-(n-1-k)}
= a_0+\sum_{k=0}^{n-1} a_{k+1}\\
\stackrel{[k:=k-1]}= a_0+\sum_{k=1}^n a_k
= \sum_{k=0}^0 a_k + \sum_{k=1}^n a_k
= \sum_{k=0}^n a_k.\,\qedsymbol
\end{gather*}
Dienlich ist die Umkehrung beispielsweise bei der folgenden klassischen
Herleitung der Summenformel $\sum_{k=0}^n k = \tfrac{n}{2}(n+1)$. 
Aufbauend auf den bereits gezeigten Regeln stellt diese Herleitung ebenfalls
einen strengen Beweis dar. Das doppelte der Summe betrachtet man hierbei
als Addition der Summe zu sich selbst, wobei die Reihenfolge der zweiten
Summe umgekehrt wird. Man erhält%
\begin{align*}
2\sum_{k=0}^n k &= \sum_{k=0}^n k+\sum_{k=0}^n k
= \sum_{k=0}^n k + \sum_{k=0}^n (n-k)\\
&= \sum_{k=0}^n (k+n-k) = \sum_{k=0}^n n
= n\sum_{k=0}^n 1 = n(n+1).
\end{align*}
Ganz allgemein dürfen die Summanden einer endlichen Summe beliebig
umgeordnet werden; dies geht aus wiederholter Anwendung des Kommutativgesetzes
und des Assoziativgesetzes hervor. Die Formalisierung dieser Tatsache
geschieht vermittels einer frei wählbaren Permutation des Indexbereichs.
Der zuvor diskutierte Sachverhalt stellt sich dahingehend als
Spezialfall mit der auf der Menge $\{0,\ldots,n\}$ definierten
Permutation $\pi(k) = n-k$ dar.

\begin{Satz}[Permutation der Reihenfolge]%
\label{Summe-Permutation-Index}\newlinefirst
Sei $M=\{k\in\Z\mid m\le k\le n\}$. Für jede Permutation
$\pi\colon M\to M$ gilt
\[\sum_{k=m}^n a_k = \sum_{k=m}^n a_{\pi(k)}.\]
\end{Satz}
\begin{Beweis} Induktion über $n$. Sei ohne Beschränkung der Allgemeinheit $m=1$.
Im Induktionsanfang $n=0$ und $n=1$ ist die Gleichung offenkundig erfüllt.

Induktionsschritt. Induktionsvoraussetzung sei die Gültigkeit für $M$.
Zu zeigen ist die Gültigkeit für $M\cup\{n+1\}$.

Sei $t$ ein fester Parameter mit $1\le t\le n+1$.
Im Fall $\pi(t) = n+1$ geht man wie folgt vor.
Man setze $\sigma(k):=\pi(k)$ für $1\le k\le t-1$. Man setze
$\sigma(k):=\pi(k+1)$ für $t\le k\le n$. Weil $n+1$ kein Wert von
$\sigma$ ist, muss $\sigma$ eine Permutation $\sigma\colon M\to M$ sein.
Ergo gilt
\begin{align*}
\sum_{k=1}^{n+1} a_{\pi(k)} &= \sum_{k=1}^{t-1} a_{\pi(k)}
+ a_{\pi(t)} + \sum_{k=t+1}^{n+1} a_{\pi(k)}
= a_{\pi(t)} + \sum_{k=1}^{t-1} a_{\pi(k)}
+ \sum_{k=t}^n a_{\pi(k+1)}\\
&= a_{n+1} + \sum_{k=1}^{t-1} a_{\sigma(k)}
+ \sum_{k=t}^n a_{\sigma(k)}
= a_{n+1} + \sum_{k=1}^n a_{\sigma(k)}\\
&\stackrel{\mathrm{IV}}= a_{n+1} + \sum_{k=1}^n a_k
= \sum_{k=1}^{n+1} a_k.
\end{align*}
Man beachte, dass in den beiden Randfällen $t=1$ und $t=n+1$ die
jeweilige Randsumme den Wert null hat und somit verschwindet.\,\qedsymbol
\end{Beweis}

\noindent
Dass die Reihenfolge der Summierung keine Rolle spielt, gibt Anlass
zu einer abstrakten Fassung der Summierung selbst, die von vornherein
auf die Festlegung einer bestimmten Reihenfolge verzichtet.

\begin{Satz}
Sei $M$ eine endliche Menge, sei $a\colon M\to\R$. Die Festlegung
\[\sum_{k\in M} a_k := \sum_{i=m}^n a_{f(i)},\;\; |M| = n-m+1,\;\;
f\colon\{m,\ldots,n\}\to M\;\text{bijektiv}\]
ist wohldefiniert, wobei $f$ frei gewählt werden darf.
\end{Satz}
\begin{Beweis}
Es gilt zu zeigen, dass die Summe auch wirklich unabhängig von der
gewählten Bijektion ist. Sind $f,g$ zwei solcher Bijektionen, so ist
$\pi := g^{-1}\circ f$ ebenfalls bijektiv, also eine Permutation des
Indexbereichs. Vermittels Satz \ref{Summe-Permutation-Index} gilt daher
\[\sum_{i=m}^n a_{f(i)} = \sum_{i=m}^n a_{g(\pi(i))}
= \sum_{i=m}^n a_{g(i)}.\,\qedsymbol\]
\end{Beweis}

\noindent
Im Weiteren ist man nun befähigt, die Begrenzung des Indexbereichs
durch Ungleichungen streng zu definieren. Man legt fest
\[\sum_{m\le k\le n} a_k := \sum_{k\in M} a_k,
\quad M := \{k\in\Z\mid m\le k\land k\le n\}.\]
Somit ist streng begründet, dass die Ungleichungen $m\le k$
und $k\le n$ einer Äquivalenzumformung unterzogen werden dürfen,
denn von diesen bleibt die Menge $M$, und somit die Summe unberührt.

Ein letztes Detail darf aber nicht vergessen werden. Im Vorfeld der
oder Anschluss an die Umformung der Ungleichungen wird ggf. eine
Substitution der Laufvariable durchgeführt. Streng genommen muss noch
abgeklärt werden, dass dies unter allen Umständen unverfänglich ist.
Gewissheit schafft der

\begin{Satz}[Substitutionsregel]
Ist $\varphi\colon M'\to M$ eine Bijektion, gilt
\[\sum_{k\in M} a_k = \sum_{k'\in M'} a_{\varphi(k')}.\]
\end{Satz}
\begin{Beweis}
Zur Bijektion $f\colon\{1,\ldots,|M|\}\to M$ existiert die Bijektion
$g$ mit $f = \varphi\circ g$, nämlich ist dies $g:=\varphi^{-1}\circ f$.
Infolge gilt
\[\sum_{k\in M} a_k = \sum_{i=1}^{|M|} a_{f(i)}
= \sum_{i=1}^{|M|} a_{\varphi(g(i))}
= \sum_{k'\in M'} a_{\varphi(k')}.\,\qedsymbol\]
\end{Beweis}
Für die allgemeine Form der Summierung gelten anloge Regeln. Genauer
lassen sich die bereits bewiesenen Regeln direkt übertragen.
Man prüft unschwer nach, dass gilt
\[\sum_{k\in M} ca_k = c\sum_{k\in M}a_k,\qquad\sum_{k\in M} (a_k+b_k)
= \sum_{k\in M} a_k + \sum_{k\in M} b_k.\]
Zur Aufteilung sei $M = A\cup B$ mit $A\cap B = \emptyset$, also
eine disjunkte Zerlegung von $M$ in $A,B$. Dann gilt
\[\sum_{k\in M} a_k = \sum_{k\in A} a_k + \sum_{k\in B} a_k.\]
Bezüglich der Indikatorfunktion $1_A\colon M\to\{0,1\}$
zu $A\subseteq M$ gilt nun
\[\sum_{k\in M} 1_A(k)a_k = \sum_{k\in A} a_k,\]
denn $M=A\cup (M\setminus A)$ ist eine disjunkte Zerlegung, so dass gilt
\[\sum_{k\in M} 1_A(k)a_k = \sum_{k\in A} \underbrace{1_A(k)}_{=1}a_k
+ \sum_{k\in M\setminus A} \underbrace{1_A(k)}_{=0}a_k
= \sum_{k\in A} a_k.\]
Ist $A\subseteq\Z$ endlich, dürfen wir daher definieren
\[\sum_k 1_A(k)a_k = \sum_{k\in\Z} 1_A(k)a_k := \sum_{k\in A} a_k,\]
da diese Summe in Wirklichkeit endlich ist, also niemals über
unendlich viele Indizes läuft. Diese Sichtweise bietet unter Umständen
Vorteile bei der Umformung von Summen.

\subsection{Endliche Produkte}

Endliche Produkte weisen analoge Rechenregeln zu den endlichen Summen
auf, kommen allerdings weit weniger häufig vor als diese. Infolgedessen
verlaufen die Beweise ebenso analog, weshalb ich darauf verichten will,
diese weitläufig aufzuführen. Verbreitet sind endliche Produkte von
Zahlen vor allem in der Kombinatorik, zum Beispiel in Form von Faktoriellen
und Binomialkoeffizienten.

Für eine Folge $a\colon\Z\to\R$ definiert man rekursiv
\[\prod_{k=m}^{m-1} a_k := 1,\quad\;
\prod_{k=m}^n a_k := a_n\prod_{k=m}^{n-1} a_k.\]
Sei $c\in\Z_{\ge 0}$ eine Konstante. Es gilt
\[\prod_{k=m}^n a_k^c = \Big(\prod_{k=m}^n a_k\Big)^c,\qquad
\prod_{k=m}^n (a_k b_k) = \Big(\prod_{k=m}^n a_k\Big)\Big(\prod_{k=m}^n b_k\Big).\]
Die Aufteilung von Produkten geht vonstatten gemäß
\[\prod_{k=m}^n a_k = \Big(\prod_{k=m}^{p-1} a_k\Big)
\Big(\prod_{k=p}^n a_k\Big), \quad m\le p\le n+1.\]
Für jede Konstante $c\in\R$ gilt
\[\prod_{k=m}^n c = c^{n+1-m},\qquad
\prod_{k=m}^n (ca_k) = c^{n+1-m}\prod_{k=m}^n a_k.\]
Bei manchen Umformungen hilft das Teleskopprodukt
\[\prod_{k=m}^n\frac{a_{k+1}}{a_k} = \frac{a_{n+1}}{a_m},\quad
\text{sofern}\;\,\forall k\in\{m,\ldots,n\}\colon a_k\ne 0.\]

\subsection{Kombinatorische Grundaufgaben}

Abstrakt gesehen thematisiert die abzählende Kombinatorik die Frage
nach der Anzahl der Elemente endlicher Mengen. Die Mengenlehre
hilft bei der Formalisierung, und infolgedessen beim strengen Beweis
kombinatorischer Gesetzmäßigkeiten.

Zur Einführung stellt man häufig die Frage, wie viele mögliche Tupel
eine Produktmenge enthält. Üblicherweise formuliert man diese
unkomplizierte Aufgabe aber nicht in der Gestalt der Mengenlehre,
sondern fragt, wie viele Zahlen eine Binärsequenz aus $k$ Bits enthält,
oder wie viele Zeichenketten der Länge $k$ mit 26 Buchstaben
formulierbar sind. Abstrakt gesehen geht es bei den Binärsequenzen um
die Anzahl der Elemente von $\{0,1\}^k$. Analog geht es bei den
Zeichenketten um die Anzahl der Elemente von
\[\{\verb|'a'|,\verb|'b'|,\ldots,\verb|'z'|\}^k.\]
Die Elemente sind jeweils Tupel der Länge $k$. Bezüglich $B^k$ will
ich $|B|$ so wie bei Stellenwertsystemen als die \emph{Basis}
bezeichnen. Dass $|A\times B| = |A|\cdot |B|$ für beliebige Mengen
gilt, wurde bereits aufgezeigt. Demnach gilt
\[|B^k| = |B|^k = n^k, \quad n:=|B|.\]
Mit den $n=2$ Binärwerten sind dies 256 Binärsequenzen der Länge $k=8$.
Oder mit $n=26$ Buchstaben sind es 676 Zeichenketten der Länge $k=2$,
und schon 17576 Zeichenketten der Länge $k=3$.

In der Vergangenheit empfahlen Dienste häufig oder erzwangen penetrant,
Passwörter dadurch sicherer zu machen, indem neben Buchstaben auch noch
Ziffern und kryptische Sonderzeichen eingegeben werden sollten. Man
empfahl also, die Basis $n$ zu vergrößern, um die Länge gering zu halten.
Nimmt man an, dass so ein Passwort rein zufällig ausgewürfelt wird, ist
es bei bereits hinreichend großer Basis aber wesentlich sinnvoller,
abermals den Exponenten $k$ zu erhöhen. Bei Hinzunahme der 10 Ziffern zu
den 26 Buchstaben bekommt man mit $n=36$ für $k=3$ eine Zahl von 46656
Zeichenketten. Unterscheidet man zudem zwischen kleinen und großen
Buchstaben, sind es immerhin 238328. Mit $n=26$ und $k=4$ sind es aber
bereits 456976.

In der professionellen Kryptografie nennt man $n^k$ die
Schlüsselraumgröße, die bei modernen symmetrischen Kryptosystemen
mindestens $2^{128}$ beträgt. Und der Schlüssel wird am besten rein
zufällig gewählt. Zur Kodierung in der Basis $n$ muss man die Gleichung
$2^{128} = n^k$ nach $k$ umformen, das macht
\[k = 128\,\tfrac{\mathrm{lg}(2)}{\mathrm{lg}(n)}.\]
Das sind acht Blöcke je vier Hexadezimalziffern. Oder 22 Zeichen
zur Basis 62. Es ist davon auszugehen, dass dies ebenso sehr ein sicheres
Passwort darstellt.

Zur Erinnerung, ein Tupel aus $B^k$ darf auch als eine endliche Folge
interpretiert werden, das ist eine Abbildung
$\{0,\ldots,k-1\}\to B.$
Dahingehend ergibt sich die abstrakte Fassung

\begin{Satz}[Anzahl der Abbildungen]\newlinefirst
Seien $X,Y$ endliche Mengen mit $|X| = k$ und $|Y|=n$. Die Menge
der Abbildungen $X\to Y$ enthält $n^k$ Elemente.
\end{Satz}
\begin{Beweis}
Induktion über $k$. Im Anfang $k=0$ ist $X=\emptyset$. Es gibt genau
eine Abbildung $\emptyset\to Y$, nämlich die leere Abbildung.
Gleichermaßen ist $n^0=1$.

Zum Induktionsschritt. Induktionsvoraussetzung sei die Gültigkeit
für $k-1$. Es  sei $|X|=k$ und $|Y|=n$. Gesucht ist die Anzahl
der Möglichkeiten zur Festlegung der Abbildung $f\colon X\to Y$.
Sei $x\in X$ fest. Für die Festlegung $f(x)=y$ bestehen nun genau $n$
Möglichkeiten, nämlich so viele, wie es Elemente $y\in Y$ gibt.
Die Festlegung der übrigen Werte geschieht gemäß der Einschränkung
von $f$ auf $X\setminus\{x\}$, und die Zahl dieser Abbildungen beträgt
$n^{k-1}$ gemäß der Induktionsvoraussetzung. Man hat also $n$ mal
$n^{k-1}$ Möglichkeiten, das sind $n^k$.\,\qedsymbol
\end{Beweis}
Gleichwohl die dargelegte Beweisführung recht genau erscheint,
bezieht sie sich im Induktionsschritt an der wesentlichen Stelle
auf rationale Überlegungen, die allerdings nicht aus bereits
formal bewiesenen Regeln der Mengenlehre entstammen. Zur Schließung der
Lücke muss man sich daran erinnern, dass
\[|Y^{A\cup B}| = |Y^A\times Y^B| = |Y^A|\cdot |Y^B|,\;
\text{sofern}\; A\cap B=\emptyset.\]
Zusammen mit $|Y^{\{x\}}|=n$ wird der Induktionsschritt damit präzisiert zu
\[|Y^X| = |Y^{(X\setminus\{x\})\cup\{x\}}| = |Y^{X\setminus\{x\}}|\cdot |Y^{\{x\}}|
\stackrel{\mathrm{IV}}= n^{k-1}\cdot n = n^k.\]

\section{Zur elementaren Zahlentheorie}

\subsection{Kongruenzen}

Die elementare Zahlentheorie beschäftigt sich mit der Teilbarkeit von
Zahlen, den Resten, die bei der Ganzzahldivision übrig bleiben
und der Auffindung ganzzahliger Lösungen von Gleichungen. Lange Zeit
als ein Gebiet der reinen Mathematik angesehen, stellte sie sich
als bedeutsam für die praktische Informatik und die Kryptologie heraus. 

Die \emph{modulare Arithmetik}, auch \emph{Restklassenarithmetik}
genannt, ist ein wichtiges Hilfsmittel der elementaren Zahlentheorie.
Sie geht aus der gewöhnlichen Arithmetik hervor, indem zwei ganze
Zahlen als gleich angesehen werden, wenn sie bei Division durch eine
vorab gewählte feste Zahl denselben Rest lassen. Diese Form der
Gleichheit nennt man \emph{Kongruenz}, die feste Zahl den \emph{Modul}.
Kongruenz zweier Zahlen ist damit gleichbedeutend, dass die eine
Zahl zur anderen um ein Vielfaches des Moduls verschoben liegt, was nichts
anderes bedeutet, als dass der Modul die Differenz der beiden Zahlen teilt.

Man vermittelt das Konzept gern am Lauf von Uhrzeigern. Hierbei werden
die Stunden $0, 12, 24, 36$ usw. als gleich angesehen. Entsprechend
werden die Stunden $1, 13, 25, 37$ usw. als gleich angesehen. Es
verbleiben nur noch $12$ unterschiedliche Zeitpunkte, die vollen
Stunden $0, 1, \ldots, 11$. Ein Uhrzeiger auf $11$ Uhr trifft
zwei Stunden später auf $1$ Uhr. Anders ausgedrückt sind die Zahlen
$11+2$ und $1$ kongruent Modulo 12, was nach Carl Friedrich Gauß
in der Form
\[11 + 2\equiv 1\pmod{12}\]
notiert wird.

\begin{Definition}[Kongruenz]\index{Kongruenz}\newlinefirst
Zwei ganze Zahlen $a,b$ heißen kongruent modulo $m$, wenn ihre Differenz
$a-b$ durch $m$ teilbar ist,%
\[a\equiv b\pmod{m}\;:\bicond\;\exists k\in\Z\colon a-b=km.\]
\end{Definition}
Statt »$(\mathrm{mod}\;m)$« schreibt man beim Rechnen meist
kürzer »$(m)$«.

\begin{Satz}
Die Kongruenz ist eine Äquivalenzrelation, das heißt, es gilt
\begin{align*}
&a\equiv a\pmod{m},&&\text{(Reflexivität)}\\
&a\equiv b\cond b\equiv a\pmod{m},&&\text{(Symmetrie)}\\
&a\equiv b\land b\equiv c\cond a\equiv c\pmod{m}.&&\text{(Transitivität)}
\end{align*}
\end{Satz}
\begin{Beweis}
Für die Reflexivität ist ein $k$ mit $0=a-a=km$
zu finden. Setze $k=0$.

Bei der Symmetrie gibt es nach Voraussetzung
ein $k$ mit $a-b=km$. Dann ist $b-a=-km$. Setze $k'=-k$.
Es gibt also $k'$ mit $b-a=k'm$, somit gilt $b\equiv a$.

Bei der Transitivität gibt es nach Voraussetzung ein $k$ mit
$a-b=km$ und ein $l$ mit $b-c=lm$. Das heißt, es gilt
$km+lm = a-c$, \text{ergo} $a-c = (k+l)m$. Setze $k'=k+l$. Es gibt also
$k'$ mit $a-c=k'm$. Somit gilt $a\equiv c$.\;\qedsymbol
\end{Beweis}

\begin{Satz}\label{Kongruenz-add-sub}
Sind $a,b,c$ ganze Zahlen, dann gilt
\begin{align*}
a\equiv b\pmod{m}&\iff a+c\equiv b+c\pmod{m},\\
a\equiv b\pmod{m}&\iff a-c\equiv b-c\pmod{m}.
\end{align*}
\end{Satz}
\begin{Beweis}
Unter Beachtung von $(b+c)-(a+c)=b-a$ findet man
\begin{gather*}
a\equiv b\pmod{m}
\iff (\exists k\in\Z\colon b-a=km)\\
\iff (\exists k\in\Z\colon (b+c)-(a+c)=km)\\
\iff a+c\equiv b+c\pmod{m}.
\end{gather*}
Für die Subtraktion von $c$ ist die Überlegung analog.\;\qedsymbol
\end{Beweis}

%\newpage
\begin{Satz}\label{Kongruenz-mul}
Sind $a,b,c$ ganze Zahlen, dann gilt
\[a\equiv b\pmod{m} \implies ac\equiv bc\pmod{m}.\]
\end{Satz}
\begin{Beweis}
Unter der Voraussetzung $a\equiv b\pmod{m}$ gibt es ein
$k$ mit $b-a=km$. Es gilt
\[b-a=km\iff (b-a)c=kcm \iff bc-ac=k'm\]
mit $k':=kc$. Man hat also
\[(\exists k'\in\Z\colon bc-ac=k'm)\iff ac\equiv bc\pmod{m}.\;\qedsymbol\]
\end{Beweis}

\begin{Satz}
Gilt $a\equiv a'\pmod{m}$ und
$b\equiv b'\pmod{m}$, dann gilt auch
\begin{align*}
a+b&\equiv a'+b'\pmod{m},\\
a-b&\equiv a'-b'\pmod{m},\\
ab&\equiv a'b'\pmod{m}.
\end{align*}
\end{Satz}
\strong{Beweis.} Man findet
\[\left.\begin{aligned}
a\equiv a'&\implies a+b\equiv a'+b\\
b\equiv b'&\implies a'+b\equiv a'+b'
\end{aligned}\right\}
\implies a+b\equiv a'+b\equiv a'+b'\pmod{m}.\]
Für die Subtraktion ist die Überlegung analog. Für die Multiplikation
ebenfalls:%
\[\left.\begin{aligned}
a\equiv a'&\implies ab\equiv a'b\\
b\equiv b'&\implies a'b\equiv a'b'
\end{aligned}\right\}
\implies ab\equiv a'b\equiv a'b'\pmod{m}.\;\qedsymbol\]

\begin{Satz}
Addition des Moduls führt auf eine kongruente Zahl:%
\[a\equiv a+m\equiv a-m\pmod{m}.\]
\end{Satz}
\begin{Beweis}
Es gilt
\[a\equiv a+m\pmod{m}\iff (\exists k\in\Z\colon km=(a+m)-a=m).\]
Setze $k=1$. Bei
\[a\equiv a-m\pmod{m}\iff (\exists k\in\Z\colon km=(a-m)-a=-m)\]
setze $k=-1$.\;\qedsymbol
\end{Beweis}

\subsection{Teilbarkeit}

Die Beziehung »$m$ teilt $a$« ist definiert als
\[m\mid a\defiff (\exists k\in\Z\colon a=km) \iff a\equiv 0\pmod{m}.\]
Während die Kongruenz eine Äquivalenzrelation ist, stellt die
Teilbarkeit eine Halbordnung dar, dergestalt dass sie die drei Axiome
\begin{align*}
&\forall a\in\Z\colon a\mid a, &&\text{(Reflexivität)}\\
&\forall a,b\in\Z\colon a\mid b\land b\mid a\cond a = b,&&\text{(Antisymmetrie)}\\
&\forall a,b,c\in\Z\colon a\mid b\land b\mid c\cond a\mid c.&&\text{(Transitivität)}
\end{align*}
erfüllt. Teilt $a$ die Zahl $b$, ist $a$ also in gewisser Weise kleiner
als $b$.

\subsection{Restklassenringe}

Wir könnten nun beginnen, mit der modularen Arithmetik interessante
Probleme zu lösen. Zunächst möchte ich aber erläutern, wie die
modulare Arithmetik mit dem Restklassenring zusammenhängt. Unter
diesem Blickwinkel bekommen wir ein tieferes Verständnis und können
Mittel der Ringtheorie und Gruppentheorie anwenden.

Zu einer ganzen Zahl $a$ ist die Restklasse modulo $m$ definiert als 
\[[a]_m := \{x\mid x\equiv a\pmod m\}.\]
Eine alternative Schreibweise für $[a]_m$ ist $a+m\Z$. Weil die
Kongruenz eine Äquivalenzrelation ist, handelt es sich bei den
Restklassen um Äquivalenzklassen. Wir betrachten nun die
Quotientenmenge
\[\Z/m\Z := \{[a]_m\mid a\in\Z\}.\]
Nun können wir die Addition und Multiplikation von Restklassen
definieren.
\begin{Satz}
Auf $\Z/m\Z$ sind die beiden Operationen
\begin{align*}
[a]_m + [b]_m &:= [a+b]_m,\\
[a]_m\cdot [b]_m &:= [ab]_m
\end{align*}
wohldefiniert.
\end{Satz}
\strong{Beweis.} Zu zeigen ist, dass $a+b\equiv x+y$ gilt, sofern
$a\equiv x$ und $b\equiv y$ ist. Gemäß Satz \ref{Kongruenz-add-sub} gilt
\begin{align*}
a\equiv x &\iff a+b\equiv x+b,\\
b\equiv y &\iff x+b\equiv x+y.
\end{align*}
Aus den beiden Prämissen erhalten wir demzufolge
$a+b\equiv x+b\equiv x+y$.
Die Argumentation zur Multiplikation ist analog, wobei
Satz \ref{Kongruenz-mul} zur Anwendung kommt.\,\qedsymbol

Die Struktur $(\Z/m\Z,+,\cdot)$ nennt man den \emph{Restklassenring}
zum Modul $m$.

\begin{Satz}
Jeder Restklassenring ist ein kommutativer unitärer Ring.
\end{Satz}
\begin{Beweis}
Bereits bewiesen wurde, dass die ganzen Zahlen einen
kommutativen unitären Ring bilden. Aufgrund der Wohldefiniertheit der
Addition und Multiplikation ist Kongruenz modulo $m$ eine
Kongruenzrelation. Satz \ref{Kongruenz-Ring-Quotient} zeigt somit die
Behauptung.\,\qedsymbol
\end{Beweis}

\noindent
Zudem ist die Quotientenabbildung
\[\varphi\colon\Z \to \Z/m\Z,\quad \varphi(a):=[a]_m.\]
ein Eins-erhaltender Ringhomomorphismus, wie aus dem Beweis
von Satz \ref{Kongruenz-Ring-Quotient} hervorgeht.

%\newpage
\subsection{Euklidische Division}

\begin{Satz}[Lemma zur euklidischen Division]\mbox{}\\*
Zu je zwei ganzen Zahlen $a,b$ mit $b\ne 0$ gibt es zwei eindeutig
bestimmte Zahlen $q,r$ mit $0\le r<|b|$, so dass $a=bq+r$. Man nennt
$q$ den \emph{Quotient} und $r$ den \emph{Rest}.
\end{Satz}
\begin{Beweis}[Beweis der Existenz]
Betrachten wir zunächst den Fall, bei dem $a\ge 0$
und $b>0$ ist. Man kann dann sooft $b$ von $a$ abziehen, bis sich eine
Zahl $\ge 0$ und $<b$ ergibt. Formal bilden wir die
Folge $r_k := a-bk$. Nun muss für irgendein $k\ge 0$ schließlich
$0\le r_k<b$ sein. Damit ist $q=k$ und $r=r_k$ gefunden.

Sei nun $a<0$. Wie bereits gezeigt gibt es $q,r$ mit $-a = bq+r$.
Für $r=0$ haben wir dann mit $q':=-q$ und $r':=0$ einen Quotient
und einen Rest. Sei nun $r\ne 0$. Dann gilt
\[a = -bq-r = -(q+1)b + b - r.\]
Mit $q':=-(q+1)$ und $r':=b-r$ gibt es somit auch in diesem Fall einen
Quotient und einen Rest. Der Rest erfüllt auch die gewünschte
Ungleichung, denn aus $r<b$ ergibt sich $0<r'$ und aus $0<r$ ergibt
sich $r'<b$.

Sei nun $a$ beliebig und $b<0$. Dann gibt es $q,r$ mit $a=(-b)q+r$.
Setze also $q':=-q$ und $r':=r$. Damit gilt $a=bq'+r'$, womit
auch in diesem Fall ein Quotient und ein Rest gefunden ist.\,\qedsymbol
\end{Beweis}
\begin{Beweis}[Beweis der Eindeutigkeit]
Das Paar $q,r$ erfülle
$a=bq+r$ und $q',r'$ erfülle ebenfalls $a=bq'+r'$. Dann gilt
\[bq+r = bq'+r',\iff b(q-q') = r'-r,\implies |b| |q-q'| = |r'-r|.\]
Aus $0<r<|b|$ und $0<r'<|b|$ erhält man außerdem $|r'-r|<|b|$. Somit
muss $|b| |q-q'| < |b|$ sein, also $|q-q'| < 1$. Eine nichtnegative
ganze Zahl kann aber nur dann kleiner als eins sein, wenn sie null
ist. Damit hat man
\[|q-q'|=0,\iff q-q'=0,\iff q=q'.\]
Entsprechend folgt $r=r'$.\,\qedsymbol
\end{Beweis}

Bei der euklidischen Division $a:b$ ist $(a\bmod b)$ eine geläufige
Schreibweise für den Rest. Das Lemma zur euklidischen Division sagt uns,
dass jede Restklasse von $\Z/m\Z$ einen kanonischen Repräsentant
besitzt. Nämlich besitzt die Restklasse $[a]_m$ den kanonischen
Repräsentant $r=(a\bmod m)$, denn $a=mq+r$ bedeutet dass $a-r$ durch
$m$ teilbar ist, also
\[a\equiv r\pmod m,\quad\text{bzw.}\quad [a]_m = [r]_m.\]


\subsection{Rundung}

Manche Formeln oder Rechnungen verlangen das Runden von Zahlen auf
eine ganze Zahl oder auf eine bestimmte Zahl von Nachkommastellen.
Im Algorithmus, der Vektorgrafiken als Rastergrafik aus Pixeln rendern
soll, muss zum Beispiel zwangsläufig an irgendeiner Stelle gerundet
werden. Außerdem ist Runden in vielen händischen Rechnungen
allgegenwärtig, sei es in wissenschaftlichen, technischen, gewerblichen
oder alltäglichen.

Als nächstes will ich daher erklären, wie Funktionen zum Runden formal
definiert werden, und welche Eigenschaften sie besitzen. Als wichtige
Grundfunktionen treten dabei die \emph{Abrundungsfunktion}, engl.
\emph{floor}, und die \emph{Aufrundungsfunktion}, engl. \emph{ceil},
auf, die vereinzelt auch in der Kombinatorik und in der Zahlentheorie
vorkommen.

\begin{Definition}[Floor]\label{def:floor}\index{Floorfunktion}
Für $x\in\R$ definiert man
\[y = \lfloor x\rfloor\,:\Leftrightarrow\, y\in\Z\land 0\le x-y < 1.\]
\end{Definition}

\begin{Definition}[Ceil]\label{def:ceil}\index{Ceilfunktion}
Für $x\in\R$ definiert man
\[y = \lceil x\rceil\,:\Leftrightarrow\, y\in\Z\land 0\le y-x < 1.\]
\end{Definition}

\begin{Satz}\label{floor-add-int}
Für jede ganze Zahl $k$ gilt $\lfloor k + x\rfloor = k + \lfloor x\rfloor$.
\end{Satz}
\begin{Beweis} Aufgrund der Prämisse $k\in\Z$ ist $y\in\Z$ äquivalent
zu $y-k\in\Z$. Unter dieser Gegebenheit findet sich mit
Def. \ref{def:floor} die äquivalente Umformung
\begin{align*}
y = \lfloor k+x\rfloor &\iff y\in\Z\land 0\le (k+x)-y < 1\\
&\iff y-k\in\Z\land 0\le x-(y-k) < 1\\
&\iff y - k = \lfloor x\rfloor \iff y = k + \lfloor x\rfloor.\,\qedsymbol
\end{align*}
\end{Beweis}

\begin{Satz}\label{floor-is-zero}
Für $0\le x < 1$ gilt $\lfloor x\rfloor = 0$.
\end{Satz}
\begin{Beweis}
Dies folgt unmittelbar aus Def. \ref{def:floor}.\,\qedsymbol
\end{Beweis}

\begin{Satz}
Es gilt
\[\begin{array}{@{}l@{\,}c@{\,}c@{}l@{\,}c@{\,}c@{}l}
\lfloor x\rfloor &=& \max &\{k\in\Z\mid k\le x\} &=& \min &\{k\in\Z\mid x < k + 1\},\\[2pt]
\lceil x\rceil &=& \min &\{k\in\Z\mid x\le k\} &=& \max &\{k\in\Z\mid k - 1 < x\}.
\end{array}\]
\end{Satz}
\begin{Beweis}
Mit bezüglich $y=\lfloor x\rfloor$ entfalteten Def. \ref{def:floor},
\ref{def:max-min} lautet die Aussage
\[y\in\Z\land 0\le x-y<1\iff
y\in\Z\land y\le x\land (\forall k\in\Z\colon k\le x\Rightarrow k\le y).\]
Für die Implikation von links nach rechts ist im Wesentlichen
\[k\in\Z, y\in\Z, k\le x, x < y+1\vdash k\le y\]
zu zeigen. Man erhält zunächst $k < y+1$ per Transitivgesetz.
Wegen $k,y\in\Z$ folgt daraus $k\le y$. Für die Implikation von rechts
nach links ist im Wesentlichen
\[y\in\Z,(\forall k\in\Z\colon k\le x\Rightarrow k\le y)\vdash x < y + 1\]
zu zeigen. Angenommen, es wäre $y+1\le x$. Die Allaussage wird mit
$k:=y+1$ spezialisiert. Per Modus ponens erhält man den Widerspruch
$y+1\le y$. Ergo muss die Annahme falsch sein, was äquivalent zu
$x<y+1$ ist. Die restlichen Beweise verlaufen analog.\,\qedsymbol
\end{Beweis}

\begin{Satz}\label{floor-div-floor}
Für $x\in\R$ und $n\in\Z_{\ge 1}$ gilt
\[\left\lfloor\frac{x}{n}\right\rfloor
= \left\lfloor\frac{\lfloor x\rfloor}{n}\right\rfloor.\]
\end{Satz}
\begin{Beweis}
Sei $a:=x-\lfloor x\rfloor$. Per Def. \ref{def:floor} gilt $0 \le a < 1$.
Laut dem Lemma zur euklidischen Division gilt außerdem
$\lfloor x\rfloor = qn+r$
mit $q=\lfloor\frac{\lfloor x\rfloor}{n}\rfloor$
und $0\le r\le n - 1$. Infolge gilt $0\le a+r<n$, also
$0\le\frac{a+r}{n}<1$ und somit $\lfloor\frac{a+r}{n}\rfloor = 0$
laut Satz \ref{floor-is-zero}.
Es findet sich%
\[\left\lfloor\frac{x}{n}\right\rfloor
= \left\lfloor\frac{a+\lfloor x\rfloor}{n}\right\rfloor
= \left\lfloor\frac{a+qn+r}{n}\right\rfloor
\stackrel{\text{(1)}}= q + \left\lfloor\frac{a+r}{n}\right\rfloor
\stackrel{\text{(2)}}= q =
\left\lfloor\frac{\lfloor x\rfloor}{n}\right\rfloor,\]
wobei (1) laut Satz \ref{floor-add-int} gilt, und (2)
wie soeben ausgeführt.\,\qedsymbol
\end{Beweis}

\begin{Satz}
Für $x\in\R$ und $m,n\in\Z$ mit $n\ge 1$ gilt
\[\left\lfloor\frac{m+x}{n}\right\rfloor
= \left\lfloor\frac{m+\lfloor x\rfloor}{n}\right\rfloor.\]
\end{Satz}
\begin{Beweis}
Laut Satz \ref{floor-add-int} ist $m+\lfloor x\rfloor = \lfloor m+x\rfloor$.
Die Aussage folgt nun als Korollar aus Satz \ref{floor-div-floor}.\,\qedsymbol
\end{Beweis}

\noindent
Das Runden einer Zahl auf eine ganze Zahl kann vermittels floor als
\[\mathrm{round}(x) := \mathrm{sgn}(x)\lfloor |x|+\tfrac{1}{2}\rfloor\]
beschrieben werden. Das Runden auf die $n$-te Nachkommastelle daraufhin als
\[\mathrm{round}(x, n) := \frac{\mathrm{round}(10^n x)}{10^n}.\]


\chapter{Elemente der Stochastik}

\section{Grundbegriffe}

\subsection{Ereignisse}

Die Wahrscheinlichkeitstheorie beschäftigt sich mit
\emph{Zufallsexperimenten}. Darunter versteht man ein Experiment
mit zufälligem Ausgang, das, um der Wissenschaftlichkeit genüge zu tun,
unter genau definierten Versuchsbedingungen durchgeführt wird.
Der Ausgang führt immer zu einem \emph{Ergebnis}. Alle erreichbaren
Ergebnisse fasst man zur \emph{Ergebnismenge} zusammen. Allgemeiner
genügt es, wenn jedes Ergebnis in der Ergebnismenge liegt, wobei diese
aber auch Elemente enthalten darf, die das Experiment niemals abwirft.
Jede Teilmenge der Ergebnismenge nennt man ein \emph{Ereignis}\index{Ereignis}.
Die Potenzmenge der Ergebnismenge heißt \emph{Ereignisraum}, sie
besteht aus allen denkbaren Ereignissen. Man sagt, ein Ereignis sei
eingetreten, wenn das Ergebnis des Versuchs im diesem Ereignis liegt.

Zu beachten ist, dass wir dabei eine endliche oder höchstens
abzählbar unendliche Ergebnismenge voraussetzen. Bei überabzählbaren
Ergebnismengen kommt es zu Unwägbarkeiten, deren Klärung Gegenstand der
Maßtheorie ist.

Ein schlichtes Experiment bietet der Wurf des Spielwüfels, ein mit
Augenzahlen beschriftetes regelmäßiges Hexaeder. Die Ergebnismenge
wird als
\[\Omega := \{1, 2, 3, 4, 5, 6\}\]
festgelegt. Betrachten wir die drei Ereignisse
\[A := \{2\},\quad B := \{1,2\},\quad C:=\{1,3\}.\]
Ist $\omega=2$ das Ergebnis des Versuchs, sind die Ereignisse $A,B$
eingetreten. Zwei Ereignisse, die niemals gleichzeitig eintreten,
heißen \emph{disjunkt}. So sind die $A,C$ disjunkt, weil ihre
Schnittmenge leer ist.

\subsection{Wahrscheinlichkeiten}

Man kann nicht voraussagen, wie ein Experiment ausgehen wird. Das
Wahrscheinlichkeitsmaß liefert dennoch ein Maß dafür, wie sicher der
Eintritt eines Ereignisses ist. Wahrscheinlichkeit wird tiefergründig
verständlich, wenn dasselbe Zufallsexperiment abermals wiederholt wird.
Wir zählen, wie häufig ein Elementarereignis eingetreten ist.

Es sei ein Versuch $n$ mal durchgeführt worden, was zu den Ergebnissen
$a_i$ für $i=1$ bis $i=n$ geführt hat. Wir definieren die \emph{relative
Häufigkeit} des Ereignisses $A$ als die Zahl
\[r_{n,a}(A) := \frac{1}{n}|\{i\in\{1,\ldots,n\}\mid a_i\in A\}|.\]
Relative Häufigkeiten bieten bei hinreichend großem $n$ eine Näherung
für die Wahrscheinlichkeit. Zur Vermessung eines Würfels wird man
diesen also möglichst oft werfen wollen. Man erhält so die relativen
Häufigkeiten der Elementarereignisse, und damit näherungsweise auch
ihre Wahrscheinlichkeiten. So lässt sich feststellen, ob ein
Würfel gezinkt wurde.

Fassen wir $a$ als Funktion $i\mapsto a_i$ auf, können wir schreiben
\[\{i\mid a_i\in A\} = \{i\mid i\in a^{-1}(A)\} = a^{-1}(A).\]
Für disjunkte Ereignisse $A,B$ erhält man nun
\begin{align*}
r_{n,a}(A\cup B) &= \tfrac{1}{n}|a^{-1}(A\cup B)|
= \tfrac{1}{n}|a^{-1}(A)\cup a^{-1}(B)|\\
&= \tfrac{1}{n}|a^{-1}(A)| + \tfrac{1}{n}|a^{-1}(B)|
= r_{n,a}(A) + r_{n,a}(B).
\end{align*}

\subsection{Zufallsgrößen}

Eine \emph{Zufallsgröße}\index{Zufallsgroesse@Zufallsgröße} darf man
sich als eine Funktion $X\colon\Omega\to\Omega'$
vorstellen, die eine kausale Verbindung zwischen den Ergebnismengen
$\Omega,\Omega'$ schafft. Ein Ergebnis $\omega\in\Omega$ führt
zu $X(\omega)$. Ursächlich für ein $x\in\Omega'$ sind daher all die
$\omega$ mit $x=X(\omega)$. Das heißt, ursächlich für das
Elementarereignis $\{x\}$ ist dessen Urbild $X^{-1}(\{x\})$.
Infolge muss die Wahrscheinlichkeit von $\{x\}$ die es Urbildes sein.
Insofern definiert man auf $\Omega'$ das Wahrscheinlichkeitsmaß%
\[P_X\colon\mathcal P(\Omega')\to [0,1],\quad P_X(A) := P(X^{-1}(A)).\]
Man nennt $P_X$ die \emph{Verteilung}\index{Verteilung} von $X$. Mit der
identischen Zufallsgröße
\[\id\colon\Omega\to\Omega,\quad \id(\omega) := \omega\]
versteht sich auch das ursprüngliche Maß $P$ als die Verteilung $P=P_{\id}$.

Geläufig sind die Schreibweisen
\begin{align*}
P(X=x) &:= P(X^{-1}(\{x\})), & \{X=x\} &:= X^{-1}(\{x\}),\\
P(X\in A) &:= P(X^{-1}(A)), & \{X\in A\} &:= X^{-1}(A).
\end{align*}
Es ist $\{X=x\}$ dasselbe wie $\{X\in\{x\}\}$. 
Ist $P$ die Gleichverteilung auf $\Omega$, ergibt sich
\[P(X\in A) = \frac{|\{X\in A\}|}{|\Omega|}.\]
Standardbeispiel. Wir werfen zwei Spielwürfel.
Die Ergebnismenge sei%
\[\Omega := \{1,\ldots,6\}\times\{1,\ldots,6\},\]
und jedes der 36 elementaren Ereignisse sei gleich wahrscheinlich, habe
also die Wahrscheinlichkeit $\tfrac{1}{36}$.
Es bezeichne $\omega_1$ das Ergebnis des ersten, und $\omega_2$
das des zweiten Wurfs. Wir betrachten die Zufallsgröße%
\[X\colon\Omega\to\{2,\ldots,12\},\quad
X(\omega_1,\omega_2) := \omega_1+\omega_2.\]
Gesucht sei $P(X=4)$. Man ermittelt
\[\{X=4\} = \{(1,3), (2,2), (3,1)\},\quad\text{ergo}\;P(X=4) = \tfrac{3}{36}.\]
Allgemein zerfällt ein Ereignis $A$ ja in seine disjunkten Elementarereignisse
$\{x\}$, so dass $A = \bigcup_{x\in A} \{x\}$ gilt. Weil nun die
Fasern $X^{-1}(\{x\})$ ebenfalls disjunkt sind, muss $P(X\in A)$ die
Summe der $P(X=x)$ mit $x\in A$ sein. Das heißt, man rechnet%
\[P(X\in A) = P(X^{-1}(\bigcup_{x\in A} \{x\})) = P(\bigcup_{x\in A} X^{-1}(\{x\}))
= \sum_{x\in A} P(X=x).\]
Die Verteilung $P_X$ ist demzufolge bereits eindeutig bestimmt,
sobald $P(X=x)$ für jedes $x\in\Omega'$ vorliegt. Dies motiviert
uns, die Funktion
\[p_X\colon\Omega\to[0,1],\quad p_X(x) := P(X=x)\]
zu definieren, genannt die \emph{Wahrscheinlichkeitsfunktion} der
Zufallsgröße $X$.

\newpage
\section{Mehrstufige Experimente}

\subsection{Bedingte Wahrscheinlichkeiten}

Es findet ein zweistufiges Experiment statt, welches sich sich
aus einem ersten und einem zweiten Wurf eines Spielwürfels
zusammensetzt. Bei jedem der Würfe bestehe eine Gleichverteilung.
Zur Frage steht, wie wahrscheinlich das Ereignis $\{(6,6)\}$ ist.
Ein Paar $(\omega_1,\omega_2)$ fasse hierbei das Ergebnis $\omega_1$
des ersten und $\omega_2$ des zweiten Wurfs zusammen.

Die Wahrscheinlichkeit der ersten Sechs beträgt $\tfrac{1}{6}$,
die der zweiten ebenfalls $\tfrac{1}{6}$. Sie multiplizieren sich
zu zu $\tfrac{1}{36}$, richtig?

Es wäre doch möglich, dass zwischen den beiden Würfen eine,
sagen wir, geisterhafte Beziehung besteht, dergestalt dass
der zweite Wurf niemals in einer Sechs resultiert, sofern das
Ergebnis des ersten eine war. Trotzdem sind die Wahrscheinlichkeiten bei
jedem der Würfe für sich allein gesehen gleichverteilt. Dafür muss man
nicht unbedingt die Wirklichkeit manipulieren. Das Phänomen ist bereits
bei der Erzeugung von Zufallszahlen im Computer beobachtbar. War die
erste Zufallszahl eine Sechs, braucht der Generator die zweite lediglich
solange zu verwerfen, wie sie eine Sechs sein sollte. Umstände dieser
Art stellen nicht nur ein Gedankenspiel dar, so dass wir uns notgedrungen
mit ihnen auseinandersetzen müssen. Sie führen zum Begriff der
\emph{bedingten Wahrscheinlichkeit}.

Bisher wurde immer nur die Verteilung der Wahrscheinlichkeiten eines
Würfels für sich allein betrachtet. Das war modelliert durch die Größe
\[X_0\colon\Omega\to\Omega,\quad X_0(\omega):=\omega,\quad \Omega := \{1,\ldots,6\},\]
mit der Gleichverteilung $P_0$, so dass $P_0(X=6)=\tfrac{1}{6}$.

Wir modellieren das zweistufige Experiment durch die Zufallsgröße
\[X\colon\Omega^2\to\Omega^2,\quad X(\omega) := (X_1,X_2)(\omega)
= (X_1(\omega), X_2(\omega)),\]
die sich mit $\omega = (\omega_1,\omega_2)$ aus den zwei Zufallsgrößen
\begin{gather*}
X_1\colon\Omega^2\to\Omega,\quad X_1(\omega_1,\omega_2) := \omega_1,\\
X_2\colon\Omega^2\to\Omega,\quad X_2(\omega_1,\omega_2) := \omega_2
\end{gather*}
zusammensetzt. Es stellt $X_1(\omega)$ das Ergebnis des ersten und
$X_2(\omega)$ das des zweiten Wurfs dar. Wie gewünscht gilt
\[(X_1(\omega),X_2(\omega)) = (X_0(\omega_1),X_0(\omega_2)) = (\omega_1,\omega_2).\]
Es bezeichne $P$ die Verteilung auf $\Omega^2$. Wir wissen hier allerdings
lediglich
\begin{gather*}
P(X_1=\omega_1) = P_0(X_0=\omega_1) = \tfrac{1}{6},\\
P(X_2=\omega_2) = P_0(X_0=\omega_2) = \tfrac{1}{6}.
\end{gather*}
Die Fehlannahme besteht nun darin, dass per se
\[P(\{X_1=\omega_1\}\cap\{X_2=\omega_2\}) = P(X_1=\omega_1)P(X_2=\omega_2)\]
gelten müsse. Ist diese Gleichung erfüllt, nennt man die
Zufallsgrößen $X_1,X_2$ \emph{unabhängig}. In der bisherigen Sichtweise,
wo wir nur $X_0$ mit $P_0$ gesehen haben, war es uns nicht möglich,
stochastische Abhängigkeit zu beschreiben. Man notiert allgemein%
\[P(X=x,Y=y) := P(\{X=x\}\cap\{X=y\}) = P(X=x)P(Y=y\mid X=x).\]
Der letzte Faktor bezeichne hierbei die bedingte Wahrscheinlichkeit
für das Ereignis $\{Y=y\}$, unter der Bedingung, dass $\{X=x\}$
bereits eingetreten ist.
\begin{Definition}[Bedingte Wahrscheinlichkeit]\newlinefirst
Die bedingte Wahrscheinlichkeit für den Eintritt von $A$ unter der
Bedingung $B$ ist für $P(B)\ne 0$ definiert gemäß
\[P(A\mid B) := \frac{P(A\cap B)}{P(B)}.\]
\end{Definition}
Wir setzen speziell $B:=\{X=x\}$ und $A:=\{Y=y\}$ ein,
das macht
\[P(Y=y\mid X=x) = \frac{P(X=x,Y=y)}{P(X=x)}.\]
Sind $X,Y$ unabhängig, gilt also
\[P(Y=y\mid X=x) = P(Y=y).\]
Mit der geisterhaften Beziehung zwischen den Würfeln wäre allerdings
\[0 = P(X_2=6\mid X_1=6) \ne P(X_1=6) = \tfrac{1}{6}.\]

\documentclass[8pt,fleqn,aspectratio=169]{beamer}
\usetheme{Antibes}
\useinnertheme{rectangles}
\useoutertheme{infolines}
\usepackage[utf8]{inputenc}
\usepackage[T1]{fontenc}
\usepackage[ngerman]{babel}

% Patch the look of +, = in arev
\usefonttheme{serif}

\usepackage{arev}
% Patch punctuation to be upright
\DeclareMathSymbol{.}{\mathpunct}{operators}{`.}
\DeclareMathSymbol{,}{\mathpunct}{operators}{`,}

\usepackage{amsmath}
\usepackage{amssymb}
\usepackage{booktabs}

\setbeamertemplate{footline}{%
\begin{beamercolorbox}[ht=3.0ex,dp=1ex]{title in head/foot}
\hfill\footnotesize\insertpagenumber\enspace\enspace\end{beamercolorbox}}

% \documentclass[8pt,aspectratio=169]{beamer}
\setbeamersize{text margin left=6em}
\setbeamersize{text margin right=6em}

\definecolor{bluegreen1}{rgb}{0.0,0.20,0.28}
\definecolor{bluegreen2}{rgb}{0.0,0.20,0.28}
\setbeamercolor*{palette primary}{fg=white,bg=bluegreen1}
\setbeamercolor*{palette secondary}{fg=white,bg=bluegreen2}
\setbeamercolor*{palette tertiary}{fg=white,bg=bluegreen2}
\setbeamercolor{itemize item}{fg=black}
\setbeamercolor{block title}{bg=bluegreen2}
\newcommand{\modest}[1]{{\small\color{gray}#1}}
\hypersetup{colorlinks,urlcolor=magenta}

\usepackage{listings}
\lstset{basicstyle=\ttfamily}
\lstdefinelanguage{IMP}{sensitive=true, keywords={
  true, false, skip, if, then, else, end, while, do, not, and, or}}
\lstdefinelanguage{Coq}{sensitive=true, keywords={
  Require, Import, Definition, Inductive, Fixpoint, Theorem, Proof, Qed,
  Goal, Module, End, Open, Scope,
  forall, fun, match, with, end, if, then, else}}

\newcommand{\inferrulewidth}{0.4608pt}
\usepackage{ebproof}
\ebproofset{rule margin = 0.5ex}
\ebproofset{label separation = 0.3em}
\ebproofset{right label template={\scriptsize\inserttext}}
\ebproofset{rule thickness=\inferrulewidth}

\newcommand{\ttq}{\textquotesingle}
\newcommand{\unit}[1]{\mathrm{#1}}
\newcommand{\strong}[1]{\textsf{\textbf{#1}}}
\newcommand{\defiff}{\quad:\Longleftrightarrow\quad}
\newcommand{\infernote}[1]{\!\text{\footnotesize #1}}
\renewcommand{\qedsymbol}{\ensuremath{\Box}}
\newcommand{\discharge}[1]{$\sim$#1}
\newcommand{\centerheadline}[1]{%
  \begin{center}\strong{#1}\end{center}}
\newcommand{\parspace}{\vspace{0.8em}}
\newcommand{\cond}{\rightarrow}

\newcommand{\kw}[1]{\textbf{\texttt{#1}}}
\newcommand{\code}[1]{{\texttt{#1}}}
\newcommand{\qb}[1]{[\!\![#1]\!\!]}
\newcommand{\Bool}{\mathrm{Bool}}
\newcommand{\Int}{\mathrm{Int}}
\newcommand{\Loc}{\mathrm{Loc}}
\newcommand{\Aexp}{\mathrm{Aexp}}
\newcommand{\Bexp}{\mathrm{Bexp}}
\newcommand{\Com}{\mathrm{Com}}
\newcommand{\evA}{\mathcal A}
\newcommand{\evB}{\mathcal B}
\newcommand{\evC}{\mathcal C}

\newcommand{\N}{\mathbb N}
\newcommand{\Z}{\mathbb Z}
\newcommand{\Abb}{\mathrm{Abb}}

\title{Programmverifikation}
% \subtitle{Untertitel}
\date{}

\begin{document}

\begin{frame}
\maketitle
\end{frame}

\begin{frame}
\centerheadline{Die Sprache IMP}
\end{frame}

\begin{frame}
Grundlegende Aspekte der Semantik von Programmiersprachen klären wir
anhand von IMP, einer minimalistischen imperativen Programmiersprache,
die uns als Studienobjekt dienen wird.
\end{frame}

\begin{frame}[fragile]

Beispiel für ein Programm in IMP:
\begin{lstlisting}[language=IMP, xleftmargin=\mathindent]
y := 1; k := n;
while not k = 0 do
   y := y*x;
   k := k - 1
end
\end{lstlisting}\pause
Es berechnet zu ganzen Zahlen $\code{x},\code{n}$ mit $\code{n}\ge 0$
die Potenz $\code{y} = \code{x}^\code{n}$. Klar ersichtlich.\pause

\parspace
Das \emph{schnelle Exponentiation} genannte Programm
\begin{lstlisting}[language=IMP, xleftmargin=\mathindent]
y := 1; a := x; k := n;
while 2 <= k do
   q := k/2; r := k - 2*q;
   if r = 1 then y := y*a else skip end;
   k := q; a := a*a
end;
if k = 1 then y := y*a else skip end
\end{lstlisting}
berechnet die Potenz ebenfalls. Immer noch klar ersichtlich?

\parspace
Wir würden dies gern \emph{beweisen}.
\end{frame}

\begin{frame}
Zunächst wird eine formale Definition der Sprache IMP unternommen.\pause

\parspace
Es sei $\Int$ die Menge der ganzen Zahlen im Dezimalsystem. Während also
$\Z$ die Menge der ganzen Zahlen bezeichnet, besteht $\Int$ aus konkreten
syntaktischen Darstellungen der Zahlen aus $\Z$.\pause

\parspace
Entsprechend sei $\Bool:=\{\kw{false},\kw{true}\}$, wobei die
Symbole \kw{false} und \kw{true} syntaktische Darstellungen der
Wahrheitswerte seien.
\end{frame}

\begin{frame}
Ein \strong{arithmetischer Ausdruck} $a$, kurz \strong{Term}, sei durch die
folgenden Produktionsregeln festgelegt.
\begin{itemize}
\item Eine ganze Zahl aus $\Int$ ist ein Term.
\item Eine Variable aus $\Loc$ ist ein Term.
\item Sind $a,a'$ Terme, so sind auch $a+a'$, $a-a'$, $a*a'$, $a\,/\,a'$ Terme.
\item Nichts anderes ist ein Term.
\end{itemize}
{\footnotesize Wir fassen die Terme hierbei als abstrakte Syntaxbäume auf.
Auf die genaue Grammatik und die Konstruktion eines Parsers will ich
an dieser Stelle nicht näher eingehen, damit der Fokus auf die
eigentliche Problemstellung nicht verloren geht. Die Operationen sollen
auf die gewöhnliche Art geschrieben werden, die Operatoren dabei die
gewöhnliche Rangfolge besitzen.}
\end{frame}

\begin{frame}
Ein \strong{boolescher Ausdruck} $b$, kurz \strong{Ausdruck}, sei durch
die folgenden Produktionsregeln festgelegt.
\begin{itemize}
\item Die Symbole \kw{false} und \kw{true} sind Ausdrücke.
\item Sind $a,a'$ Terme, so sind $a\;\code{=}\;a'$ und $a \;\code{<=}\; a'$ Ausdrücke.
\item Ist $b$ ein Ausdruck, so ist auch $\kw{not}\;b$ ein Ausdruck.
\item Sind $b,b'$ Ausdrücke, so sind auch $b\;\kw{and}\;b$ und $b\;\kw{or}\;b$ Ausdrücke.
\item Nichts anderes ist ein Ausdruck.
\end{itemize}
\end{frame}

\begin{frame}
Ein \strong{Kommando} $c$, auch \strong{Programm} genannt, sei durch die
folgenden Produktionsregeln festgelegt.
\begin{itemize}
\item Das Symbol \kw{skip} ist ein Kommando.
\item Ist $X$ eine Variable aus $\mathrm{Loc}$ und $a$ ein Term, so ist $X:=a$ ein Kommando.
\item Sind $c,c'$ Kommandos, so ist auch $c$; $c'$ ein Kommando.
\item Ist $b$ ein Ausdruck und sind $c,c'$ Kommandos, so ist auch\\
  \kw{if} $b$ \kw{then} $c$ \kw{else} $c'$ \kw{end} ein Kommando.
\item Ist $b$ ein Ausdruck und $c$ ein Kommando, so ist auch\\
  \kw{while} $b$ \kw{do} $c$ \kw{end} ein Kommando.
\item Nichts anderes ist ein Kommando.
\end{itemize}
\end{frame}

\begin{frame}
\centerheadline{Denotationelle Semantik}
\end{frame}

\begin{frame}
Wir bezeichnen
\begin{itemize}
\item mit $\Aexp$ die Menge der arithmetischen Ausdrücke,
\item mit $\Bexp$ die Menge der booleschen Ausdrücke,
\item mit $\Com$ die Menge der Kommandos.
\end{itemize}\pause
Will man nun einen arithmetischen Ausdruck auswerten, denkt
man sich dafür eine Funktion $\evA\colon\Aexp\to\Int$. Es ist
\[\evA\qb{1+2} = 3\]
usw. Nun kann ein arithmetischer Ausdruck aber auch Variablen enthalten,
womit bspw. auch der Wert $\evA\qb{\code{x}+1}$ bezüglich $\code{x}\in\Loc$
bestimmt sein muss.
\end{frame}

\begin{frame}
Die Auswertungsfunktion $\evA$ wird daher parametrisiert durch den aktuellen
Zustand $s$. Wir haben also $\evA\colon\Aexp\to (S\to\Int)$
bzw. $\evA\colon\Aexp\times S\to\Int$.\pause

\parspace
Einen \strong{Zustand} $s$ modellieren wir schlicht als die Belegung
der verfügbaren Variablen, also als eine Funktion $s\in S$ mit
$S:=\Abb(\Loc,\Int)$.\pause

\parspace
Zum Beispiel
\[s(X) := \begin{cases}
1, & \text{wenn $X=\code{x}$},\\
2, & \text{wenn $X=\code{y}$},\\
0 & \text{sonst}. 
\end{cases}\]
\end{frame}

\begin{frame}
Wir legen $\evA\colon\Aexp\to (S\to\Int)$ rekursiv fest gemäß
\begin{gather*}
\evA\qb{n}(s) := n,\\
\evA\qb{X}(s) := s(X),\\
\evA\qb{a+a'}(s) := \evA\qb{a}(s) + \evA\qb{a'}(s),\\
\evA\qb{a-a'}(s) := \evA\qb{a}(s) - \evA\qb{a'}(s),\\
\evA\qb{a*a'}(s) := \evA\qb{a}(s) \cdot \evA\qb{a'}(s),\\
\evA\qb{a\,/\,a'}(s) := \evA\qb{a}(s)\,/\, \evA\qb{a'}(s).
\end{gather*}
bezüglich $n\in\Int$, $X\in\Loc$ und $a,a'\in\Aexp$.
Die Operationen auf der rechten Seite werden hierbei auf die übliche
Art und Weise berechnet. Wir verwenden euklidische Ganzzahldivision
und setzen $n\,/\,0:=0$ für jedes $n\in\Int$.
\end{frame}

\begin{frame}
Wir legen $\evB\colon\Bexp\to (S\to\Bool)$ rekursiv fest gemäß
\begin{gather*}
\evB\qb{\kw{false}}(s) := \kw{false},\\
\evB\qb{\kw{true}}(s) := \kw{true},\\
\evB\qb{a \;\code{=}\; a'}(s) := (\evA\qb{a}(s) = \evA\qb{a'}(s)),\\
\evB\qb{a \;\code{<=}\; a'}(s) := (\evA\qb{a}(s)\le \evA\qb{a'}(s)),\\
\evB\qb{\kw{not}\;b} := \lnot\evB\qb{b}(s),\\
\evB\qb{b\;\kw{and}\;b'} := \evB\qb{b}(s)\land\evB\qb{b'}(s),\\
\evB\qb{b\;\kw{or}\;b'} := \evB\qb{b}(s)\lor\evB\qb{b'}(s),
\end{gather*}
bezüglich $a,a'\in\Aexp$ und $b,b'\in\Bexp$. Der Wahrheitswert der
Relationen bzw. Verknüpfungen auf der rechten Seite wird hierbei auf die
übliche Art und Weise berechnet.
\end{frame}

\begin{frame}
Wir legen $\evC\colon\Com\to (S\rightharpoonup S)$ rekursiv fest gemäß
\begin{gather*}
\evC\qb{\kw{skip}}(s) := s,\\
\evC\qb{X:=a}(s) := s[X:=A\qb{a}(s)],\\
\evC\qb{c; c'}(s) := \evC\qb{c'}(\evC\qb{c}(s)) = (\evC\qb{c'}\circ \evC\qb{c})(s),\\
\evC\qb{\text{\kw{if} $b$ \kw{then} $c$ \kw{else} $c'$ \kw{end}}}(s) := \begin{cases}
\evC\qb{c}(s),& \text{wenn $\evB\qb{b}(s)=\kw{true}$},\\
\evC\qb{c'}(s)& \text{sonst}.
\end{cases}\\
\evC\qb{\text{\kw{while} $b$ \kw{do} $c$ \kw{end}}}(s) := \begin{cases}
\varphi_{b,c}(\evC\qb{c}(s)),& \text{wenn $\evB\qb{b}(s)=\kw{true}$},\\
s & \text{sonst}
\end{cases}
\end{gather*}
mit $\varphi_{b,c}(s) := \evC\qb{\text{\kw{while} $b$ \kw{do} $c$ \kw{end}}}(s)$.\pause

\parspace
Bei $\evC\qb{c}$ handelt es sich um eine partielle Funktion, da die rekursive Auswertung
der while"=Schleife unter Umständen nicht terminiert -- zum Beispiel bei
\[\texttt{\kw{while} \kw{true} \kw{do} \kw{skip} \kw{end}}.\]
\end{frame}

\begin{frame}
\centerheadline{Die Zusicherungssprache}
\end{frame}

\begin{frame}
\emph{Zusicherungen} sind Aussagen, die man in einem bestimmten Zustand als
erfüllt sehen will. Zum Beispiel ist nach der Ausführung des Kommandos
$\code{y := x*x}$ die Zusicherung $\code{y}\ge 0$
erfüllt.\pause

\parspace
Zusicherungen sind logische Formeln, die der Sprache der einsortigen Logik
erster Stufe entstammen sollen. Die logische Sprache wird passend zur
Programmiersprache IMP definiert, dergestalt dass das Diskursuniversum
$\Int$ sei und die Variablen aus $\Loc$ in den Termen auftauchen dürfen.
Benötigte Funktionssymbole der Signatur $\Int^n\to\Int$ und
Relationssymbole der Signatur $\Int^n\to\Bool$ zu $n\in\N_{\ge 0}$ kann man
je nach Bedarf in ihrer üblichen Bedeutung hinzufügen; die Operatoren
von IMP sollen dabei aber mindestens verfügbar sein.\pause

\parspace
Wir müssen nun allerdings zwischen Variablen $x,y,z$ und Variablen
$\code{x},\code{y},\code{z}\in\Loc$ unterscheiden. Die kursiven tauchen
als freie und gebundene Variablen in den Formeln auf. Die aufrechten
verhalten sich dagegen wie Konstantensymbole, über sie kann nicht
quantifiziert werden.
\end{frame}

\begin{frame}
Die Notation $I,s\models A$ stehe für die Aussage, dass
die Interpretation $I=(\mathcal M,\beta)$ und der Zustand $s$
die Formel $A$ erfüllen. Hierbei ist $\mathcal M$ eine Struktur, die
die Bedeutung der Funktions- und Relationssymbole festlegt,
und $\beta$ eine Belegung der kursiven Variablen. Der Zustand $s$
belegt die aufrechten Variablen. {\footnotesize Die Definition der
Erfüllung geschieht analog zur gewöhnlichen Logik erster Stufe, weshalb
ich sie hier nicht näher ausführen will.}\pause

\parspace
Wir betrachten nur das Modell $\mathcal M_0$, das die Symbole mit ihrer üblichen
Bedeutung versieht. Daher sei $\mathcal I_0:=\{(\mathcal M,\beta)\mid\mathcal M=\mathcal M_0\}$,
das heißt, $\mathcal I_0$ sei die Menge der Interpretationen, deren Struktur
$\mathcal M_0$ ist.\pause

\parspace
Wir schreiben später $s\models A$ als Abkürzung für $I,s\models A$,
sofern sich dadurch keine Zweideutigkeiten ergeben.
\end{frame}

\begin{frame}
\centerheadline{Der Hoare-Kalkül}
\end{frame}

\begin{frame}[fragile]
An ein Kommando $c$ können wir Zusicherungen machen. Wir notieren
$\{A\}c\{B\}$ für die Aussage, dass die Aussage $B$ unter allen Umständen
nach der Ausführung von $c$ erfüllt ist, sofern die Aussage $A$ zuvor
erfüllt war. Man nennt $A$ diesbezüglich eine \emph{Vorbedingung}
und $B$ eine \emph{Nachbedingung} von $c$.\pause

\parspace
Beispiel für ein allgemeingültiges Tripel:
\begin{lstlisting}[language=IMP, xleftmargin=\mathindent, mathescape]
$\{\kw{true}\}$
y := x*x
$\{\code{y}\ge 0\}$
\end{lstlisting}
\end{frame}

\begin{frame}
Die Allgemeingültigkeit des Tripels $\{A\}c\{B\}$ wird dahingehend
definiert als
\[(\models\{A\}c \{B\}) \;:\Leftrightarrow\; \forall I\in\mathcal I_0\colon\forall s\in S\colon
(I,s\models A)\Rightarrow\forall s'\in S\colon \evC\qb{c}(s)=s'\Rightarrow (I,s'\models B).\]\pause
{\footnotesize In Bezug auf IMP ist der Wert $s'$,
sofern $\evC\qb{c}(s)$ existiert, eindeutig bestimmt.
Denken wir uns nun den ungültigen Zustand $\bot$, der sich ergeben soll,
wenn $c$ eine nicht terminierende Schleife ist, bekommt man eine totale Funktion
$\evC\qb{c}\colon S\cup\{\bot\}\to S\cup\{\bot\}$, wobei $\evC\qb{c}(\bot):=\bot$
gesetzt wird. Diesbezüglich verkürzt sich die Allgemeingültigkeit zu
\[(\models\{A\}c \{B\}) \;\Leftrightarrow\; \forall I\in\mathcal I_0\colon\forall s\in S\colon
(I,s\models A)\Rightarrow (I,\evC\qb{c}(s)\models B).\]
Hierbei verlangt man, dass $I,\bot\models A$ für jede Formel $A$ gilt.}
\end{frame}

\begin{frame}
\strong{Bemerkung.}
Wir können die Zusicherungssprache erweitern um eine modale Operation
$\Box_c B$ je Kommando $c$, üblicherweise $[c]B$ geschrieben.\pause{} Deren Semantik sei
\[(I,s\models [c]B) \;:\Leftrightarrow\; \forall s'\in S\colon R_c(s,s')\Rightarrow (I,s'\models B)\]
mit der Zugänglichkeitsrelation $R_c(s,s')\,:\Leftrightarrow\, \evC\qb{c}(s)=s'$.\pause

\parspace
Diesbezüglich sind $\{A\}c\{B\}$ und $A\Rightarrow [c]B$ semantisch äquivalent.\pause

\parspace
Die so erweiterte Logik nennt man die \emph{dynamische Logik} von IMP.
Die denotationelle Semantik nimmt hierbei die Rolle einer Kripke"=Semantik
ein, wobei die Zustände die Kripke"=Welten sind.
\end{frame}

\begin{frame}
Ein Tripel $\{A\}c\{B\}$ fordert allerdings nicht, dass $c$ terminiert;
lediglich dass $B$ gilt, \emph{falls} $c$ terminiert. Man nennt ein
solches Programm \emph{partiell korrekt} in Bezug auf $A,B$.\pause

\parspace
Terminiert das Programm zusätzlich, nennen wir es \emph{total korrekt}
in Bezug auf $A,B$ und notieren dies $[A]c[B]$.\pause

\parspace
Da die Semantik von IMP eine deterministische ist, ist der Zustand $s'=\evC\qb{c}(s)$,
sofern dieser existiert, eindeutig bestimmt; es gibt also keine weiteren Zustände,
an die wir $B$ fordern müssten. Daher ergibt sich
\[(\models [A]c [B]) \;\Leftrightarrow\; \forall I\in\mathcal I_0\colon\forall s\in S\colon
(I,s\models A)\Rightarrow\exists s'\in S\colon \evC\qb{c}(s)=s'\land (I,s'\models B).\]\pause
In Bezug auf die Modaloperation $\Diamond_c B$ bzw. $\langle c\rangle B$ mit der Semantik
\[(s\models\langle c\rangle B)\;:\Leftrightarrow\; \exists s'\in S\colon R_c(s,s')\land (s'\models B).\]
sind nun $[A]c[B]$ und $A\Rightarrow\langle c\rangle B$ semantisch äquivalent.

\parspace
{\footnotesize Außerdem bestehen die Äquivalenzen
$\langle c\rangle B \Leftrightarrow \lnot [c]\lnot B$
und $[c]B\Leftrightarrow\lnot\langle c\rangle\lnot B$.}
\end{frame}

\begin{frame}
Die Formalismen stehen im engen Bezug zum dijkstraschen wp-Kalkül.
Dijkstra notiert $\mathrm{wp}(c,B)$ für die schwächste Vorbedingung,
engl. \emph{weakest precondition}, unter der das Programm $c$ in einem
Zustand terminiert, in dem $B$ erfüllt ist. Für jede Vorbedingung $A$
gilt daher $A\Rightarrow\mathrm{wp}(c,B)$.\pause

\parspace
Analog steht $\mathrm{wlp}(c,B)$ für die schwächste liberale
Vorbedingung, die lediglich $B$ fordert, falls das Programm terminiert.
Demnach sind $[c]B$ und $\mathrm{wlp}(c,B)$ bedeutungsgleich.
Entsprechend drücken $\{A\}c\{B\}$ und $A\Rightarrow\mathrm{wlp}(c,B)$
dasselbe aus.\pause

\parspace
Aufgrund der deterministischen Natur von IMP sind des Weiteren
$\langle c\rangle B$ und $\mathrm{wp}(c,B)$ bedeutungsgleich.
Entsprechend drücken $[A]c[B]$ und $A\Rightarrow\mathrm{wp(c,B)}$
dasselbe aus.
\end{frame}

\begin{frame}[fragile]
Wir diskutieren nun die \strong{Schlussregeln} des Kalküls.\pause

\parspace
\begin{block}{Regel zur Zuweisung}
\[\dfrac{}{\vdash\{A[X:=a]\}X:=a\{A\}}\]
\end{block}
Eine Schlussregel ohne Prämissen. Mit $A[X:=a]$ ist hierbei die
Formel gemeint, die aus $A$ hervorgeht, indem jedes Vorkommen von
$X$ in $A$ durch den Term $a$ ersetzt wird.\pause

\parspace
Zum Beispiel erkennt man das Tripel
\begin{lstlisting}[language=IMP, xleftmargin=\mathindent, mathescape]
$\{\code{x}\ge 0\}$
x := x + 1
$\{\code{x}\ge 1\}$
\end{lstlisting}
unschwer als allgemeingültig. Dieses fällt unter die Fittiche der
Regel, indem $X:=\code{x}$, $a:=\code{x}+1$ und $A:=(\code{x}\ge 1)$
gesetzt wird. Damit ergibt sich $A[X:=a]$ zu $\code{x} + 1\ge 1$,
was logisch äquivalent zu $\code{x}\ge 0$ ist.
\end{frame}

\begin{frame}
Ob die Rechnung stimmt, prüft man so: Man betrachtet die Nachbedingung,
wendet auf diese die mit der Zuweisung übereinstimmende Substitution an,
und dies muss dann in der Vorbedingung resultieren.\pause

\parspace
\strong{Beweis der Gültigkeit der Regel.} Wir wollen
\[\models\{A[X:=a]\}X:=a\{A\}\]
zeigen.\pause{} Sei dazu $s$ fest, aber beliebig. Es gelte $s\models A[X:=a]$.
Des Weiteren gelte $\evC\qb{X:=a}(s)=s'$. Zu zeigen ist $s'\models A$.\pause

\parspace
Gemäß der denotationellen Semantik gilt $\evC\qb{X:=a}(s)=s[X:=a]$, womit
wir $s'=s[X:=a]$ haben. Schließlich folgt die Behauptung vermittels
der Äquivalenz
\[(s\models A[X:=a])\;\Leftrightarrow\; (s[X:=a]\models A),\]
die man unschwer als richtig erkennt oder pedantisch per struktureller
Induktion über den Aufbau von $A$ beweisen kann.\,\qedsymbol
\end{frame}

\begin{frame}
\begin{block}{Regel zum leeren Kommando}
\[\dfrac{}{\vdash\{A\}\kw{skip}\{A\}}\]
\end{block}\pause
\strong{Beweis ihrer Gültigkeit.}
Wir wollen $\models\{A\}\kw{skip}\{B\}$ zeigen.
Dazu sei $s$ fest, aber beliebig. Es gelte $s\models A$. Des Weiteren
gelte $\evC\qb{\kw{skip}}(s)=s'$. Zu zeigen ist $s'\models A$.\pause

\parspace
Gemäß der denotationellen Semantik gilt $\evC\qb{\kw{skip}}(s)=s$, womit
wir $s'=s$ und somit bereits die Behauptung haben.\,\qedsymbol
\end{frame}

\begin{frame}
\begin{block}{Regel zur Sequenz von Kommandos}
\[\dfrac{\vdash\{A\}c\{B\}\qquad \vdash\{B\}c'\{C\}}{\vdash\{A\}c; c'\{C\}}\]
\end{block}\pause
\strong{Beweis ihrer Gültigkeit.} Wir wollen $\models\{A\} c; c'\{C\}$ zeigen.
Dazu sei $s$ fest, aber beliebig. Es gelte $s\models A$. Des Weiteren
gelte $\evC\qb{c;c'}(s)=s''$. Zu zeigen ist $s''\models C$.\\
{\footnotesize Die letzten beiden C's stehen für Unterschiedliches; aber eine
Verwechslung ist ausgeschlossen, denke ich.}\pause

\parspace
Gemäß der denotationellen Semantik existiert $s'=\evC\qb{c}(s)$ mit
\[\evC\qb{c;c'}(s) = \evC\qb{c'}(s')=s''.\]
Aus $s\models A$ und $\models\{A\}c\{B\}$ folgt nun zunächst $s'\models B$.
Mit $\models\{B\}c'\{C\}$ folgt daraufhin $s''\models C$ aus
$s'\models B$.\,\qedsymbol
\end{frame}

\begin{frame}
\begin{block}{Regel zur Verzweigung}
\[\dfrac{\vdash\{A\land b\}c\{C\}\qquad\vdash\{A\land\lnot b\}c'\{C\}}{
\vdash\{A\}\text{\kw{if} $b$ \kw{then} $c$ \kw{else} $c'$ \kw{end}}\{C\}}\]
\end{block}\pause
\strong{Beweis ihrer Gültigkeit.} Sei $s$ fest, aber beliebig.
Es gelte $s\models A$. Des Weiteren gelte
\[\evC\qb{\text{\kw{if} $b$ \kw{then} $c$ \kw{else} $c'$ \kw{end}}}(s) = s'.\]
Zu zeigen ist $s'\models C$.\pause{} Fallunterscheidung. Im Fall
$\evB\qb{b}(s)=\kw{true}$ gilt $s\models b$, also $s\models A\land b$.
Außerdem ergibt sich in diesem Fall die Vereinfachung
\[\evC\qb{\text{\kw{if} $b$ \kw{then} $c$ \kw{else} $c'$ \kw{end}}}(s) = \evC\qb{c}(s).\]
Vermittels $\models\{A\land b\}c\{C\}$ erhalten wir somit $s'\models C$.
Die Argumentation im Fall $\evB\qb{b}(s)=\kw{false}$ verläuft analog.\,\qedsymbol
\end{frame}

\begin{frame}
\begin{block}{Regel zur Schleife}
\[\dfrac{\vdash\{A\land b\}c\{A\}}{
\vdash\{A\}\text{\kw{while} $b$ \kw{do} $c$ \kw{end}}\{A\land\lnot b\}}\]
\end{block}\pause
Sofern $A$ also sowohl vor der Schleife gilt, als auch vor und hinter
dem Schleifenrumpf, muss $A$ auch hinter der Schleife gelten. Man nennt
$A$ hierbei eine \emph{Schleifeninvariante}. Die Auffindung einer
zielführenden Invariante stellt eine wesentliche Schwierigkeit bei der
Programmverifikation dar.
\end{frame}

\begin{frame}
\strong{Beweis ihrer Gültigkeit.} Sei $s$ fest, aber beliebig.
Es gelte $s\models A$. Des Weiteren gelte
\[\evC\qb{\text{\kw{while} $b$ \kw{do} $c$ \kw{end}}}(s) = s''.\]
Zu zeigen ist $s''\models A\land\lnot b$.\pause{} Induktion über die
Anzahl der Durchläufe. Im Anfang findet kein Durchlauf statt. Das
geht aber nur, wenn $\evB\qb{b}(s)=\kw{false}$, also $s\models\lnot b$.
Da dem Zustand $s$ nach der Semantik der Schleife in diesem Fall keine Änderung
widerfährt, gilt $s=s''$. Demnach gilt $s''\models A$ und $s''\models\lnot b$,
also $s''\models A\land\lnot b$.\pause

\parspace
Zum Induktionsschritt. Da die Schleife durchlaufen wird, existiert
$s'$ mit $\evC\qb{c}(s)=s'$. Die restlichen null der mehr Schleifendurchläufe
führen dann zum Zustand $s''$, das heißt, $\varphi_{b,c}(s')=s''$.
Die Induktionsvoraussetzung ist, dass $s''\models A\land\lnot b$
aus $s'\models A$ folgt. Es verbleibt also $s'\models A$ zu zeigen.
Weil der erste Durchlauf stattfindet, muss des Weiteren $\evB\qb{b}(s)=\kw{true}$,
also $s\models b$ gelten, womit wir $s\models A\land b$ haben.
Vermittels $\models\{A\land b\}c\{A\}$, was ja Kraft der Prämisse der Regel
zur Verfügung steht, erhält man schließlich $s'\models A$.\,\qedsymbol
\end{frame}

\begin{frame}
\begin{block}{Regel zur Verstärkung der Vorbedingung}
\[\dfrac{\vdash A'\Rightarrow A\qquad\vdash\{A\}c\{B\}}{
\vdash\{A'\}c\{B\}}\]
\end{block}\pause
\strong{Beweis ihrer Gültigkeit.} Es gelte $s\models A'$.
Des Weiteren gelte $\evC\qb{c}(s)=s'$. Zu zeigen ist $s'\models B$.\pause{}
Mit $\models A'\Rightarrow A$ erhält man zunächst $s\models A$.
Mit $\models\{A\}c\{B\}$ daraufhin $s'\models B$.\,\qedsymbol
\end{frame}

\begin{frame}
\begin{block}{Regel zur Abschwächung der Nachbedingung}
\[\dfrac{\vdash\{A\}c\{B\}\qquad\vdash B\Rightarrow B'}{
\vdash\{A\}c\{B'\}}\]
\end{block}\pause
\strong{Beweis ihrer Gültigkeit.} Es gelte $s\models A$.
Des Weiteren gelte $\evC\qb{c}(s)=s'$. Zu zeigen ist $s'\models B'$.\pause{}
Mit $\models\{A\}c\{B\}$ erhält man zunächst $s'\models B$. Mit
$\models B\Rightarrow B'$ erhält man daraufhin $s'\models B'$.\,\qedsymbol.
\end{frame}

\begin{frame}
\begin{block}{Ersetzungsregeln}
\[\dfrac{\vdash A\Leftrightarrow A'\qquad\vdash\{A\}c\{B\}}{
  \vdash\{A'\}c\{B\}},\qquad
\dfrac{\vdash B\Leftrightarrow B'\qquad\vdash\{A\}c\{B\}}{
  \vdash\{A\}c\{B'\}}\]
\end{block}\pause
\strong{Beweis.} Folgt unmittelbar aus der Regel zur Verstärkung der
Vorbedingung bzw. Abschwächung der Nachbedingung.\,\qedsymbol
\end{frame}

\begin{frame}
\centerheadline{Verifikation des ersten Programms}
\end{frame}

\begin{frame}[fragile]
Gesucht ist ein Beweis des Tripels:
\begin{lstlisting}[language=IMP, xleftmargin=\mathindent, mathescape]
$\{\mathtt n\ge 0\}$
y := 1; k := n;
while not k = 0 do
   y := y*x;
   k := k - 1
end
$\{\mathtt{y = x^n}\}$
\end{lstlisting}
\end{frame}

\begin{frame}[fragile]
Vermittels der Regel zur Zuweisung -- zuzüglich Reflexivität der
Gleichheit und der Ersetzungsregel -- ergibt sich erst einmal die Ableitung:
\[\begin{prooftree}
  \infer0{\vdash \mathtt n\ge 0\,\Leftrightarrow\,\mathtt n\ge 0\land 1 = 1}
  \infer0{\vdash\{\mathtt n\ge 0\land 1 = 1\}\;\code{y := 1}\;\{\mathtt n\ge 0\land\mathtt y = 1\}}
\infer2{\vdash\{\mathtt n\ge 0\}\;\code{y := 1}\;\{\mathtt n\ge 0\land\mathtt y = 1\}}
\end{prooftree}\]\pause
Wir fügen bereits ermittelte Zusicherungen in den Quelltext ein:
\begin{lstlisting}[language=IMP, xleftmargin=\mathindent, mathescape]
$\{\mathtt n\ge 0\}$
y := 1; k := n;
$\{\mathtt k\ge 0\land\mathtt y = 1\}$
while not k = 0 do
   y := y*x;
   k := k - 1
end
$\color{gray}\{\mathtt{y = x^n}\}?$
\end{lstlisting}
\end{frame}

\begin{frame}[fragile]
Die Zielführende Schleifeninvariante ist hier $\mathtt{y =
x^{n-k}}$. Vor der Schleife gilt ja $\mathtt{n=k}$.
Mit $\mathtt{x}^0 = 1$ kommt man somit auf $\mathtt y = 1$.
{\footnotesize Bezüglich $0^0:=1$ gilt dies auch im Fall $\code{x}=0$.}\pause{}

\parspace
Wir durchziehen das Programm nun gemäß den Regeln sukzessive mit
Zusicherungen und gelangen daraufhin zum Abschluss der Verifikation:
\begin{lstlisting}[language=IMP, xleftmargin=\mathindent, mathescape]
$\{\mathtt n\ge 0\}$
y := 1; k := n;
$\{\mathtt{k\ge \mathrm 0\land y = \mathrm 1\land y = x^{n-k}}\}$
while not k = 0 do
   $\{\mathtt{y = x^{n-k}}\}$
   $\{\mathtt{y\cdot x = x^{n-k}\cdot x}\}$
   y := y*x;
   $\{\mathtt{y = x^{n-k}\cdot x = x^{n-(k - 1)}}\}$
   k := k - 1
   $\{\mathtt{y = x^{n-k}}\}$
end
$\{\mathtt{y = x^{n-k}\land k=\mathrm 0}\}$
$\{\mathtt{y = x^n}\}$
\end{lstlisting}
\end{frame}

\begin{frame}
\centerheadline{Verifikation des zweiten Programms}
\end{frame}

\begin{frame}[fragile]
Gesucht ist der Beweis von:
\begin{lstlisting}[language=IMP, xleftmargin=\mathindent, mathescape]
$\{\mathtt n\ge 0\}$
y := 1; a := x; k := n;
while 2 <= k do
   q := k/2; r := k - 2*q;
   if r = 1 then y := y*a else skip end;
   k := q; a := a*a
end;
if k = 1 then y := y*a else skip end
$\{\mathtt{y = x^n}\}$
\end{lstlisting}\pause

Wir betrachten das Programm zunächst vom Ende aus. Dort findet sich das Tripel
\begin{lstlisting}[language=IMP, xleftmargin=\mathindent, mathescape]
$\{\mathtt{(y = x^n\land k = \mathrm 0)\lor (y\cdot a = x^n\land k = \mathrm 1)}\}$
if k = 1 then y := y*a else skip end
$\{\mathtt{y = x^n}\}.$
\end{lstlisting}
Nimmt man bei der Vorbedingung nämlich die Fallunterscheidung in die
beiden Seiten der Disjunktion vor, gelangt man in beiden Fällen
zur Nachbedingung.
\end{frame}

\begin{frame}[fragile]
Am Anfang der Schleife gilt $\mathtt{a^k = x^n}$, was etwas mit der
Schleifeninvariante zu tun haben könnte. Ist diese Gleichung so
modifizierbar, dass unter ihr die Zusicherung
\[\mathtt{(y = x^n\land k = \mathrm 0)\lor (y\cdot a = x^n\land k = \mathrm 1)}\]
gilt?\pause{} Ja, nämlich zu $\mathtt{y\cdot a^k = x^n}$, denn
\begin{gather*}
\text{im Fall $\mathtt k=0$ gilt $\mathtt{y\cdot a^k = y}$},\\
\text{im Fall $\mathtt k=1$ gilt $\mathtt{y\cdot a^k = y\cdot a}$}.
\end{gather*}\pause
Die Gleichung $\mathtt{y\cdot a^k = x^n}$ gilt ebenfalls vor der
Schleife, da dort $\mathtt y = 1$ ist. Tatsächlich stellt sie eine
zielführende Schleifeninvariante dar. Es verbleibt also zu bestätigen,
dass sie auch am Ende des Schleifenrumpfs gilt.\pause

\parspace
Zunächst ergibt sich
\begin{lstlisting}[language=IMP, xleftmargin=\mathindent, mathescape]
$\{\mathtt{y\cdot a^k = x^n}\}$
q := k/2; r := k - 2*q
$\{\mathtt{y\cdot a^k = x^n\land q = \big\lfloor\frac{k}{\mathrm 2}\big\rfloor
\land r = k - \mathrm{2}q}\}.$
\end{lstlisting}
\end{frame}

\begin{frame}[fragile]
Wegen $r\in\{0,1\}$ findet sich bei der Verzweigung nun
\begin{lstlisting}[language=IMP, xleftmargin=\mathindent, mathescape]
$\{\mathtt{y\cdot a^k = x^n\land q = \big\lfloor\frac{k}{2}\big\rfloor
\land r = k - \mathrm 2q}\}$
if r = 1 then y := y*a else skip end
$\{\mathtt{((y\cdot a^{k-\mathrm 1} = x^n\land r = \mathrm 1)\lor
(y\cdot a^k = x^n\land r = \mathrm 0))
\land r = k - \mathrm 2q}\}.$
\end{lstlisting}\pause
Die lange Nachbedingung kann man kompakter fassen als
\[\mathtt{y\cdot a^{k - r} = x^n\land r = k - \mathrm 2q},\]
wobei sich des Weiteren $\mathtt{k - r = \mathrm 2q}$ ergibt.\pause

\parspace
Zum Rest des Schleifenrumpfs findet sich somit schließlich
\begin{lstlisting}[language=IMP, xleftmargin=\mathindent, mathescape]
$\{\mathtt{y\cdot a^{\mathrm 2q} = x^n}\}$
k := q
$\{\mathtt{y\cdot a^{\mathrm 2k} = x^n}\}$
a := a*a
$\{\mathtt{y\cdot a^k = x^n}\}.$
\end{lstlisting}
Quod erat demonstrandum.
\end{frame}

\begin{frame}
\centerheadline{Maschinengestützte Formalisierung des Kalküls}
\end{frame}

\begin{frame}
Im Anschluss will ich noch eine Formalisierung der Überlegungen in Coq
ausführen. Dies setzt allerdings Grundwissen über die Sprache Gallina
voraus, und wie logisches Schließen vermittels Taktiken abläuft.

\parspace
Der vollständige Quelltext findet sich im Anhang. Um den Ablauf der
Beweise zu betrachten, muss CoqIDE genutzt werden.
\end{frame}

\begin{frame}[fragile]
Wir indizieren die Programmvariablen durch die natürlichen Zahlen.\\
Hierbei steht \code{loc 0}, \code{loc 1}, \code{loc 2} usw. für
$\mathtt{x}_0,\mathtt{x}_1,\mathtt{x}_2$ usw.

\parspace
Ein Zustand stellt sich daher schlicht als Funktion $s\colon\N\to\mathrm{Int}$
dar. Für die ganzen Zahlen $\mathrm{Int}$ verwenden wir hierbei den Typ
\code{Z} aus \code{ZArith}.

\parspace
\begin{lstlisting}[language=Coq, xleftmargin=\mathindent]
Require Import ZArith.ZArith.
Require Import Bool.Bool.

Definition Loc := nat.
Definition State := Loc -> Z.
\end{lstlisting}
\end{frame}

\begin{frame}[fragile]
Als nächstes erfolgt die Erklärung der Sprache IMP vermittels
induktiver Tyen.\\
{\footnotesize Die Division will ich der Prägnanz halber erst einmal
entfallen lassen. Wir können sie später ggf. nachtragen.}

\begin{tabular}{@{\hspace{\mathindent}}l@{\qquad\qquad}l@{}}
\vtop{\null\hbox{
\begin{lstlisting}[language=Coq]
Inductive Aexpr :=
| int (n: Z)
| loc (X: Loc)
| add (a1 a2: Aexpr)
| sub (a1 a2: Aexpr)
| mul (a1 a2: Aexpr).

Inductive Bexpr :=
| btrue
| bfalse
| beq (a1 a2: Aexpr)
| ble (a1 a2: Aexpr)
| bnot (b: Bexpr)
| band (b1 b2: Bexpr)
| bor (b1 b2: Bexpr).
\end{lstlisting}
}}
&
\vtop{\null\hbox{
\begin{lstlisting}[language=Coq]
Inductive Com :=
| Skip
| Assign (X: Loc) (a: Aexpr)
| Seq (c1 c2: Com)
| If (b: Bexpr) (c1 c2: Com)
| While (b: Bexpr) (c: Com).
\end{lstlisting}
}}
\end{tabular}
\end{frame}

\begin{frame}[fragile]
Nun die Auswertungsfunktionen $A$, $B$. Für die Werte von $B$ nehmen wir
\code{bool} statt \code{Bexpr}, weil dieses bereits Funktionalität zum
Argumentieren mit sich bringt.

\begin{small}
\begin{lstlisting}[language=Coq, xleftmargin=\mathindent]
Fixpoint evA (a: Aexpr) (s: State): Z :=
  match a with
  | int n => n
  | loc X => s X
  | add a1 a2 => evA a1 s + evA a2 s
  | sub a1 a2 => evA a1 s - evA a2 s
  | mul a1 a2 => evA a1 s * evA a2 s
  end.

Fixpoint evB (b: Bexpr) (s: State): bool :=
  match b with
  | btrue => true
  | bfalse => false
  | beq a1 a2 => Z.eqb (evA a1 s) (evA a2 s)
  | ble a1 a2 => Z.leb (evA a1 s) (evA a2 s)
  | bnot b => negb (evB b s)
  | band b1 b2 => andb (evB b1 s) (evB b2 s)
  | bor b1 b2 => orb (evB b1 s) (evB b2 s)
  end.
\end{lstlisting}
\end{small}
\end{frame}

\begin{frame}[fragile]
Es wäre nun denkbar, die Terme und Formeln der Zusicherungssprache
ebenfalls vermittels induktiver Typen zu erklären. Dies führt allerdings
zu einem Rattenschwanz von Umständlichkeiten, da es zu einer
sogenannten \emph{tiefen Einbettung} des logischen Systems führt.
Wir müssten, wie für jede neue logische Sprache, einen Kalkül des natürlichen
Schließens formulieren und streng genommen auch noch dessen Korrektheit
beweisen.\pause

\parspace
Wir unternehmen stattdessen den Ansatz, die Zusicherungen als
Prädikate darzustellen, die Zustände als Argument bekommen. Eine
Zusicherung $A$ sei also eine Funktion
\[A\colon\texttt{State}\to\texttt{Prop}.\]\pause
Die Variante des Zustandes $s$ wird auf die übliche Weise erklärt.
Im Folgenden steht \code{variant s X n} also für $s[X:=n]$ bzw. $s[n/X]$.

\parspace
\begin{lstlisting}[language=Coq, xleftmargin=\mathindent]
Definition variant (s: State) (X: Loc) (n: Z): State :=
  fun Y => if Nat.eqb X Y then n else s Y.
\end{lstlisting}
\end{frame}

\begin{frame}
Die Übergangsfunktion $\evC$ ist wie gesagt eine partielle. Als Folge
dessen tut sich das Problem auf, dass ihre Implementierung unmöglich
wird, da in der Typentheorie nur solche Rekursionen fassbar
sind, deren Terminierung geklärt ist.

\parspace
Wir lösen dieses Problem in zwei Schritten. Erstens wird die partielle
Funktion $\evC$ als Relation formuliert. Statt $\evC\qb{c}(s)=s'$ notieren
wir also $\evC\qb{c}(s)(s')$ bzw. \code{evC c s s\ttq}. Zweitens beschränken
wir die rekursive Auswertung der Schleife über eine rekursive
Konstruktion \code{iter} auf eine maximale Tiefe $N$. Da nicht klar
ist, wie groß $N$ bei
\[\evC\qb{\kw{while}\; b\; \kw{do}\; c\; \kw{end}}(s)(s')\]
mindestens sein muss, wird lediglich die Existenz von $N$ gefordert.
\end{frame}

\begin{frame}[fragile]
\begin{lstlisting}[language=Coq, xleftmargin=\mathindent]
Fixpoint iter (ev: State -> State -> Prop)
(N: nat) (b: Bexpr) (s s1: State): Prop :=
  match N with
  | 0 => False
  | S N => if evB b s then
      exists s0, ev s s0 /\ iter ev N b s0 s1
    else s1 = s
  end.

Fixpoint evC (c: Com) (s s1: State): Prop :=
  match c with
  | Skip => s1 = s
  | Assign X a => s1 = variant s X (evA a s)
  | Seq c1 c2 => exists s0, evC c1 s s0 /\ evC c2 s0 s1
  | If b c1 c2 => if evB b s then evC c1 s s1 else evC c2 s s1
  | While b c => exists N, iter (evC c) N b s s1
  end.
\end{lstlisting}
\end{frame}

\begin{frame}[fragile]
Nun kommen die Erklärungen zur Gültigkeit von Zusicherungen.

\begin{lstlisting}[language=Coq, xleftmargin=\mathindent]
Definition Assertion := State -> Prop.
Definition sat (s: State) (A: Assertion): Prop := A s.

Definition valid (A: Assertion) (c: Com) (B: Assertion) :=
  forall s, sat s A -> forall s1, evC c s s1 -> sat s1 B.

Definition subst (A: Assertion) (X: Loc) (a: Aexpr): Assertion :=
  fun s => A (variant s X (evA a s)).
\end{lstlisting}

Diese Festlegungen sind folgendermaßen zu verstehen:
\begin{itemize}
\item Es steht \code{sat s A} für die Erfüllung $s\models A$.
\item Es steht \code{valid A c B} für die Gültigkeit des Tripels $\{A\}c\{B\}$.
\item Es steht \code{subst A X a} für die Substitution $A[X:=a]$, die direkt
als $s\mapsto A(s[X:=\evA\qb{a}(s)])$ definiert wird, insofern die Formel
von $A$ verborgen bleibt und das Substitutionslemma somit entfallen muss.
{\footnotesize Die beiden $A$ sind von unterschiedlicher Bedeutung --
das zweite steht für \code{evA}, siehe oben.}
\end{itemize}
\end{frame}

\begin{frame}[fragile]
Die Formalisierung des Beweises der Gültigkeit der jeweiligen Schlussregel
kann nun unternommen werden.

\parspace
Zum leeren Kommando findet sich:

\begin{small}
\begin{lstlisting}[language=Coq, xleftmargin=\mathindent]
Theorem skip_intro_is_valid A:
  valid A Skip A.
Proof.
  unfold valid. intros s hs. intros s1 hs1.
  simpl evC in hs1. rewrite hs1. exact hs.
Qed.
\end{lstlisting}
\end{small}

Zur Zuweisung findet sich:

\begin{small}
\begin{lstlisting}[language=Coq, xleftmargin=\mathindent]
Theorem assign_intro_is_valid A X a:
  valid (subst A X a) (Assign X a) A.
Proof.
  unfold valid. intros s hs. intros s1 hs1.
  simpl evC in hs1.
  unfold sat in hs. unfold subst in hs.
  rewrite hs1. exact hs.
Qed.
\end{lstlisting}
\end{small}
\end{frame}

\begin{frame}[fragile]
Zur Sequenz von Kommandos findet sich:

\begin{small}
\begin{lstlisting}[language=Coq, xleftmargin=\mathindent]
Theorem seq_intro_is_valid A B C c1 c2:
  valid A c1 B -> valid B c2 C -> valid A (Seq c1 c2) C.
Proof.
  intros h1 h2. unfold valid. intros s hs. intros s2 hs2.
  simpl evC in hs2.
  destruct hs2 as (s1, (h11, h12)).
  unfold valid in h1. unfold valid in h2.
  assert (h1 := h1 s hs s1 h11).
  exact (h2 s1 h1 s2 h12).
Qed.
\end{lstlisting}
\end{small}
\end{frame}

\begin{frame}[fragile]
Zur Verzweigung findet sich:

\begin{small}
\begin{lstlisting}[language=Coq, xleftmargin=\mathindent]
Definition Conj (A: Assertion) (b: Bexpr) :=
  fun s => sat s A /\ evB b s = true.

Theorem if_intro_is_valid A b C c1 c2:
  valid (Conj A b) c1 C -> valid (Conj A (bnot b)) c2 C ->
  valid A (If b c1 c2) C.
Proof.
  intros h1 h2. unfold valid. intros s hs s1 hs1.
  simpl evC in hs1.
  destruct (evB b s) eqn:heq.
  * unfold valid in h1. apply (h1 s).
    - unfold sat. unfold Conj. rewrite heq.
      exact (conj hs (eq_refl true)).
    - exact hs1.
  * unfold valid in h2. apply (h2 s).
    - unfold sat. unfold Conj.
      simpl evB. rewrite heq. simpl.
      exact (conj hs (eq_refl true)).
    - exact hs1.
Qed.
\end{lstlisting}
\end{small}
\end{frame}

\begin{frame}[fragile]
Zur Schleife findet sich:

\begin{small}
\begin{lstlisting}[language=Coq, xleftmargin=\mathindent]
Theorem while_intro_is_valid A b c:
  valid (Conj A b) c A -> valid A (While b c) (Conj A (bnot b)).
Proof.
  intro h. unfold valid. intros s hs. intros s2 hs2.
  simpl evC in hs2. destruct hs2 as (N, hiter).
  revert s s2 hs hiter.
  induction N as [| N ih].
  * intros s s2 hs hiter. simpl iter in hiter.
    exfalso. exact hiter.
  * intros s s2 hs hiter. simpl iter in hiter.
    destruct (evB b s) eqn:heq.
    - destruct hiter as (s1, (h11, h12)).
      apply (ih s1 s2). clear ih.
      -- unfold valid in h. apply (h s). clear h.
         --- unfold sat. unfold Conj. exact (conj hs heq).
         --- exact h11.
      -- exact h12.
    - rewrite hiter. unfold sat. unfold Conj. split.
      -- exact hs.
      -- simpl evB. rewrite heq. simpl. reflexivity.
Qed.
\end{lstlisting}
\end{small}
\end{frame}

\begin{frame}[fragile]
Zur Verstärkung der Vorbedingung findet sich:

\begin{small}
\begin{lstlisting}[language=Coq, xleftmargin=\mathindent]
Theorem strengthen_precondition_is_valid {A1 A2 B c}:
  (forall s, A2 s -> A1 s) -> valid A1 c B -> valid A2 c B.
Proof.
  intros h1 h2. unfold valid. intros s hs s1 hs1.
  unfold sat in hs. apply (h1 s) in hs. clear h1.
  unfold valid in h2. apply (h2 s).
  * unfold sat. exact hs.
  * exact hs1.
Qed.
\end{lstlisting}
\end{small}

Zur Abschwächung der Nachbedingung findet sich:

\begin{small}
\begin{lstlisting}[language=Coq, xleftmargin=\mathindent]
Theorem weaken_postcondition_is_valid {A B1 B2 c}:
  (forall s, B1 s -> B2 s) -> valid A c B1 -> valid A c B2.
Proof.
  intros h1 h2. unfold valid. intros s hs s1 hs1.
  unfold sat. apply (h1 s1). clear h1.
  fold (sat s1 B1). unfold valid in h2. apply (h2 s).
  * exact hs.
  * exact hs1.
Qed.
\end{lstlisting}
\end{small}
\end{frame}

\begin{frame}[fragile]
\strong{Kurzes Anwendungsbeispiel}

\parspace
Beweis der Gültigkeit des Tripels
$\{\mathtt x = 0\}\; \code{x := x + 1}\; \{\mathtt x = 1\}$.\pause

\begin{small}
\begin{lstlisting}[language=Coq, xleftmargin=\mathindent]
Module Example1.
  Open Scope Z_scope.
  Definition x: nat := 0.

  Goal valid
    (fun s => s x = 0)
    (Assign x (add (loc x) (int 1)))
    (fun s => s x = 1).
  Proof.
    assert (h1 := assign_intro_is_valid
      (fun s => s x = 1) x (add (loc x) (int 1))).
    unfold subst in h1. unfold variant in h1. simpl in h1.
    assert (h2: forall s, s x = 0 -> s x + 1 = 1). {
      intros s heq. rewrite <- (Z.add_cancel_r _ _ 1) in heq.
      simpl in heq. exact heq.
    }
    exact (strengthen_precondition_is_valid h2 h1).
  Qed.
End Example1.
\end{lstlisting}
\end{small}
\end{frame}

\begin{frame}
\strong{Literatur}
\begin{itemize}
\item Glynn Winskel:
  \emph{The Formal Semantics of Programming Languages: An Introduction}.
  The MIT Press, 1993.
\item Krzysztof R. Apt, Ernst-Rüdiger Olderog:
  \emph{Fifty years of Hoare’s logic}.\\
  In: \emph{Formal Aspects of Computing}.
  Band 31, Nr. 6, 2019, S. 751--807.\\
  \href{https://doi.org/10.1007/s00165-019-00501-3}{doi:10.1007/s00165-019-00501-3}.
\item Edsger W. Dijkstra:
  \emph{A Discipline of Programming}. Prentice Hall, 1976.
\item Benjamin C. Pierce u.\,a.:
  \href{https://softwarefoundations.cis.upenn.edu/}{\emph{Software Foundations}}.
\item Nicolas Troquard, Philippe Balbiani:
  \href{https://plato.stanford.edu/entries/logic-dynamic/}{\emph{Propositional Dynamic Logic}}.
  In: \emph{The Stanford Encyclopedia of Philosophy}.
\item David Harel, Dexter Kozen, Jerzy Tiuryn:
  \emph{Dynamic Logic}. The MIT Press, 2000.
\end{itemize}
\end{frame}

\begin{frame}
\strong{Anlagen}

\begin{itemize}
\item\href{https://github.com/JohnBSmith/misc/tree/master/cs/imp}{%
IMP-Interpreter} -- Führt ein IMP-Programm gemäß der denotationellen
Semantik aus. In Python verfasst, in unter 300 Zeilen Quelltext.

\item\href{https://github.com/JohnBSmith/misc/tree/master/cs/imp/proofs}{%
Quelltext der formalisierten Beweise}.
\end{itemize}
\end{frame}

\begin{frame}
Ende.
\vfill\hfill\modest{Januar 2025}\\
\hfill\modest{Creative Commons CC0 1.0}
\end{frame}

\end{document}


\chapter{Anhang}
\section{Griechisches Alphabet}

\begin{tabular}{l|l}
\begin{tabular}[t]{lll}
$\mathrm A$ & $\alpha$   & Alpha\\
$\mathrm B$ & $\beta$    & Beta\\
$\Gamma$    & $\gamma$   & Gamma\\
$\Delta$    & $\delta$   & Delta\\
\noalign{\vspace{1em}}
$\mathrm E$ & $\varepsilon$ & Epsilon\\
$\mathrm Z$ & $\zeta$    & Zeta\\
$\mathrm H$ & $\eta$     & Eta\\
$\Theta$    & $\theta$   & Theta\\
\noalign{\vspace{1em}}
$\mathrm I$ & $\iota$    & Jota\\
$\mathrm K$ & $\kappa$   & Kappa\\
$\Lambda$   & $\lambda$  & Lambda\\
$\mathrm M$ & $\mu$      & My
\end{tabular}
&
\begin{tabular}[t]{lll}
$\mathrm N$ & $\nu$      & Ny\\
$\Xi$       & $\xi$      & Xi\\
$\mathrm O$ & $o$        & Omikron\\
$\Pi$       & $\pi$      & Pi\\
\noalign{\vspace{1em}}
$\mathrm R$ & $\varrho$  & Rho\\
$\Sigma$    & $\sigma$   & Sigma\\
$\mathrm T$ & $\tau$     & Tau\\
$\mathrm Y$ & $\upsilon$ & Ypsilon\\
\noalign{\vspace{1em}}
$\Phi$      & $\varphi$  & Phi\\
$\mathrm X$ & $\chi$     & Chi\\
$\Psi$      & $\psi$     & Psi\\
$\Omega$    & $\omega$   & Omega 
\end{tabular}
\end{tabular}

\section{Frakturbuchstaben}
\begin{tabular}{l|l}
\begin{tabular}[t]{l@{\hskip 2pt}ll@{\hskip 2pt}l}
A & a & $\mathfrak A$ & $\mathfrak a$\\
B & b & $\mathfrak B$ & $\mathfrak b$ \\
C & c & $\mathfrak C$ & $\mathfrak c$\\
D & d & $\mathfrak D$ & $\mathfrak d$\\
\noalign{\vspace{1em}}
E & e & $\mathfrak E$ & $\mathfrak e$\\
F & f & $\mathfrak F$ & $\mathfrak f$\\
G & g & $\mathfrak G$ & $\mathfrak g$\\
H & h & $\mathfrak H$ & $\mathfrak h$\\
\noalign{\vspace{1em}}
I & i & $\mathfrak I$ & $\mathfrak i$\\
J & j & $\mathfrak J$ & $\mathfrak j$\\
K & k & $\mathfrak K$ & $\mathfrak k$\\
L & l & $\mathfrak L$ & $\mathfrak l$\\
\noalign{\vspace{1em}}
M & m & $\mathfrak M$ & $\mathfrak m$\\
N & n & $\mathfrak N$ & $\mathfrak n$
\end{tabular}
&
\begin{tabular}[t]{l@{\hskip 2pt}ll@{\hskip 2pt}l}
O & o & $\mathfrak O$ & $\mathfrak o$\\
P & p & $\mathfrak P$ & $\mathfrak p$\\
Q & q & $\mathfrak Q$ & $\mathfrak q$\\
R & r & $\mathfrak R$ & $\mathfrak r$\\
\noalign{\vspace{1em}}
S & s & $\mathfrak S$ & $\mathfrak s$\\
T & t & $\mathfrak T$ & $\mathfrak t$\\
U & u & $\mathfrak U$ & $\mathfrak u$\\
V & v & $\mathfrak V$ & $\mathfrak v$\\
\noalign{\vspace{1em}}
W & w & $\mathfrak W$ & $\mathfrak w$\\
X & x & $\mathfrak X$ & $\mathfrak x$\\
Y & y & $\mathfrak Y$ & $\mathfrak y$\\
Z & z & $\mathfrak Z$ & $\mathfrak z$
\end{tabular}
\end{tabular}

\section{Mathematische Konstanten}
\begin{enumerate}
\item Kreiszahl\\
$\pi = 3{,}14159\;26535\;89793\;23846\;26433\;83279\ldots$

\item Eulersche Zahl\\
$\ee = 2{,}71828\;18284\;59045\;23536\;02874\;71352\ldots$

\item Euler-Mascheroni-Konstante\\
$\gamma = 0{,}57721\;56649\;01532\;86060\;65120\;90082\ldots$

\item Goldener Schnitt, $(1+\sqrt{5})/2$\\
$\varphi = 1{,}61803\;39887\;49894\;84820\;45868\;34365\ldots$

\item 1. Feigenbaum-Konstante\\
$\delta = 4{,}66920\;16091\;02990\;67185\;32038\;20466\ldots$

\item 2. Feigenbaum-Konstante\\
$\alpha = 2{,}50290\;78750\;95892\;82228\;39028\;73218\ldots$
\end{enumerate}

\newpage
\section{Physikalische Konstanten}
\begin{enumerate}
\item Lichtgeschwindigkeit im Vakuum\\
$c=299\;792\;458\;\unit{m/s}$

\item Elektrische Feldkonstante\\
$\varepsilon_0 = 8{,}854\;187\;817\;620\;39\times 10^{-12}\:\unit{F/m}$

\item Magnetische Feldkonstante\\
$\mu_0 = 4\pi\times 10^{-7}\:\unit{H/m}$

\item Elementarladung\\
$e = 1{,}602\;176\;6208\;(98)\times 10^{-19}\,\unit{C}$

\item Gravitationskonstante\\
$G = 6{,}674\;08\;(31)\times 10^{-11}\,\unit{m^3/(kg\,s^2)}$

\item Avogadro-Konstante\\
$N_A = 6{,}022\;140\;857\;(74)\times 10^{23}/\unit{mol}$

\item Boltzmann-Konstante\\
$k_B = 1{,}380\;648\;52\;(79)\times 10^{-23}\,\unit{J/K}$

\item Universelle Gaskonstante\\
$R = 8{,}314\;4598\;(48)\:\unit{J/(mol\,K)}$

\item Plancksches Wirkungsquantum\\
$h = 6{,}626\;070\;040\;(81)\times 10^{-34}\,\unit{Js}$

\item Reduziertes planksches Wirkungsquantum\\
$\hbar = 1{,}054\;571\;800\;(13)\times 10^{-34}\,\unit{Js}$

\item Masse des Elektrons\\
$m_e = 9{,}109\;383\;56\;(11)\times 10^{-31}\,\unit{kg}$

\item Masse des Neutrons\\
$m_n = 1{,}674\;927\;471\;(21)\times 10^{-27}\,\unit{kg}$

\item Masse des Protons\\
$m_p = 1{,}672\;621\;898\;(21)\times 10^{-27}\,\unit{kg}$
\end{enumerate}

\newpage
\section{Einheiten}
\subsection{SI-System}
Newton (Kraft):
\begin{equation}
\unit{N}=\unit{kg\,m/s^2}.
\end{equation}
Watt (Leistung):
\begin{equation}
\unit{W}=\unit{kg\,m^2/s^3}=\unit{VA}.
\end{equation}
Joule (Energie):
\begin{equation}
\unit{J}=\unit{kg\,m^2/s^2}=\unit{Nm}=\unit{Ws}=\unit{VAs}.
\end{equation}
Pascal (Druck):
\begin{equation}
\unit{Pa}=\unit{N/m^2} = 10^{-5}\,\unit{bar}.
\end{equation}
Hertz (Frequenz):
\begin{equation}
\unit{Hz} = \unit{1/s}.
\end{equation}
Coulomb (Ladung):
\begin{equation}
\unit{C} = \unit{As}.
\end{equation}
Volt (Spannung):
\begin{equation}
\unit{V} = \unit{kg\,m^2/(A\,s^3)}
\end{equation}
Tesla (magnetische Flussdichte):
\begin{equation}
\unit{T} = \unit{N/(A\,m)} = \unit{Vs/m^2}.
\end{equation}

\subsection{Nicht-SI-Einheiten}
\begin{tabular}{l|r|l}
\hline
\thbf{Einheit} & \thbf{Symbol} & \thbf{Umrechnung}\pstrut{1pt}\\
\hline
\multicolumn{3}{l}{\thbf{Zeit:}\pstrut{4pt}}\\
\hline
Minute & min & = 60\,s\pstrut{2pt}\\
Stunde & h & = 60\,min = 3600\,s\\
Tag & d & = 24\,h = 86\,400\,s\\
Jahr & a & = 356,25\,d\\
\hline
\multicolumn{3}{l}{\thbf{Druck:}\pstrut{4pt}}\\
\hline
bar & bar & $= 10^5\,\unit{Pa}$\pstrut{2pt}\\
mmHg & mmHg & = 133,322 Pa\\
\hline
\multicolumn{3}{l}{\thbf{Fläche:}\pstrut{4pt}}\\
\hline
Ar & a & $= 100\,\unit{m^2}$\pstrut{2pt}\\
Hektar & ha & = 100\,a = $10\,000\,\unit{m^2}$\\
\hline
\multicolumn{3}{l}{\thbf{Masse:}\pstrut{4pt}}\\
\hline
Tonne & t & = 1000\,kg\pstrut{2pt}\\
\hline
\multicolumn{3}{l}{\thbf{Länge:}\pstrut{4pt}}\\
\hline
Liter & L & = $10^{-3}\,\unit{m^3}$\pstrut{2pt}
\end{tabular}

\subsection{Britische Einheiten}
\begin{tabular}{l|r|l}
\thbf{Einheit} & \thbf{Abk.} & \thbf{Umrechnung}\\
inch & in. & = 2,54\,cm\\
foot & ft. & = 12\,in. = 30,48\,cm\\
yard & yd. & = 3\,ft. = 91,44\,cm\\
chain & ch. & = 22\,yd. = 20,1168\,m\\
&\\[-4pt]
furlong & fur. & = 10\,ch. = 201,168\,m\\
mile & mi. & = 1760\,yd. = 1609,3440\,m
\end{tabular}




\begin{thebibliography}{00}

\bibitem{Hoffmann}
Dirk W. Hoffmann: »Grenzen der Mathematik«.
Springer"=Verlag, Berlin 2011.

\bibitem{Klein}
Felix Klein: »Elementarmathematik vom höheren Standpunkte aus«.\\
Springer"=Verlag, Berlin 1933.

\bibitem{Rubin}
Jean E. Rubin: »Set Theory for the Mathematician«.
Holden"=Day, San Francisco 1967.

\bibitem{Companion}
Timothy Gowers (ed.): »The Princeton Companion to Mathematics«.
Princeton University Press, Princeton 2008.

\end{thebibliography}


\printindex

\end{document}


